\chapter{Representatives, Extensions and Comparison Laws}
An extension is witnessed by a homotopy class together with its image.
Two extension equations with the same source label may nevertheless use
different representatives: their difference lies in higher Adams
filtration, and its image can alter the target. The central question of
this chapter is when representatives can be chosen to satisfy several
equations at once. Crossing conditions control this ambiguity.

Finite lambda quotients make the correction process finite. Their
associated grades occupy a bounded strip, so a residual either produces
a crossing or can be pushed beyond the strip and vanishes.
We use this mechanism first for commutative squares. Applied to the
quotient-boundary square, it gives the generalized Leibniz rule;
applied to a cofiber triangle, it gives the Mahowald rule. Full homotopy representatives are obtained by lifting with
convergence, separation and completeness, as in
Theorems~\ref{sem:Literature:complete-permanent-lifts}
and~\ref{sem:Literature:complete-boundary-lift}.

We use the classical and synthetic Adams systems, quotient maps and
comparison isomorphisms of Chapter~2. Statements quoted from
\cite{LWX126} concern $2$-completed, connective, finite-type spectra.
The level-indexed rules, page-extension comparisons and composition
theorems also allow $2$-completed, bounded-below finite-type spectra, in
the range of Definition~\ref{sem:AdamsSynthetic:paper-spectrum-scope}.
Statements for general filtered groups or convergent towers state their
own hypotheses. We write $F^s\pi_{m,w}U$ for the Adams filtration of a
synthetic object $U$, and $Q_qU=U/\lambda^q$ for its finite lambda
quotient. For a classical spectrum $X$, this gives $Q_q(\nu X)$.
A finite page extension is an extension equation for a normalized map
on one of these quotients, with specified source and target labels.
Definition~\ref{paper:5.4} gives its page and length bounds; the
coordinates following it explain the lambda levels and target
indeterminacies.

For the sphere argument, the main tools are the finite quotient window
in Theorem~\ref{sem:ProofTools:finite-strip}, the boundary comparison in
Proposition~\ref{paper:4.6}, and the nu-cofiber calculation in
Theorem~\ref{sem:ProofTools:nu-zero-layer}. The Leibniz and Mahowald rules
explain how these tools fit together. The later composition criteria
eliminate higher-filtration ambiguity directly in finite quotients.
Numerical examples use the
bounded sphere comparison of \ref{sem:Computation:bounded-comparison}
and the independent classical differentials of
Theorem~\ref{sem:Literature:two-line-classical-inputs}.

\section{Common labels and representative calculus}

\begin{definition}[Common labels and differential equations]
\label{sem:ProofTools:common-labels}
\level{coarse}
For a classical spectrum $X$, a common label in bidegree $(s,t)$ is an element of $E_2^{s,t}(X)$. For $c\geq1$, it is a cycle through $d_c$, respectively a boundary through $d_c$, if it lies in $Z_c^{s,t}(X)$, respectively $B_c^{s,t}(X)$. A permanent label is an element of $Z_\infty^{s,t}(X)$; the zero label is included. The index $c$ is distinct from the lambda level of a synthetic label.

For $c\geq1$, labels $x,y\in E_2^{s,t}(X)$ are congruent modulo $B_c$ if $x-y\in B_c^{s,t}(X)$. For $r\geq2$, the equation $d_r(x)=y$ means that $x$ and $y$ represent classes on page $r$ and their classes satisfy the Adams differential equation, with source bidegree $(s,t)$ and target bidegree $(s+r,t+r-1)$. The target may be zero.
\lean{KIP126.ProofTools.Ambient,KIP126.ProofTools.IsCycle,KIP126.ProofTools.IsBoundary,KIP126.ProofTools.IsPermanent,KIP126.ProofTools.Congruent,KIP126.ProofTools.DifferentialAt}
\end{definition}

For $r\geq2$, a common label lies in $Z_{r-1}$ if and only if it reaches page $r$. If $x$ represents $x_r$ on page $r$ and $y-x\in B_{r-1}$, then $y$ represents $x_r$. If $d_r(x)=y$ and $y'-y\in B_{r-1}$, then $d_r(x)=y'$.

For $c\geq1$, $B_c\subseteq Z_c$, and $B_1=0$. If a common label represents zero on page $r\geq2$, it lies in $B_{r-1}$. Projection of a common cycle to the initial page recovers its $E_2$-label.

Throughout this chapter, \emph{reaching} a page allows the represented
class to be zero. \emph{Nonzero survival} additionally requires that
class to be nonzero. A permanent label may therefore be a true boundary;
its image at infinity is then zero. The notation $d_r(x)=y$ always
includes the page qualifications of both labels.

\begin{theorem}[Filtration, associated grades and detection]
\label{sem:ProofTools:filtration-representatives}
\level{coarse}
For a convergent synthetic Adams tower on $X$, let
\[E_\infty^{s,m+s,w}(X)\cong\operatorname{gr}^s\pi_{m,w}X\]
be its convergence isomorphism.

Every element of $F^s\pi_{m,w}X$ lifts to the $s$-th stage of the Adams tower. If $g:X\to Y$ is a map of convergent towers, then
\[g_*\bigl(F^s\pi_{m,w}X\bigr)\subseteq F^s\pi_{m,w}Y.\]

Let $g:X\to Y$ be a map of convergent towers, $a\in F^s\pi_{m,w}X$ and $b\in F^s\pi_{m,w}Y$. If $g_*(a)=b$, their associated-graded classes satisfy $g_\infty([a])=[b]$. Conversely, this graded equality implies
\[b-g_*(a)\in F^{s+1}\pi_{m,w}Y.\]
If $g_\infty$ is injective in this degree and $g_*(a)\in F^{s+1}\pi_{m,w}Y$, then $a\in F^{s+1}\pi_{m,w}X$.

Let $g:X\to Y$ be a map of convergent towers. If $x$ detects $\alpha$ on $X$, then $g_2(x)$ detects $g_*(\alpha)$ on $Y$. Conversely, suppose that $g_\infty$ is injective in filtration $s$, that $x$ represents an infinite-page class, and that $\alpha\in F^s\pi_{m,w}X$. If $g_2(x)$ detects $g_*(\alpha)$, then $x$ detects $\alpha$.
\uses{sem:AdamsSynthetic:selected-convergence,sem:SpectralSequence:filtration,sem:SpectralSequence:detection,sem:AdamsSynthetic:families}
\lean{KIP126.ProofTools.Internal.filtration_representative_exists,KIP126.ProofTools.Internal.map_preserves_filtration,KIP126.ProofTools.Internal.grade_natural,KIP126.ProofTools.Internal.grade_difference,KIP126.ProofTools.Internal.next_filtration_of_infinity_injective,KIP126.ProofTools.Internal.detects_map,KIP126.ProofTools.Internal.detects_reflect}
\end{theorem}
\begin{proof}
Filtration classes are images of the corresponding tower stages, so they have representatives there and a map of towers preserves them. The kernel of passage from $F^s$ to its associated grade is $F^{s+1}$. Applying this observation to $b-g_*(a)$ proves naturality and the next-filtration ambiguity; injectivity on the associated grade gives the stated reflection.

Under the convergence isomorphisms, detection is precisely equality of associated-graded classes. Naturality transports that equality, and the assumed injectivity reflects it. Functoriality of the tower maps ensures that the resulting infinite-page class is represented by the image of the original homotopy class.
\end{proof}

\begin{theorem}[Representatives for a tower extension equation]
\label{sem:ProofTools:extension-representatives}
\level{coarse}
Let $g:X\to Y$ be a map of convergent towers. Fix $(m,w)$ and classes $x\in E_\infty^{s,m+s,w}(X)$ and $y\in E_\infty^{s+n,m+s+n,w}(Y)$. An extension equation of length $n$ from $x$ to $y$ holds if and only if $n\geq0$ and there are representatives
\[a\in F^s\pi_{m,w}X,\qquad b\in F^{s+n}\pi_{m,w}Y,\qquad g_*(a)=b\]
whose associated-graded classes correspond to $x$ and $y$ under the convergence isomorphisms.
\uses{sem:ProofTools:filtration-representatives,paper:2.3}
\lean{KIP126.ProofTools.Internal.tower_extension_iff}
\end{theorem}
\begin{proof}
Unpack the extension relation in the filtered homotopy groups of the two towers. Its source and target grades are represented by $a$ and $g_*(a)$ in the displayed filtrations; conversely such representatives give the extension equation. The convergence isomorphisms identify these grades with $x$ and $y$.
\end{proof}

\section{Finite filtration correction and quotient windows}

\begin{theorem}[Finite correction of a residual]
\label{sem:ProofTools:finite-correction}
\level{coarse}
Let $A,A',B$ be abelian groups, let $N\in\mathbb N$, and let
$\rho:A'\to A$ and $g:A\to B$ be homomorphisms. For each integer $k$,
let $E_k\leq A$ and $L_k\leq A'$, with $L$ decreasing, and let
$T\leq B$. Fix an integer $s$. Suppose that for each $s\leq k<s+N$, a residual $\delta\in E_k$ can either be corrected by some $\varepsilon\in L_k$ so that $\delta-\rho\varepsilon\in E_{k+1}$, or already satisfies $g\delta\in T$. Suppose also that $g(E_{s+N})\subseteq T$. If $\alpha_0-\rho\alpha\in E_s$, there is $\alpha'$ with
\[\alpha'-\alpha\in L_s,\qquad g(\alpha_0-\rho\alpha')\in T.\]
\lean{KIP126.ProofTools.finite_filtration_correction}
\end{theorem}
\begin{proof}
Induct over the finite interval. If the current residual already has image in $T$, keep the current correction. Otherwise subtract the supplied $\rho\varepsilon$ and add $\varepsilon$ to the source representative, thereby moving the residual to the next filtration. Since $L$ decreases, every accumulated change lies in $L_s$. At the terminal stage the hypothesis on $g(E_{s+N})$ puts its image in $T$. This finite induction proves both conclusions.
\end{proof}

\begin{theorem}[Finite quotient windows and restriction]
\label{sem:ProofTools:finite-strip}
\level{coarse}
For $q>0$, the infinite page of $\nu X/\lambda^q$ vanishes outside $0\leq t-w<q$. Every $\alpha\in\pi_{t-s,t}(\nu X/\lambda^q)$ has filtration at least $s$. If $q\leq m+s-w$, an element of $\pi_{m,w}(\nu X/\lambda^q)$ in filtration at least $s$ is zero.

For $0<i\leq j$ and $t-w<i$, the quotient reduction $\rho:\nu X/\lambda^j\to\nu X/\lambda^i$ induces an injective map on the infinite-page component in tridegree $(s,t,w)$.
\uses{sem:Literature:infinite-finite-window,sem:Literature:infinite-restriction,sem:AdamsSynthetic:finite-quotient-separation,sem:ProofTools:filtration-representatives}
\lean{KIP126.ProofTools.Internal.Mahowald.finite_infinity_zero_outside,KIP126.ProofTools.Internal.Mahowald.finite_bottom_filtration,KIP126.ProofTools.Internal.Mahowald.finite_high_filtration_zero,KIP126.ProofTools.Internal.rho_infty_injective}
\end{theorem}
\begin{proof}
The finite-window comparison gives the stated support of the infinite page. Vanishing associated grades identify successive filtration subgroups; iterating over a finite interval gives the corresponding filtration bounds. Separation of the finite-quotient filtration then makes the subgroup above the upper edge of the strip zero.

In the retained strip, the two quotient windows have the same boundary denominator, and reduction is induced by inclusion of the classical cycle groups. It is therefore injective on the indicated associated grade. This associated-grade injectivity gives the filtration reflection used below.
\end{proof}

\begin{theorem}[Finite target labels and sphere detection]
\label{sem:ProofTools:sphere-label-detection}
\level{coarse}
For $q>0$ and $0\leq k<q$, let $y\in Z_{q-k}^{s,t}(X)$. Its finite target label at weight $t-k$ represents $e\in E_\infty^{s,t,t-k}(\nu X/\lambda^q)$ if and only if the finite-window comparison carries $e$ to the class of $y$ in
\[Z_{q-k}^{s,t}(X)/B_{1+k}^{s,t}(X).\]

The first-quotient comparison carries the first label of $y\in E_2^{s,t}(S^0)$ to $y$. Let $\alpha\in\pi_{t-s,t}(\nu S^0)$ have first-quotient image equal to this label, using the canonical zero-shift identification. For $q>0$ and $k\geq0$, the finite target label of $y$ in tridegree $(s,t,t-k)$ detects the image of $\lambda^k\alpha$ in $\nu S^0/\lambda^q$.
\uses{sem:ProofTools:filtration-representatives,sem:ProofTools:finite-strip,sem:Literature:infinite-finite-window,sem:Literature:infinite-shift,sem:AdamsSynthetic:e2-comparison,sem:AdamsSynthetic:first-quotient-shifts,sem:AdamsSynthetic:ess-shifts}
\lean{KIP126.ProofTools.Internal.Mahowald.finite_target_label_iff,KIP126.ProofTools.Internal.Mahowald.sphere_first_quotient_label,KIP126.ProofTools.Internal.Mahowald.sphere_finite_lambda_detects}
\end{theorem}
\begin{proof}
The finite-window isomorphism identifies the finite target label with the displayed classical cycle coset. At the first quotient this is simply $y$. Quotient projection and multiplication by $\lambda$ carry its representative to the image of $\lambda^k\alpha$, so naturality gives detection at weight $t-k$. The resulting associated grade may be zero.
\end{proof}

\section{Crossings and common representatives}

A change of representative adds a higher-filtration source. Its image
may change the target grade, but the relevant obstruction depends on
which shorter images have already been quotiented out. We will use four
crossing conditions, with different quantifiers:
\begin{enumerate}
\item Actual essential no-crossing concerns one filtered synthetic map at
one stem and weight. It excludes \emph{essential} higher-source
extensions in a filtration triangle, including length zero.
\item The representative no-crossing condition of
Definition~\ref{paper:2.7} excludes \emph{every} higher-source extension
with nonzero target in the stated range, including inessential ones.
For representative relations of the same filtered map and range, this
implies actual essential no-crossing. For a classical extension spectral sequence, use the representative
comparison specified after that definition.
\item An Adams crossing in Definition~\ref{paper:4.14} is an essential
(nonzero) classical Adams differential with the displayed source and target bounds.
Proposition~\ref{paper:4.16} compares it with crossings of the corresponding
quotient-boundary extension.
\item A page-extension crossing in Definition~\ref{paper:5.9} is an
essential extension on a lower effective page, with an additional
classical boundary exclusion. Its general identification with a
synthetic crossing fails, as shown in
Theorem~\ref{sem:ProofTools:finite-crossing-counterexample}.
\end{enumerate}
These distinctions matter: a nonzero target label need not give an
essential extension, and a change of quotient or weight can change its
target indeterminacy. Each theorem states which condition it uses.

\begin{definition}[Actual essential no-crossing]
\label{sem:ProofTools:mh-actual-no-crossing}
\level{coarse}
Fix a map $F:U\to V$ of convergent synthetic Adams towers and one stem and weight $(p,w)$. For a source filtration $s$ and length $a$, actual essential no-crossing means
\[
\text{there is no essential }d_{v-u}^{F}(\xi)=\eta
\quad\text{with }s<u\leq v\leq s+a
\]
in this same map and bidegree. Essentiality is taken in the actual extension spectral sequence, including all shorter extension images. The diagonal $u=v$ is included.
\uses{sem:AdamsSynthetic:selected-convergence,sem:SpectralSequence:filtration}
\lean{KIP126.ProofTools.Internal.Mahowald.ActualNoCrossing}
\end{definition}

\begin{theorem}[A common source from actual no-crossing]
\label{sem:ProofTools:mh-common-source}
\level{coarse}
Let $F:U\to V$ and $G:U\to W$ be maps of convergent towers with the same source convergence. Fix the same stem and weight, and suppose
\[
d_a^F(\xi)=\eta,\qquad d_b^G(\xi)=\zeta.
\]
If either equation has actual no-crossing in its full higher-source triangle, then there are actual representatives $u,v,w$ of $\xi,\eta,\zeta$, respectively, with
\[F(u)=v,\qquad G(u)=w.\]
The representatives lie in filtrations $s,s+a,s+b$. Zero labels, inessential input equations, and length zero are allowed.

The two input equations may initially be witnessed by different representatives of $\xi$. Their difference is in $F^{s+1}$, so changing one representative to the other can change either target. The no-crossing alternative allows that difference to be corrected without changing the prescribed target grades. This is the common-representative question of \cite[Remark 6.18]{LWX126}.
\uses{sem:ProofTools:mh-actual-no-crossing,sem:ProofTools:filtration-representatives,sem:ProofTools:finite-correction}
\lean{KIP126.ProofTools.Internal.Mahowald.actual_common_source_representative}
\end{theorem}
\begin{proof}
\textit{Proof sketch.}
Choose representatives for the two given extension equations. Their source representatives differ in the next source filtration. Synchronize them over the finite higher-source triangle: a first nonzero image grade of a residual that cannot be removed gives an essential crossing of the edge on which actual no-crossing is assumed. The finite correction argument therefore supplies a common source while retaining the two prescribed image grades. The same argument includes diagonal extensions and zero labels.
\end{proof}

Assume that suspension and bigraded shifts are coherent with the homotopy functors. The canonical isomorphism
\[\Sigma^{1,-q}Q\cong\Sigma(\Sigma^{0,-q}Q)\]
identifies the induced suspension and bigraded-shift maps. The inverse and composition isomorphisms for shifts induce the corresponding inverse and composition maps on homotopy groups.
For all integers $m,n,k,l$ and $f,g\in\pi_{m,n}X$, suspension invariance is additive:
\[
\operatorname{susp}_{k,l}(f+g)=\operatorname{susp}_{k,l}(f)+\operatorname{susp}_{k,l}(g)
\quad\text{in }\pi_{m+k,n+l}(\Sigma^{k,l}X).
\]

\section{No-crossing criteria for filtered extensions}

A filtration bound on a map first rules out short extensions. We then
ask when a change of source representative can alter the target of a
longer extension. The criterion below expresses this question entirely
in the filtered homotopy groups.

% AimPaper: 2.6 corollary 327 330
\begin{corollary}[Vanishing below the Adams filtration]
\label{paper:2.6}
\source{AimPaper Corollary 2.6}
\level{coarse}
From \cite{LWX126}, Corollary 2.6.\par

If the Adams filtration of $f$ is $k$, then
    $d^f_i=0$ for $i<k$.

\uses{paper:2.5,sem:AdamsSynthetic:filtration-comparison}

\lean{KIP126.ProofTools.Paper.adamsFiltrationVanishing}
\end{corollary}
\begin{proof}
\textit{(See \cite{LWX126}, Corollary 2.6.)} \href{paper/main.tex.html#L332}{Original proof}.

When $k>0$, it is clear that $d^f_0=0$, which means that for any $s$ and $[x]\in F_s\pi_*(X)$, we have $f([x])\in F_{s+1}\pi_*(Y)$. Since $\AF(f)=k$, we know that there exist $k$ maps
    $$X=X_0\fto{f_1}X_1\fto{f_2}\cdots\fto{f_k}X_k=Y$$
    such that each $f_i$ induces trivial map in $\HF$-homology (hence $\AF(f_i)>0$) and $f$ is homotopic to the composition $f_k\circ \cdots\circ f_1.$
    Then for any $s$ and $[x]\in F_s\pi_*(X)$, inductively we have
    $$f_1([x])\in F_{s+1}\pi_*(X_1)$$
    $$f_2\circ f_1([x])\in F_{s+2}\pi_*(X_2)$$
    $$\cdots$$
    $$f([x])=f_k\circ\cdots\circ f_1([x])\in F_{s+k}\pi_*(X_k)=F_{s+k}\pi_*(Y).$$
    By Proposition \ref{prop:8154e6f1}.(1) this means that $d_i^fx=0$ for any $x\in E_\infty^{*,*}(X)$ and $i<k$.
\end{proof}

% AimPaper: 2.7 definition 344 352
\begin{definition}[Crossings of extensions]
\label{paper:2.7}
\source{AimPaper Definition 2.7}
\level{coarse}
From \cite{LWX126}, Definition 2.7.\par

\label{def:98skj23}
\begin{enumerate}[leftmargin=*]
    \item For $x\in E_\infty^{s,t}(X)$, we say that an $f$-extension $d_n^{f}(x)=y$ has a crossing that hits Adams filtration $p$ for some $p \leq s+n$, if there exists an element $x^\prime\in E_\infty^{s+a,t+a}(X)$ with $a>0$ and an $f$-extension $d_m^{f}(x^\prime)=y^\prime$ for $0\neq y^\prime\in E^{s+a+m, t+a+m}_\infty(Y)$ such that
    $$p = \textup{AF}(y^\prime) = s+a+m\le \textup{AF}(y) = s+n.$$
    \item We say that an $f$-extension $d_n^{f}(x)=y$ has no crossing that hits the range AF $\ge p$ if there does not exist an element $x^\prime\in E_\infty^{s+a,t+a}(X)$ with $a>0$ and an $f$-extension $d_m^{f}(x^\prime)=y^\prime$ for $0\neq y^\prime\in E^{s+a+m, t+a+m}_\infty(Y)$ such that
    $$p\le \textup{AF}(y^\prime) = s+a+m\le \textup{AF}(y) = s+n.$$
    \item We say that an $f$-extension $d_n^{f}(x)=y$ has no crossing if it has no crossing that hits the range AF$\ge \AF(x) + 1$.
\end{enumerate}


Let $f$ be a map of convergent towers or filtered abelian groups. An extension from source filtration $s$ to target filtration $s+n$ has no higher-source crossing over $[u,s+n]$ if every extension with source filtration greater than $s$ and target filtration in $[u,s+n]$ has zero target label. All extensions, including inessential ones, occur in this quantification.
\uses{paper:2.3}


A classical $f$-extension spectral sequence application also requires
that its differential crossings be identified with these
representative-level conditions. Proposition~\ref{paper:2.10} uses
that same comparison when interpreted in the classical ESS.

\lean{KIP126.ProofTools.Paper.FilteredNoCrossingRange,KIP126.ProofTools.Paper.ExtensionNoCrossingRange}

\sketch
\end{definition}

\begin{proposition}[Representatives and crossing-free extensions]
\label{sem:ProofTools:representative-criterion}
\label{paper:2.10}
\label{prop:i8r47oe}
\source{AimPaper Proposition 2.10}
\level{coarse}
\uses{sem:ProofTools:extension-representatives,paper:2.7,sem:ProofTools:finite-correction,paper:2.5}
\lean{KIP126.ProofTools.Paper.extension_noCrossingRange_iff,KIP126.ProofTools.Paper.extension_noCrossing_iff,KIP126.ProofTools.Paper.extension_zero_noCrossing_iff,KIP126.ProofTools.Paper.extension_noCrossing_of_filtration_bound}
\lean{KIP126.ProofTools.Paper.filtered_extension_noCrossingRange_iff,KIP126.ProofTools.Paper.filtered_extension_noCrossing_iff,KIP126.ProofTools.Paper.filtered_zero_extension_noCrossing_iff,KIP126.ProofTools.Paper.filtered_extension_noCrossing_of_filtration_bound}

Fix a map $f$ of convergent towers or filtered abelian groups, and an
extension from $x$ in filtration $s$ to $y$ in filtration $s+n$.
For $p\le s+n$, absence of higher-source extensions with nonzero target
in $[p,s+n]$ is equivalent to
\[
 a\in\{x\},\quad f(a)\in F^p
 \quad\Longrightarrow\quad f(a)\in\{y\}.
\]
Here $\{x\}$ and $\{y\}$ are the sets of representatives with the indicated
associated grades. At $p=s$ the image-filtration premise is automatic.
In particular:
\begin{enumerate}
\item The extension has no crossing if and only if every representative
of $x$ maps to a representative of $y$.
\item For target label zero, no crossing means
$f(a)\in F^{s+n+1}$ for every $a\in\{x\}$; in classical notation this is
$\AF(fa)>\AF(x)+n$.
\item A map increasing every filtration by at least $n$ satisfies this
criterion for its length-$n$ extensions. Thus, when $\AF(f)=n$, its
length-$n$ extensions have no crossing.
\end{enumerate}
This is the representative formulation of \cite{LWX126}, Proposition 2.10.
Reading it as a statement about a general classical $f$-ESS requires the
comparison discussed in Definition~\ref{paper:2.7}.
\end{proposition}
\begin{proof}
Choose representatives of the given extension. Two representatives of
$x$ differ in higher source filtration. If their image grades differ in
$[p,s+n]$, the first nonzero grade of their difference gives a crossing
in that finite interval. Conversely, adding a representative of a
crossing source preserves the source grade and changes the required
image grade. This proves the equivalence. Reading the target grade as
zero gives the strict filtration inequality, and a uniform filtration
increase gives the final case.

Compare the \href{paper/main.tex.html#L418}{original proof} of
\cite{LWX126}, Proposition 2.10.
\end{proof}

% AimPaper: 2.10 proposition 408 416


\section{Commutative squares and extension comparison}

The square rule is a representative chase. Choose one source with the
two prescribed outgoing images, then compare the two composites.
No-crossing on one outgoing edge makes the first choice possible;
no-crossing on the bottom edge determines the common image. A known
zero extension on the right can replace the bottom condition by a
condition only above a filtration cutoff.

% AimPaper: 2.12 theorem 461 481
\begin{theorem}[Commutative squares, triangles and composition of extensions]
\label{paper:2.12}
\label{thm:4114f70c}
\label{paper:2.15}
\label{cor:0012nik}
\label{paper:2.16}
\label{cor:166dc180}
\label{paper:2.17}
\label{cor:290d35ce}
\source{AimPaper Theorem 2.12}
\source{AimPaper Corollary 2.15}
\source{AimPaper Corollary 2.16}
\source{AimPaper Corollary 2.17}
\level{coarse}
\uses{paper:2.10}
\lean{KIP126.ProofTools.Paper.extensionSquare}
\lean{KIP126.ProofTools.Paper.noCrossingExtensionSquare}
\lean{KIP126.ProofTools.Paper.extensionTriangle}
\lean{KIP126.ProofTools.Paper.extensionComposition}

\textit{Square with a filtration cutoff.}
Consider a homotopy commutative diagram of spectra
    $$\xymatrix{
    X \ar[r]^{f}\ar[d]_{p} & Y\ar[d]^{q}\\
    Z \ar[r]_{g} & W
    }$$
    Suppose $m,n,l\ge 0$, $0<k\le m+l-n$,
    $$\begin{array}{ll}
        x\in E^{s,t}_\infty(X) & y\in E^{s+n,t+n}_\infty(Y),\\
        z\in E^{s+m,t+m}_\infty(Z) & w\in E^{s+m+l,t+m+l}_\infty(W),
    \end{array}$$
    and
    \begin{enumerate}
        \item $d^{f}_{n}(x)=y$,
        \item $d^{p}_{m}(x)=z$,
        \item the differential in (1) or the differential in (2) has no crossing,
        \item $d^{g}_{l}(z)=w$, and has no crossing that hits the range AF $\ge s+n+k$,
        \item $d^{q}_{k-1} y=0$, and has no crossing.
    \end{enumerate}
    Then $d^{q}_{m+l-n}(y)=w$. (See Figure \ref{fig:4f154a8}.)

\textit{Square without crossings.}
In the same diagram and bidegrees, suppose instead that
\[
 d_n^f(x)=y,\qquad d_m^p(x)=z,\qquad d_l^g(z)=w,
\]
with $m,n,l\ge0$, that one of the first two extensions has no crossing,
and that the $g$-extension has no crossing. Then
\[
 d_{m+l-n}^q(y)=w.
\]
This is \cite{LWX126}, Corollary 2.15; no right-edge zero extension is
needed under these full no-crossing hypotheses.

\textit{Triangle and composition cases.}
For $X\xrightarrow fY\xrightarrow qZ$ and $p=q\circ f$, let
$x\in E_\infty^{s,t}(X)$,
$y\in E_\infty^{s+n,t+n}(Y)$ and
$z\in E_\infty^{s+m,t+m}(Z)$, with $m,n\ge0$.
If $d_n^f(x)=y$ and $d_m^p(x)=z$ and one of these has no crossing, then
$d_{m-n}^q(y)=z$.
For $X\xrightarrow pZ\xrightarrow gW$ and $q=g\circ p$, let
$x\in E_\infty^{s,t}(X)$,
$z\in E_\infty^{s+m,t+m}(Z)$ and
$w\in E_\infty^{s+m+l,t+m+l}(W)$, with $m,l\ge0$.
If $d_m^p(x)=z$ and the extension $d_l^g(z)=w$ has no crossing, then
$d_{m+l}^q(x)=w$. These are \cite{LWX126}, Corollaries 2.16--2.17.
\end{theorem}
\begin{proof}
\textit{(See \cite{LWX126}, Theorem 2.12.)} \href{paper/main.tex.html#L483}{Original proof}.

First we find a representative $[x]$ of $x$ such that
    \begin{equation}\label{eq:9d45feac}
        p[x]\in \{z\}
    \end{equation}
    and
    \begin{equation}\label{eq:36e72616}
        f[x]\in \{y\}.
    \end{equation}
    If (1) has no crossing, by (2) we can pick $[x]\in \{x\}$ such that (\ref{eq:9d45feac}) holds;
    If (2) has no crossing, by (1) we can pick $[x]\in \{x\}$ such that  (\ref{eq:36e72616}) holds.
    In both cases the other equation also holds because of the no crossing condition.

    By (5), (\ref{eq:36e72616}) and Proposition~\ref{prop:i8r47oe}~(2) we know that $qf[x]$ has AF $\ge s+n+k$. Since the diagram on the Adams $E_\infty$-pages commutes, we have
    $$AF(gp[x])=AF(qf[x])\ge s+n+k.$$
    Combining with (4), (\ref{eq:9d45feac}) and Proposition~\ref{prop:i8r47oe} we have
    $$gp[x]\in \{w\}.$$
    Therefore, $qf[x]\in \{w\}$ and by (\ref{eq:36e72616}) there is a $q$-extension from $y$ to $w$.

For the full no-crossing square, choose a representative of $x$ with
both prescribed $f$ and $p$ images, using the no-crossing alternative.
The $g$ condition identifies $gp[x]$ with the prescribed target $w$.
Commutativity gives $qf[x]=gp[x]$, proving the asserted extension from
$y$ to $w$. The triangle case takes the lower horizontal map to be an
identity, and the composition case takes the upper horizontal map to
be an identity. These specializations retain the displayed filtration
differences and the same representatives.
\end{proof}

% AimPaper: 2.15 corollary 555 574


% AimPaper: 2.16 corollary 576 591


% AimPaper: 2.17 corollary 597 611


\section{Functoriality and exactness of extension spectral sequences}

The square rule gives maps between extension spectral sequences.
For a null composite, it also shows that an image class supports no
outgoing extension. The final result combines this observation with
completeness to turn permanence into a single representative in the
kernel, where homotopy exactness applies.

% AimPaper: 2.18 corollary 617 626
\begin{corollary}[Morphisms of extension spectral sequences]
\label{paper:2.18}
\source{AimPaper Corollary 2.18}
\level{coarse}
From \cite{LWX126}, Corollary 2.18.\par

\label{cor:e7b20ae2}
    Consider a homotopy commutative diagram of spectra
    $$\xymatrix{
    X \ar[r]^{f}\ar[d]_{p} & Y\ar[d]^{q}\\
    Z \ar[r]_{g} & W
    }.$$
    Assume $\ESS{p}_0^{*,*}=\ESS{p}_r^{*,*}$ and $\ESS{q}_0^{*,*}=\ESS{q}_r^{*,*}$ for some $r\ge 0$. Then
    $$(d^p_r, d^q_r): E_\infty^{*,*}(X)\oplus E_\infty^{*,*}(Y)\to E_\infty^{*,*}(Z)\oplus E_\infty^{*,*}(W)$$
    induces a map from the $f$-ESS to the $g$-ESS.

\uses{paper:2.15,paper:2.1}

\lean{KIP126.ProofTools.Paper.extensionSpectralSequenceMorphism}
\end{corollary}
\begin{proof}
\textit{(See \cite{LWX126}, Corollary 2.18.)} \href{paper/main.tex.html#L628}{Original proof}.

Since $\ESS{p}_0^{*,*}=\ESS{p}_r^{*,*}$ and $\ESS{q}_0^{*,*}=\ESS{q}_r^{*,*}$, we know that $d^p_r$, $d^q_r$ have no crossings. By Theorem \ref{cor:0012nik}, we have $d^g_0\circ d^p_r = d^q_r\circ d^f_0$ and therefore $(d^p_r, d^q_r)$ induces a map from $\ESS{f}_1^{*,*}$ to $\ESS{g}_1^{*,*}$. Inductively we can apply Theorem \ref{cor:0012nik} again and show that $(d^p_r, d^q_r)$ induces a map from $\ESS{f}_n^{*,*}$ to $\ESS{g}_n^{*,*}$ and $d^g_n\circ d^p_r = d^q_r\circ d^f_n$ for all $n$.
\end{proof}

% AimPaper: 2.19 corollary 632 635
\begin{corollary}[Permanent cycles from a null composite]
\label{paper:2.19}
\source{AimPaper Corollary 2.19}
\level{coarse}
From \cite{LWX126}, Corollary 2.19.\par

\label{cor:aed3d1a4}
    If the composition $g\circ f$ of two maps $f:X\to Y$, $g:Y\to Z$ is trivial, and $d_n^f(x)=y$ for $x\in E_\infty^{s,t}(X)$ and $y\in E_\infty^{s+n,t+n}(Y)$, then $y$ is a permanent cycle in the $g$-ESS, i.e., we have
    $d^g_m(y)=0$ for all $m\ge 0$.

\uses{paper:2.15,paper:2.5}

\lean{KIP126.ProofTools.Paper.permanentCycle_of_nullComposite}
\end{corollary}
\begin{proof}
\textit{(See \cite{LWX126}, Corollary 2.19.)} \href{paper/main.tex.html#L637}{Original proof}.

By Proposition~\ref{paper:2.5}, choose a representative $[x]$ for which $[y]=f[x]$ is detected by $y$. The null composite gives $g[y]=gf[x]=0$. This same representative therefore has zero image in every filtration quotient. By the representative definition of the extension spectral sequence, every outgoing $g$-extension differential on $y$ vanishes.
\end{proof}

% AimPaper: 2.20 proposition 648 654
\begin{proposition}[Exactness and extension boundaries]
\label{paper:2.20}
\source{AimPaper Proposition 2.20}
\level{coarse}
From \cite{LWX126}, Proposition 2.20.\par

\label{prop:cfe810af}
    Suppose that the sequence of spectra
    $$X\fto{f}Y\fto{g}Z$$
    induces a sequence of homotopy groups
    $$\pi_*X\fto{\pi_*f}\pi_*Y\fto{\pi_*g}\pi_*Z$$
    that is exact in $\pi_*Y$. Then all permanent cycles in $E^{*,*}_\infty(Y)$ of the $g$-ESS are boundaries in the $f$-ESS.

\uses{paper:2.1,paper:2.5,sem:SpectralSequence:completion,sem:AdamsSynthetic:selected-convergence,sem:AdamsSynthetic:completed-adams-scope}

For the $2$-completed, connective, finite-type spectra of the paper and
their bounded-below shifts, the middle homotopy filtration is complete,
its nonnegative filtration quotients are finite, and the target
filtration is separated, by
Theorem~\ref{sem:AdamsSynthetic:completed-adams-scope}.
These hypotheses allow representatives witnessing permanence on different
pages to be replaced by one representative in the kernel of $g_*$.
Zero labels are included.
\lean{KIP126.ProofTools.Paper.extensionBoundary_of_exact}
\end{proposition}
\begin{proof}
\uses{paper:2.5,sem:SpectralSequence:completion,sem:AdamsSynthetic:selected-convergence,sem:AdamsSynthetic:completed-adams-scope}
\textit{(See \cite{LWX126}, Proposition 2.20.)} \href{paper/main.tex.html#L656}{Original proof}.

The representative argument makes the convergence step explicit. Fix
$y\in E_\infty^{s,d+s}(Y)$. For each integer $k$, let $S_k$ consist of the
residue classes in $\pi_dY/F^k\pi_dY$ represented by some $b$ representing
$y$ with $g_*b\in F^k\pi_dZ$. Permanence makes every $S_k$ nonempty.
The finite nonnegative quotients also make every integer-indexed quotient
finite, and the quotient transitions carry $S_l$ into $S_k$ for $k\leq l$.
The finite inverse-system principle gives a compatible family of such
residues. Completeness lifts it to one element $b\in\pi_dY$.

The residue at $k=s+1$ shows that $b$ still represents $y$. For every
$n\geq0$, its residue at $n$ agrees with that of a representative whose
image lies in $F^n\pi_dZ$. Filtration preservation then puts $g_*b$ in
$F^n\pi_dZ$ as well. Separation gives $g_*b=0$, and homotopy exactness
supplies $c\in\pi_dX$ with $f_*c=b$. At $a=\min(0,s)$ the source Adams
filtration is the whole homotopy group, so $c$ defines an
associated-graded class $x$. Proposition~\ref{paper:2.5} gives the required
$f$-extension from $x$ to $y$, so $y$ is an $f$-extension boundary.
\end{proof}

\section{Synthetic boundary maps and extension differentials}

The quotient triangle turns an Adams differential into an extension for
its connecting map. The lambda exponent of the target records how far
the differential travels beyond the quotient length. We first calculate
the extension spectral sequences of the inclusion and restriction maps;
these have only length-zero differentials, so composing with them
transports the boundary formulas without introducing further extensions.

% AimPaper: 4.2 proposition 1025 1029
\begin{proposition}[Extensions for $\lambda^n$ and $\rho$]
\label{paper:4.2}
\source{AimPaper Proposition 4.2}
\level{coarse}
From \cite{LWX126}, Proposition 4.2.\par

\label{prop:9770ae6e}
The only nonzero differentials in the extension spectral sequences for the maps $\lambda^n$ and $\rho$ from Notation~\ref{not:deltaandrho} are the $d_0$'s:
$$d^{\lambda^n}_0=\lambda^n \ \text{and} \ d^{\rho}_0=\rho.$$
As a result, these $d_0$'s have no crossings.

\uses{paper:3.11,paper:3.12,paper:2.19}

\lean{KIP126.ProofTools.Paper.lambda_rho_extension_collapse_finite,KIP126.ProofTools.Paper.lambda_rho_extension_collapse_infinite}

\end{proposition}
\begin{proof}
\textit{(See \cite{LWX126}, Proposition 4.2.)} \href{paper/main.tex.html#L1030}{Original proof}.

Compute the three associated graded groups in the cofiber sequence using Propositions~\ref{paper:3.11} and~\ref{paper:3.12}. Write $k=t-w$. Outside $0\le k<m$ the middle group vanishes. For $0\le k<n$, the $\lambda^n$ source group vanishes and $\rho$ is the inclusion of the relevant cycle subquotients. For $n\le k<m$, the $\rho$ target group vanishes and $\lambda^n$ is the quotient map from a smaller boundary quotient onto the middle group. Thus the sequence of associated graded groups is exact in the middle in every degree.

The two composites are zero. Corollary~\ref{paper:2.19} makes the image of each preceding map a permanent cycle for the next extension spectral sequence. Since each nonzero graded $\lambda^n$ map is surjective and each nonzero graded $\rho$ map is injective in these ranges, no positive-length extension remains. Hence the only differentials are the indicated $d_0$ maps, which have no crossings. The untruncated case uses $Z_\infty$ in the same calculation.
\end{proof}

For $n>0$, the absolute connecting map of the quotient triangle is
\[\delta_n:\nu X/\lambda^n\longrightarrow\Sigma^{1,-n}\nu X.\]
For $0<n<m$, its finite version has target $\Sigma^{1,-n}(\nu X/\lambda^{m-n})$. Using the suspension isomorphism of the quotient triangle, the source label of $x$ is $\lambda^a x$ in the $n$-th quotient. For a classical differential of length $r=k+1$ and $n\leq a+k$, its target label is the shifted $\lambda^{a+k-n}y$; the finite label is its image under the target quotient projection. A crossing of the absolute boundary on page $p$ is an essential extension in that same boundary ESS, with source filtration $u$ and target filtration $v$ satisfying $s<u\leq v\leq s+r$.

% AimPaper: 4.6 proposition 1068 1086
\begin{proposition}[Boundary extensions and their lambda multiples]
\label{paper:4.6}
\label{prop:6de7d130}
\label{paper:4.10}
\label{cor:2a636737}
\source{AimPaper Proposition 4.6}
\source{AimPaper Corollary 4.10}
\level{coarse}
\uses{sem:Literature:finite-bockstein,paper:3.11,paper:3.12,paper:4.2,paper:2.16,paper:2.17}
\lean{KIP126.ProofTools.Paper.syntheticDeltaExtension_prop46_above_absolute,KIP126.ProofTools.Paper.syntheticDeltaExtension_prop46_below_absolute,KIP126.ProofTools.Paper.syntheticDeltaExtension_prop46_above_finite,KIP126.ProofTools.Paper.syntheticDeltaExtension_prop46_below_finite}
\lean{KIP126.ProofTools.Paper.syntheticDeltaExtension_absolute,KIP126.ProofTools.Paper.syntheticDeltaExtension_finite,KIP126.ProofTools.Paper.syntheticDeltaSource_vanishes,KIP126.ProofTools.Paper.syntheticDeltaFiniteTarget_vanishes}

Suppose in the classical Adams spectral sequence of $X$ we have $d_r(x)=y$, where $x\in Z_{r-1}^{s,t}(X)$ and $y\in Z_\infty^{s+r,t+r-1}(X)$. Consider the map
    $$\delta:\nu X/\lambda^n\to \Sus^{1,-n}\nu X/\lambda^{m-n}.$$
    \begin{enumerate}
        \item If $r\ge n+1$, then we view $x$ as an element of
        $$E_\infty^{s,t,t}(\nu X/\lambda^n)\iso Z_n^{s,t}(X),$$
        and $\lambda^{r-n-1} y$ as an element of
        $$E_\infty^{s+r,t+r-1, t+n}(\nu X/\lambda^{m-n})\iso Z_{m-r+1}^{s+r,t+r-1}(X)/B_{r-n}^{s+r,t+r-1}(X).$$
        We then have
        $$d^{\delta}_r(x)=\lambda^{r-n-1}y,$$
        which is trivial if $r>m$.
        \item If $r<n+1$, then we view $\lambda^{n+1-r}x$ as an element of
        $$E_\infty^{s,t,t-n-1+r}(\nu X/\lambda^n)\iso Z^{s,t}_{r-1}(X)/B^{s,t}_{n+2-r}(X),$$
        and $y$ as an element of
        $$E_\infty^{s+r,t+r-1, t+r-1}(\nu X)\iso Z_\infty^{s+r,t+r-1}(X).$$
        In this case, we have
        $$d^{\delta}_r(\lambda^{n+1-r}x)=y.$$
    \end{enumerate}

\textit{All retained lambda multiples.}
For these $x,y,\delta$, the two cases combine into
\begin{equation}\label{eq:841425e9}
 d_r^\delta(\lambda^ax)=\lambda^{a+r-n-1}y,
\end{equation}
whenever $0\le a\le n$ and $0\le a+r-n-1<m-n$.
Beyond these source or target windows the differential is trivial.
This is \cite{LWX126}, Corollary 4.10.
\end{proposition}
\begin{proof}
\textit{(See \cite{LWX126}, Proposition 4.6.)} \href{paper/main.tex.html#L1087}{Original proof}.

First take $m=\infty$ and write $\delta_n=\delta_{n,\infty}$. For $n=1$, Theorem~\ref{sem:Literature:finite-bockstein} identifies the connecting-map extension with the classical Adams differential. Induct on $n$, using the triangle comparing the quotients of lengths $n-1$ and $n$ and multiplication by $\lambda$.

For $r\le n$, the commutative-triangle extension rule transports the induction hypothesis to $\lambda^{n+1-r}x$ and gives target $y$. For $r\ge n+1$, the comparison gives $\lambda d_r^{\delta_n}(x)=\lambda^{r-n}y$ modulo classical $B_{r-1}$. The possible ambiguity in dividing this particular equality by $\lambda$ is contained in $B_{r-n+1}\subseteq B_{r-1}$, by the explicit infinity-page subquotients. This yields $d_r^{\delta_n}(x)=\lambda^{r-n-1}y$ in the required quotient.

Induction on the differential length identifies $B_{r-1}$ here with the sum of the earlier $\delta_n$-extension images: each earlier classical boundary gives its corresponding boundary extension, and the same comparison gives the converse. Finally $\delta_{n,m}$ is $\delta_n$ followed by restriction to the finite target quotient. Proposition~\ref{paper:4.2} and the composition rule transport both formulas to that quotient. The target is zero when its $\lambda$-exponent reaches the quotient length, giving the stated vanishing range.

Multiplication by $\lambda^a$ commutes with the quotient boundary by
naturality of its cofiber triangle. The $\lambda$ and restriction maps
have no positive-length extensions by Proposition~\ref{paper:4.2}.
Their action therefore gives the final formula. Its bounds are exactly
the retained source and target windows; outside those windows the
corresponding label is zero.
\end{proof}

% AimPaper: 4.10 corollary 1138 1144


\section{Crossings of Adams differentials}

The boundary formulas let us translate a crossing question back to
the classical Adams spectral sequence. The next definition gives the
source and target bounds of the possible interfering differential;
the following proposition compares those bounds with the boundary
extension.

% AimPaper: 4.14 definition 1165 1173
\begin{definition}[Crossings of Adams differentials]
\label{paper:4.14}
\source{AimPaper Definition 4.14}
\level{coarse}
From \cite{LWX126}, Definition 4.14.\par

\label{def:classicalcrossdiff}
    Suppose $r\ge n+1$ and $d_r(x)=y$, where
    $$x\in E_2^{s,t}(X), \hspace{0.5cm} y\in E_2^{s+r,t+r-1}(X).$$
    A crossing of $d_r(x)=y$ on the $E_{n+1}$-page refers to an essential Adams differential
    $$d_{r-a-b}(x^\prime)=y^\prime,$$
    where
    $$x^\prime\in E_2^{s+a,t+a}(X), \hspace{.5cm} y^\prime\in E_2^{s+r-b,t+r-b-1}(X),$$
    with $0<a\le n-1$ and $0\le b\le r-n-1$. See Figure \ref{fig:f5eae42a}.

\uses{sem:SpectralSequence:differential,sem:SpectralSequence:cycles-boundaries}

\lean{KIP126.ProofTools.HasCrossingOn}

\end{definition}

% AimPaper: 4.16 proposition 1218 1224
\begin{proposition}[Adams crossings and boundary-extension crossings]
\label{paper:4.16}
\source{AimPaper Proposition 4.16}
\level{coarse}
From \cite{LWX126}, Proposition 4.16.\par

\label{prop:cross-dr-En}
    The Adams differential $d_r(x)=y$ has a crossing on the $E_{n+1}$-page if and only if the corresponding $\delta_n$-extension
    $$d^{\delta_n}_r(x)=\lambda^{r-n-1}y$$
    for
    $$\delta_{n}:\nu X/\lambda^{n}\to \Sus^{1,-n}\nu X$$
    has a crossing.

\uses{paper:4.14,paper:2.7,paper:3.11,paper:3.12,paper:4.6,paper:4.10}

\lean{KIP126.ProofTools.Paper.classical_crossing_iff_delta_crossing}

\end{proposition}
\begin{proof}
\textit{(See \cite{LWX126}, Proposition 4.16.)} \href{paper/main.tex.html#L1225}{Original proof}.

By Propositions \ref{prop:30e8b746}, \ref{prop:59f111f} and \ref{prop:6de7d130}, a crossing of $d^{\delta_n}_r(x)=\lambda^{r-n-1}y$ takes the form
    $$d^{\delta_n}_{r-a-b}(\lambda^a x')=\lambda^{r-n-1-b}y',$$
    where $0<a\le n-1$, $0\le b\le r-n-1$,
    $$\lambda^a x'\in E_\infty^{s+a,t+a,t}(\nu X/\lambda^n)\iso Z_{n-a}^{s+a,t+a}(X)/B_{1+a}^{s+a,t+a}(X),$$
    and
    $$y'\in E_\infty^{s+r-b, t+r-b-1}(\nu X).$$
    By Proposition \ref{cor:2a636737}, we see that this crossing corresponds to the classical Adams differential
    $$d_{r-a-b}(x')=y'.$$
\end{proof}

\section{Page extensions at zero and positive lambda levels}

\begin{definition}[Level extensions and the quotient-boundary square]
\label{sem:ProofTools:level-finite-extension}
\label{sem:ProofTools:extension-level-zero}
\level{coarse}
\uses{paper:5.4,sem:AdamsSynthetic:normalized-map,sem:AdamsSynthetic:ess-shifts,sem:ProofTools:mh-actual-no-crossing}
\lean{KIP126.ProofTools.Standard.levelFiniteSourceLabel,KIP126.ProofTools.Standard.LevelFiniteExtension,KIP126.ProofTools.Standard.LevelFiniteNoCrossing,KIP126.ProofTools.Standard.LevelBoundaryNoCrossing,KIP126.ProofTools.Standard.LevelBottomNoCrossing}
\lean{KIP126.ProofTools.Standard.finiteExtensionAtLevelZero}

For $k\in\mathbb N$, multiplying a label by $\lambda^k$ lowers its
weight by $k$. For a quotient exponent $q$, a normalization shift $e$, and
$x\in E_2^{s,t}(X)$, shift the finite label $\lambda^kx$ by $(0,e)$
and identify the shifted quotient with the quotient of the shifted
spectrum. The resulting normalized source label lies in tridegree
\[(s,t,t+e-k)\quad\text{of}\quad (\Sigma^{0,e}\nu X)/\lambda^q.\]
We use the shift--quotient isomorphism commuting with quotient inclusion
and the absolute boundary. Thus all source labels below are obtained
from the unshifted labels by this same functorial construction.

Let $f:X\to Y$, $e=e(f)$, $A=\Sigma^{0,e}\nu X$, $B=\nu Y$, and $F:A\to B$ be the normalized map. For $n,k\in\mathbb N$, put $q=n+k-1$. A level-$k$ finite extension of length $m$ from $x\in E_2^{s,t}(X)$ to $y\in E_2^{s+m,t+m}(Y)$ consists of the bounds
\[2\le n,\qquad e\le m\le n-2+e,\]
and infinite-page source and target classes, represented by the normalized label $\lambda^k x$ and the quotient projection of $\lambda^{k+m-e}y$, which satisfy
\[d_m^{F_q}(\lambda^k x)=\lambda^{k+m-e}y\]
at fixed stem and weight $(t-s,t+e-k)$. The respective cycle/boundary windows are
\[Z_{n-1}^{s,t}(X)/B_{k+1}^{s,t}(X),\qquad
Z_{n-1-m+e}^{s+m,t+m}(Y)/B_{k+m-e+1}^{s+m,t+m}(Y).\]
Both essential and inessential extensions, including zero labels, are allowed.
We abbreviate it by $\mathcal E^f_{q,k}(m;x,y)$, with the coordinate
convention explained after Definition~\ref{paper:5.4}.
Its complete target indeterminacy is defined using the shorter images
of this same finite map.

For a filtered map $g$ and fixed stem and weight $p$, write $NC(g;p,b,h)$ if there is no essential extension $d_{v-u}^g(\xi)=\eta$ for any $b<u\le v\le b+h$ and any associated-graded classes $\xi,\eta$. Essentiality is taken modulo the complete target indeterminacy of that actual map. For the further lengths $r,l$, with $A,B,F,q$ as above and $p=(t-s,t+e-k)$, the three no-crossing conditions are
\[
NC(F_q;p,s,m),\qquad NC(\delta_q^A;p,s,r),\qquad
NC(\Sigma^{1,-q}F;p,s+r,l).
\]
They refer to the top, left and bottom maps in the quotient-boundary square. The bottom condition is tested on the propagated infinite extension
at its actual level $r-n$.

\textit{Zero-level interpretation.}
Let $n\ge2$ and let $m,s,t\in\mathbb Z$. An original page extension
$d_m^{f,E_n}(x)=y$ supplies
$\mathcal E^f_{n-1,0}(m;x,y)$ for the same labels and map equation.
The cycle and length qualifications are already part of that page
extension. This comparison transports the representative relations;
the no-crossing conditions are specified on the maps at that quotient and weight.
\end{definition}
\begin{proof}
Choose the actual source and target witnesses in the page extension.
The zero-level shift and quotient comparisons identify their labels
with the level-zero labels in the preceding definition. Transporting
these two representative equalities retains the same extension equation.
The page and length bounds are already part of the given witness.
\end{proof}


At fixed $n$, increasing $k$ changes both the quotient length, from
$n-1$ to $n+k-1$, and the weight, to $t+e-k$. Quotient reduction instead
keeps the weight fixed. Division by $\lambda$ is the problem of finding a smaller-layer
representative with a prescribed image. Essentiality and no-crossing are
tested using the shorter-extension indeterminacy at the resulting quotient
and weight.

The crossing conditions must be checked on the normalized map at the
retained quotient and weight. We establish the comparison and lifting
results next; the numerical nu and multiplication-by-two calculations
will then illustrate essentiality and crossings.

\section{Crossing comparisons and conditional lambda division}

The theorem below constructs a preimage under $\lambda^i$ in the required
target filtration. Its target-kernel condition controls low-filtration
lambda torsion by excluding the specified essential smaller-layer extensions.
Theorem~\ref{sem:ProofTools:finite-crossing-counterexample} compares
classical page-extension crossings with synthetic crossings in an explicit example.

% AimPaper: 5.9 definition 1437 1441

\begin{definition}[Crossings of $E_r$-extensions]
\label{paper:5.9}
\source{AimPaper Definition 5.9}
\level{coarse}
From \cite{LWX126}, Definition 5.9.\par

\label{def:41d51149}
    A crossing of the $(f,E_r)$-extension $d_{n}^{f,E_r}(x)=y$ in (\ref{eq:b98be76a}) is defined as an essential $(f, E_{r-a})$-extension from some $x'\in Z_{r-1-a}^{s+a,t+a}(X)$ to
    $$y'\in Z_{r-1-n+b+e(f)}^{s+n-b,t+n-b}(Y)\backslash B_{1+n-b-e(f)}^{s+n-b,t+n-b}(Y)$$
    (where $\backslash$ denotes the difference of sets and it means that $y'$ should survive to the classical Adams $E_{r-n+b+e(f)}$ page while it should not be hit by an Adams differential of length at most $1+n-b-e(f)$) for $0<a \le r-2$ and $0\le b\le n-a-e(f)$. See Figure \ref{fig:66241bf6}.

\uses{paper:5.4}

\lean{KIP126.ProofTools.FiniteHasCrossing}

\end{definition}

% AimPaper: 5.10 proposition 1515 1517

\begin{theorem}[Finite lambda division with a target-kernel condition]
\label{sem:ProofTools:finite-lambda-division}
\source{Lambda division with a target-kernel crossing condition; compare AimPaper Proposition 5.10}
\level{coarse}
For the lambda-division step in \cite[Proposition 5.10]{LWX126},
the following target-kernel condition suffices. The construction supplies a preimage with the prescribed image filtration.

Let $f:X\to Y$, $e=e(f)$, $A=\Sigma^{0,e}\nu X$, $B=\nu Y$, and $F:A\to B$ be the normalized map. Fix $0<i<q$, $p=(t-s,t+e)$ and a target filtration $N$. In the actual lambda square put
\[
A'=\Sigma^{0,-i}(A/\lambda^{q-i}),\quad
B'=\Sigma^{0,-i}(B/\lambda^{q-i}),\quad
F'=\Sigma^{0,-i}F_{q-i}.
\]
Let $L_A:A'\to A/\lambda^q$ and $L_B:B'\to B/\lambda^q$ be the finite lambda inclusions. The extra crossing condition states that, for every
\[s+i\le u\le v<N,\]
an essential extension $d_{v-u}^{F'}(\xi)=\eta$ at $p$ has target satisfying
\[\operatorname{gr}_v(L_B)(\eta)\ne0.\]
The initial source row $u=s+i$ is included, while the terminal target row $v=N$ is excluded. This condition controls the essential smaller-layer targets killed by lambda.

Assume $s+i\le N$ and the additional target-kernel crossing condition. If an actual representative
\[a\in F^{s+i}\pi_{t-s,t+e}(A/\lambda^q)\]
has $F_q(a)\in F^N\pi_{t-s,t+e}(B/\lambda^q)$, there exists
\[b\in F^{s+i}\pi_{t-s,t+e}(A')\]
such that
\[L_A(b)=a,\qquad L_B(F'(b))=F_q(a),\qquad F'(b)\in F^N\pi_{t-s,t+e}(B').\]
Thus $a$ has a preimage whose image lies in the requested filtration,
and the commutative lambda square preserves the equation $F_q(a)$.
The construction determines the smaller-layer source as the grade of $b$. The target-kernel condition is only below $N$,
so kernel ambiguity at the final target grade is allowed.
\uses{sem:ProofTools:finite-correction,sem:ProofTools:finite-strip,sem:AdamsSynthetic:lambda-quotients}
\lean{KIP126.ProofTools.Standard.finiteLambdaDivision}
\end{theorem}
\begin{proof}
\textit{Proof sketch.}
The target-kernel crossing condition controls low-filtration $\lambda$-torsion
by excluding the specified essential lower-filtration images.
Lift the actual source along the finite $\lambda^i$ map. Correct this lift over the finite source and target filtration intervals. A first lower target grade killed by the large-quotient lambda map would give an essential smaller-layer extension forbidden by the stated kernel condition. Finite correction therefore raises the smaller-layer image to $F^N$ while retaining the actual large-layer source and image. The two displayed map equations follow from the actual lambda square. The construction chooses the smaller source label and allows kernel ambiguity at the terminal target grade.
\end{proof}

\begin{theorem}[Failure of the unrestricted finite crossing equivalence]
\label{sem:ProofTools:finite-crossing-counterexample}
\level{coarse}
A counterexample to the universal assertion of
\cite{LWX126}, Proposition 5.10, constructed in the completed background of
\ref{sem:AdamsSynthetic:paper-spectrum-scope}.

The following map gives a counterexample to the universal equivalence between
crossings in Definition~\ref{paper:5.9} and nonzero-target synthetic extensions
in the same finite quotient and weight. Let $S$ be its sphere, let $C_4$ be the cofiber of
$4:S\to S$, and let $j:S\to C_4$ be the cofiber inclusion. Set
\[
M=S\oplus(C_4\oplus S),\qquad f(a,\theta,c)=(4c,j(a),2a).
\]
These objects lie in the paper's spectrum range, and $e(f)=0$.
For $(r,n,s,t)=(4,2,1,1)$ the zero source and target labels give a
finite extension. It has no crossing in Definition~\ref{paper:5.9},
but the actual map on $\nu M/\lambda^3$ in bidegree $(0,1)$ has an
extension from source filtration $2$ to a nonzero target in filtration $3$.
The chosen original extension may be inessential.
\uses{sem:AdamsSynthetic:paper-spectrum-scope,paper:5.4,paper:5.9}
\lean{KIP126.ProofTools.Paper.not_finiteCrossing_iff_syntheticCrossing}
\end{theorem}
\begin{proof}
Write $Q_qX=\nu X/\lambda^q$ and let $g$ be the image of the actual
synthetic $h_0$ in $\pi_{0,1}Q_3S$. The identities
$\lambda h_0^2=2h_0$ and $\lambda^2h_0^3=4h_0$ give
\[
\alpha=(2g,0,0)\in F^2\pi_{0,1}Q_3M,\qquad
f\alpha=(0,0,4g)\in F^3\pi_{0,1}Q_3M.
\]
The normalized cofiber triangle gives $2h_0j=0$. The class $h_0^3$
is nonzero, permanent, and has no incoming Adams differential, so the
finite-quotient comparison gives $4g\ne0$. The same comparison and
filtration separation give $F^4\pi_{0,1}Q_3M=0$. Consequently the target
is nonzero in the actual associated graded group and supplies the required
synthetic extension.

The bounds in Definition~\ref{paper:5.9} leave precisely
$(a,b)=(1,0),(1,1),(2,0)$. In the positive-length window on $Q_2M$,
the map $Q_1j$ is injective on the relevant sphere component. If the
matrix image has filtration $3$, its first source coordinate lies in the
kernel of $Q_2S\to Q_1S$, hence has filtration $3$ and is killed by $2$.
All elements of $\pi_{0,2}Q_2S$ are killed by $4$. Deleting that first
coordinate therefore preserves the source grade and makes the matrix image
zero. This excludes essential extensions, with their full representative
indeterminacy. In the two length-zero windows, the first-quotient matrix
has only the middle component induced by $j$. The identities above and
$2h_0j=0$ show that its degree $(2,2)$ and $(3,3)$ labels belong respectively
to $B_2$ and $B_3$, violating the required boundary exclusions.
\uses{sem:Literature:finite-quotient-labels,sem:Literature:normalized-finite-quotient-row,sem:Literature:first-quotient-products,sem:Literature:synthetic-h-zero,sem:AdamsSynthetic:finite-quotient-separation,sem:Computation:bounded-comparison,sem:Computation:named-basis-identities}
\end{proof}


% AimPaper: 6.1 theorem 1606 1621
% AimPaper: 6.12 theorem 1780 1799

\section{Classical realization and target-preserving lifting}

Inverting $\lambda$ carries a full synthetic representative and its
image to a classical representative and its image. Inverting also
enlarges the boundary indeterminacy, so essentiality may be lost.
For a finite-quotient input, a lift preserving both detections supplies
the full representative. The additive criterion below characterizes
such target-preserving lifts.

% AimPaper: 5.14 proposition 1581 1587

\begin{proposition}[From infinity-page extensions to homotopy extensions]
\label{paper:5.14}
\source{AimPaper Proposition 5.14}
\level{coarse}
From \cite{LWX126}, Proposition 5.14.\par

Consider $f: X\to Y$, $x\in Z_\infty^{s,t}(X)$ and $y\in Z_\infty^{s+n,t+n}(Y)$.
    Then
    $$d^{f,E_\infty}_n(x)=y$$
    implies that
    $$d^{f}_n(x+B_\infty^{s,t}(X))=y+B_\infty^{s+n,t+n}(Y).$$

\uses{paper:5.4,sem:Literature:lambda-inversion,paper:3.11,paper:2.5,sem:Literature:realization-coordinates,sem:Literature:lambda-inversion-adams-comparison,sem:Literature:lambda-inversion-detection}

As throughout \cite{LWX126}, spectra are understood to be 2-completed,
connective and of finite type. By the representative criterion of
Proposition~\ref{paper:2.5}, the conclusion means that
$x$ detects one $a$ and $y$ detects $f_*a$. The realization isomorphisms and Adams-tower comparison of
Definition~\ref{sem:Literature:lambda-inversion-adams-comparison}
identify these detections with the classical ones.
Zero labels, boundary labels and inessential extensions are allowed.

\textit{Positive levels and finite quotients.}
A full-object equation at level $k$ uses source $\lambda^kx$
and target $\lambda^{k+n-e}y$ at weight $t+e-k$.
Comparing it with the classical equation requires the detection
transport at that weight and the normalized-map comparison.
The displayed conclusion concerns the zero-level infinite extension
with these fixed comparisons.
For a finite-quotient equation, apply full-object transport to a
lift preserving both specified detections.
\lean{KIP126.ProofTools.Paper.infiniteExtension_to_classicalExtension}
\end{proposition}
\begin{proof}
\uses{sem:ProofTools:extension-representatives,sem:Literature:lambda-inversion-detection,sem:Literature:lambda-inversion-adams-comparison,sem:AdamsSynthetic:normalized-map,paper:3.11,paper:2.5}
\textit{(See \cite{LWX126}, Proposition 5.14.)} \href{paper/main.tex.html#L1589}{Original proof}.

Invert $\lambda$ in the filtered synthetic map $\widehat f$. The comparison with the classical map $f$ induces a morphism of extension spectral sequences. On the starting page, Proposition~\ref{paper:3.11} sends a label in $Z_\infty/B_{1+t-w}$ to its class in $Z_\infty/B_\infty=E_\infty$.

For the representative transport, put $e=e(f)$ and $w=t+e$.
The actual extension equation gives one representative
$\beta\in\pi_{t-s,w}(\Sigma^{0,e}\nu X)$ whose image under $\hat f$
has the specified target detection. Apply the detection transport theorem
to this same $\beta$ and to $\hat f_*\beta$. The prescribed $E_2$
comparisons return precisely $x$ and $y$.

Realization of the normalized factorization, followed by cancellation of
the invertible realized lambda map, gives
\[
\theta_Y\circ\operatorname{Re}(\hat f)
 =f\circ\theta_{\Sigma^{0,e}\nu X}.
\]
Thus the first realized representative is $a$, and the second is exactly
$f_*a$. The sphere-coordinate transport accounts for
$(t+n)-(s+n)=t-s$. By Proposition~\ref{paper:2.5}, these detections give
the classical $f$-extension in the conclusion. The target is taken in its
own bidegree $(s+n,t+n)$, as required for the quotient by its boundary
group. The target coset after localization uses the enlarged boundary
indeterminacy. The argument applies to zero and boundary labels and to
inessential extensions.
\end{proof}

As observed in \cite{LWX126}, Remark 5.15, the classical comparison
in Proposition~\ref{paper:5.14} forgets the original lambda layer,
the finer synthetic target coset, and the chosen representative and
its image. Recovering these data requires a target-preserving lift.
The following criterion describes the corresponding additive
obstruction.

\begin{proposition}[Lifting a specified family of finite representatives]
\label{sem:ProofTools:ml-representative-family-lift}
\level{coarse}
An additive exactness criterion for the extra lifting step needed by
classical extensions and permanence outputs.\par
Let $q:D\to A$ and $\delta:A\to C$ be homomorphisms of abelian
groups with $\ker\delta=\operatorname{im}q$. Let $V\leq A$ be the
subgroup of permitted corrections preserving the specified source and
image labels, and fix $a_0\in A$. Then
\[
 (a_0+V)\cap\operatorname{im}q\ne\varnothing
 \quad\Longleftrightarrow\quad \delta(a_0)\in\delta(V).
\]
Equivalently, the obstruction is the class of $\delta(a_0)$ in
$C/\delta(V)$. If $q$ is restriction from the full synthetic homotopy
group, its vanishing gives a full lift of a corrected representative. For $q$ a restriction between two finite quotients,
the same criterion supplies the corresponding finite lift.
\uses{sem:AdamsSynthetic:lambda-quotients}
\lean{KIP126.ProofTools.Standard.representativeFamilyLift_iff}
\end{proposition}
\begin{proof}
If $q(d)=a_0+v$, then $0=\delta(a_0)+\delta(v)$, so
$\delta(a_0)=\delta(-v)\in\delta(V)$. Conversely, write
$\delta(a_0)=\delta(v)$ with $v\in V$. Then
$\delta(a_0-v)=0$, and exactness gives $q(d)=a_0-v$ for some $d$.
Since $-v\in V$, this is a lift of a representative in the specified family.
\end{proof}

\section{The quotient-boundary square and the Leibniz rule}

In the quotient-boundary square, a classical differential supplies the
left edge. A finite $f$-extension supplies the top edge, and a full
$f$-extension supplies the bottom edge after multiplication by the
appropriate power of $\lambda$. The square rule then determines a
right boundary extension. Translating that boundary back to the Adams
spectral sequence gives a differential of length $r+l-m$.
The crossing hypotheses are precisely those needed to carry out this
representative chase at one stem and weight.

\begin{theorem}[Generalized Leibniz and its zero-level form]
\label{sem:ProofTools:generalized-leibniz-at-level}
\label{sem:ProofTools:ss-leibniz-specialization}
\level{coarse}
\uses{sem:ProofTools:level-finite-extension,paper:2.12,paper:4.6,sem:Literature:finite-bockstein,sem:Literature:infinite-finite-window,sem:Literature:infinite-nu-window,sem:Literature:infinite-shift,sem:AdamsSynthetic:ess-shifts,sem:ProofTools:extension-level-zero}
\lean{KIP126.ProofTools.Standard.generalizedLeibnizAtLevel}
\lean{KIP126.ProofTools.Standard.ssLeibnizAtLevelZero}

A level-indexed form of \cite[Theorem 6.1]{LWX126}, with the
synthetic no-crossing conditions stated below.

Let $f:X\to Y$, $e=e(f)$, and $F:\Sigma^{0,e}\nu X\to\nu Y$
be its normalized map. Suppose
\[
2\le n\le r,\qquad e\le m\le n-2+e,\qquad l\ge e.
\]
Fix $k\ge0$, put $q=n+k-1$, and take labels
\[
\begin{array}{ll}
x\in E_2^{s,t}(X),&y\in E_2^{s+m,t+m}(Y),\\
x_\infty\in Z_\infty^{s+r,t+r-1}(X),&
y_\infty\in Z_\infty^{s+r+l,t+r+l-1}(Y).
\end{array}
\]
Assume the three input relations
\[
d_r(x)=x_\infty,\qquad
 d_m^{F_q}(\lambda^kx)=\lambda^{k+m-e}y,\qquad
 d_l^F(x_\infty)=\lambda^{l-e}y_\infty.
\]
The middle relation is in the quotient modulo $\lambda^q$;
the last is an infinite extension at input level zero.
Write $A=\Sigma^{0,e}\nu X$, $\delta_q^A:A/\lambda^q\to\Sigma^{1,-q}A$,
and $p=(t-s,t+e-k)$. Require actual no-crossing
\[
\bigl(NC(\delta_q^A;p,s,r)\lor NC(F_q;p,s,m)\bigr)
\quad\text{and}\quad NC(\Sigma^{1,-q}F;p,s+r,l).
\]
Thus crossing is checked on the left or top edge, and on the bottom edge,
of the following extension diagram at the same stem and weight $p$:
\[
\begin{array}{ccc}
\lambda^kx&\xrightarrow{\ d_m^{F_q}\ }&\lambda^{k+m-e}y\\
{\scriptstyle d_r^{\delta_q^A}}\downarrow&&
\downarrow{\scriptstyle d_{r+l-m}^{\delta_q^{\nu Y}}\ (\mathrm{output})}\\
\lambda^{r-n}x_\infty&\xrightarrow{\ d_l^{\Sigma^{1,-q}F}\ }&
\lambda^{r-n+l-e}y_\infty.
\end{array}
\]
\textit{Extension-relation diagram for the generalized Leibniz rule.}

The lower labels are transported into the suspended objects
$\Sigma^{1,-q}A$ and $\Sigma^{1,-q}\nu Y$.
Then the classical Adams differential is
\[
\boxed{d_{r+l-m}(y)=y_\infty.}
\]
There is no upper bound on $k$. The left and bottom relations in the
diagram are constructed from the inputs, so they add no hypotheses.

\textit{Zero-input-level form.}
Set $k=0$ and rename $(r,m,l)$ as $(b,M,L)$. Thus
$2\le n\le b$, $e\le M\le n-2+e$, and $L\ge e$.
The finite input is $\mathcal E^f_{n-1,0}(M;x,y)$ at
$p=(t-s,t+e)$. The crossing conditions are
\[
 \bigl(NC(F_{n-1};p,s,M)\lor NC(\delta_{n-1}^A;p,s,b)\bigr)
 \land NC(\Sigma^{1,-(n-1)}F;p,s+b,L),
\]
and the conclusion is $d_{b+L-M}(y)=y_\infty$.
An input $d_M^{f,E_n}(x)=y$ supplies the finite relation by the
zero-level interpretation in Definition~\ref{sem:ProofTools:level-finite-extension}.
The other labels, permanent representatives and crossing conditions are
exactly the corresponding ones in the general statement.
\end{theorem}
\begin{proof}
\textit{Proof sketch.}
Use the quotient-boundary square
\[
\begin{array}{ccc}
A/\lambda^q&\xrightarrow{F_q}&\nu Y/\lambda^q\\
\delta_q^A\downarrow&&\downarrow\delta_q^{\nu Y}\\
\Sigma^{1,-q}A&\xrightarrow{\Sigma^{1,-q}F}&\Sigma^{1,-q}\nu Y.
\end{array}
\]
\textit{Commutative quotient–boundary square.}

The classical $d_r$ supplies the left relation. Multiplication by
$\lambda^{r-n}$ transports the infinite input to the bottom relation,
with exactly the labels displayed above. Apply the filtered-square
extension theorem with the assumed no-crossing conditions. Its right
boundary relation translates, by the weighted boundary comparison, into
$d_{r+l-m}(y)=y_\infty$. The construction retains the specified labels and every displayed power of $\lambda$.

\textit{At $k=0$.} The finite input is in the quotient of length $n-1$;
the preceding representative comparison supplies it from an
$(f,E_n)$-extension. The no-crossing hypotheses concern the indicated
synthetic maps at that weight. In \cite[Theorem 6.1]{LWX126}, they
are expressed through classical crossing predicates and the comparison
of \cite[Proposition 5.10]{LWX126}. The theorem uses the synthetic conditions at the indicated weight;
Theorem~\ref{sem:ProofTools:finite-crossing-counterexample} describes
the distinction between the two crossing predicates.

The zero-input-level form is the substitution $k=0$ and
$(r,m,l)=(b,M,L)$. Its quotient is $Q_{n-1}$ and its weight is $t+e$;
all no-crossing conditions remain conditions on those actual maps.
\end{proof}


\begin{proposition}[The cutoff extension square and its Adams consequence]
\label{sem:ProofTools:ss-cutoff-square-specialization}
\label{paper:6.5}
\label{sem:ProofTools:ss-adams-soundness}
\level{coarse}
\uses{paper:2.12,sem:ProofTools:level-finite-extension,sem:ProofTools:generalized-leibniz-at-level,paper:4.6,sem:ProofTools:filtration-representatives}
\lean{KIP126.ProofTools.Standard.ssCutoffSquareAtLevel}
\lean{KIP126.ProofTools.Standard.cutoffLeibnizSquare}
\lean{KIP126.ProofTools.Standard.cutoffAdamsDifferentialAtLevel}

\textit{Extension-square form.}
Apply the four-edge theorem \ref{paper:2.12} to the quotient-boundary square.

Fix $F:A\to B$ with $A=\Sigma^{0,e}\nu X$, $B=\nu Y$. For
$n\ge2$ and a lambda level $k\ge0$ put $q=n+k-1$. Use the actual
commutative square
\[
\begin{array}{ccc}
Q_qA&\xrightarrow{F_q}&Q_qB\\
\delta_q^A\downarrow&&\downarrow\delta_q^B\\
\Sigma^{1,-q}A&\xrightarrow{\Sigma^{1,-q}F}&\Sigma^{1,-q}B.
\end{array}
\]
Here $Q_qA=A/\lambda^q$. Fix a stem and weight $p=(a,w)$, and actual
infinite-page classes $\xi,\eta,\zeta,\omega$ in the four vertices,
at filtrations $s,s+M,s+b,s+b+L$, respectively; the remaining degree
coordinates are $(a+\mathrm{filtration},w)$.
Let $M,b,L\ge0$, $D=b+L-M$ and $1\le\kappa\le D$.
Assume
\[
d_M^{F_q}(\xi)=\eta,\qquad
d_b^{\delta_q^A}(\xi)=\zeta,\qquad
d_L^{\Sigma^{1,-q}F}(\zeta)=\omega.
\]
Require no crossing on the top or left equation. Require that the
bottom equation have no crossing hitting filtration at least
$s+M+\kappa$. Finally require the actual right-edge equation
\[
d_{\kappa-1}^{\delta_q^B}(\eta)=0
\]
with no crossing. Every crossing here is measured in the extension
spectral sequence of its actual edge, at the fixed stem and weight,
with its full shorter-image indeterminacy. Then
\[
d_D^{\delta_q^B}(\eta)=\omega.
\]
This is an extension-square conclusion. To read it as a specified
classical Adams differential, its labels must also have the required
boundary-comparison form. For example, use the classical labels,
degrees and range bounds of \ref{sem:ProofTools:ss-leibniz-specialization}.
If the actual comparisons give
\[
\begin{gathered}
\xi=\lambda^k x,\quad
\eta=\lambda^{k+M-e}y,\\
\zeta=\lambda^{b-n}u,\quad
\omega=\lambda^{b-n+L-e}v,
\end{gathered}
\]
at weight $t+e-k$, then the exponent identity
\[
D-q-1+(k+M-e)=b-n+L-e
\]
is precisely the arithmetic required for the right boundary to decode
as $d_Dy=v$. Those range bounds give $D\ge2$ and
$0\le k+M-e<q$. The qualified labels and boundary comparisons identify this exponent
calculation with the stated differential.

\textit{Classical Adams form.}
The classical form uses the extension-square conclusion above and the
weighted boundary comparisons in the proof of
\ref{sem:ProofTools:generalized-leibniz-at-level}.

Let $f:X\to Y$ be a map of $2$-completed, bounded-below finite-type spectra, $e=e(f)$,
$A=\Sigma^{0,e}\nu X$, $B=\nu Y$, and let $F:A\to B$ be the
normalized map. Take
\[
2\le n\le b,\qquad e\le M\le n-2+e,\qquad L\ge e,\qquad k\ge0,
\]
and put $q=n+k-1$, $p=(t-s,t+e-k)$ and $D=b+L-M$.
Specify
\[
\begin{gathered}
x\in E_2^{s,t}(X),\qquad y\in E_2^{s+M,t+M}(Y),\\
u\in Z_\infty^{s+b,t+b-1}(X),\qquad
v\in Z_\infty^{s+b+L,t+b+L-1}(Y),
\end{gathered}
\]
with the actual relations
\[
d_bx=u,\qquad
\mathcal E^f_{q,k}(M;x,y),\qquad
d_L^F(u)=\lambda^{L-e}v.
\]
The middle relation has specified infinite-page source and target
classes in the finite quotients, with labels
$\lambda^kx$ and $\lambda^{k+M-e}y$. Denote its target class by
$\eta$. Boundary labels and zero labels are allowed.

Choose $1\le\kappa\le D$. On the quotient-boundary square
displayed above, require
\[
NC(F_q;p,s,M)\ \lor\ NC(\delta_q^A;p,s,b).
\]
For the propagated bottom map $\Sigma^{1,-q}F$, require that there
be no essential extension whose source and target filtrations
$a,c$ satisfy
\[
s+b<a\le c\le s+b+L,\qquad s+M+\kappa\le c.
\]
Finally require the right-edge zero extension from the same $\eta$,
\[
d_{\kappa-1}^{\delta_q^B}(\eta)=0,\qquad
NC(\delta_q^B;p,s+M,\kappa-1).
\]
All these conditions are at the fixed stem and weight $p$ in the
extension spectral sequences of the displayed actual maps. The
crossing ranges include length zero and their terminal target
filtrations. Then $D\ge2$ and the specified classical labels satisfy
\[
d_D(y)=v.
\]
This includes the cycle qualifications on $E_D$.
Setting $k=0$ gives the zero-input-level rule in quotient $Q_{n-1}$;
the cutoff $\kappa$ is independent of the lambda level $k$.

The abstract version of the same rule, without choosing classical labels,
is \cite{LWX126}, Remark 6.5. Retain the three input relations and
top-or-left no-crossing. For top, left and bottom lengths $m,r,l$,
a cutoff $0<c\le r+l-m$ requires bottom no-crossing only at filtrations
at least $s+m+c$, together with the crossing-free right-edge zero
extension of length $c-1$. It returns a right extension of length
$r+l-m$. This formulation applies at the actual quotient and weight of
each level; decoding its right edge as a classical differential also
uses the qualified boundary labels above.
\end{proposition}
\begin{proof}
\textit{Extension-square form.}
Apply the four-edge theorem to this square, replacing its top length
by $M$, its left length by $b$, its bottom length by $L$, and its
cutoff by $\kappa$. The hypotheses include its right-edge zero
extension and both required crossing conditions. It returns length
$b+L-M=D$ on the right. The theorem applies to every quotient
$q=n+k-1$; setting $k=0$ gives the zero-input-level square with the
same cutoff hypotheses. The displayed exponent identity follows by
substitution of $D$ and $q$.

\textit{Classical Adams form.}
Use the actual left-edge and bottom-edge constructions from the
proof of \ref{sem:ProofTools:generalized-leibniz-at-level}. The
differential $d_bx=u$ gives the left extension from the same source
class as the finite top extension. Its target is the propagated
class labelled by $\lambda^{b-n}u$. Multiplication by
$\lambda^{b-n}$ and the actual boundary suspension transport the
given infinite extension to the bottom edge, with specified target
label $\lambda^{b-n+L-e}v$. These constructions identify the actual
source and target classes, including their suspension comparisons.

The extension-square form now gives a length-$D$ right extension
from $\eta$ to this specified bottom target. Choose a representative
$a$ witnessing this extension. It detects $\lambda^d y$ in $Q_qB$,
where $d=k+M-e$, and its own actual boundary has the specified
bottom target detection. Undoing the boundary suspension gives a
representative in $\nu Y$ with the corresponding weighted label
of $v$.

For completeness, the bounds needed by the weighted reverse
boundary comparison are
\[
\begin{gathered}
0\le d<q,\qquad
q-d=n-1-M+e\le D-1,\qquad D\ge2,\\
(t+M)-d=t+e-k,\qquad
D-q-1+d=b-n+L-e.
\end{gathered}
\]
Apply that comparison to this same $a$ and its boundary, with quotient
length $q$, input degree $(s+M,t+M)$, input lambda exponent $d$ and
differential length $D$. It gives the Adams relation
$d_Dy=v$, with the required cycle qualifications. The displayed
identities verify the range and weight of these detected
representatives; they are used after the representative comparisons,
not as substitutes for them.
\end{proof}

\subsection{Why the crossing hypothesis is needed}

The multiplication-by-two example illustrates both the crossing
hypothesis in the Leibniz rule and the role of the target coset.
We begin with its independent nonzero classical differential, then
construct the extension representatives.

\begin{proposition}[A crossing obstructing a Leibniz conclusion]
\label{paper:6.8}
\label{paper:6.6}
\label{paper:6.9}
\label{paper:6.10}
\label{exam:Leibnizno}
\label{rem:chuaerror}
\level{coarse}
From \cite{LWX126}, Example 6.8 and Remarks 6.6, 6.9--6.10.\par
Let $f=2:S^0\to S^0$, normalized by $[h_0]$. At
\[
n=r=2,\quad m=1,\quad l=2,\quad e(f)=1,\quad s=1,\quad t=16,
\]
the input relations are
\[
d_2(h_4)=h_0h_3^2,\qquad
d_1^{f,E_2}(h_4)=h_0h_4,\qquad
d_2^{f,E_\infty}(h_0h_3^2)=0.
\]
The bottom relation has the essential crossing
\[
d_1^{f,E_\infty}(d_0)=h_0d_0.
\]
Consequently, omitting the bottom no-crossing hypothesis from the original Leibniz rule would give the incorrect conclusion $d_3(h_0h_4)=0$.

The same shorter extension also makes
\[
d_2^{f,E_\infty}(h_0h_3^2)=h_0d_0
\]
inessential. Applying a rule without the required crossing control to these two choices would give the conflicting outputs
\[
d_3(h_0h_4)=0,\qquad d_3(h_0h_4)=h_0d_0.
\]
The maximal-image remark makes a stronger assertion about actual
homotopy classes. It chooses a lift of $h_0h_3^2$ in
$\pi_{14,17}S^{0,0}$ whose $[h_0]$-image is zero in
$\pi_{14,18}S^{0,0}$. A second lift has image $\lambda\beta$,
where $\beta\in\pi_{14,19}S^{0,0}$ has label $h_0d_0$.
These weights display the lambda factor in the article's bracket
notation. The zero-image choice is maximal in the sense discussed there;
the essential crossing above explains the nonzero classical differential
in this example.
\uses{paper:5.4,paper:5.7}
\lean{KIP126.ProofTools.Standard.paper6_8_complete_crossing_warning}
\end{proposition}
\begin{proof}
\uses{sem:Literature:one-line-differentials,sem:Literature:two-line-classical-inputs,sem:Computation:bounded-comparison,sem:Literature:finite-bockstein,sem:Literature:infinite-nu-window,sem:Literature:complete-permanent-lifts,sem:Literature:realization-kernel,sem:Literature:complete-detection,sem:Literature:filtered-product-detection,sem:AdamsSynthetic:finite-separation}
\textit{The nonzero classical conclusion.} The independent classical input in
\ref{sem:Literature:two-line-classical-inputs} gives a nonzero $d_3$ from
the specified $h_0h_4$ to degree $(5,19)$. The bounded Ext calculation gives
one basis element there, namely $h_0d_0$, where $d_0$ is the generator
in degree $(4,18)$ fixed by the sphere comparison. A one-dimensional $\mathbb F_2$ space has
exactly one nonzero element, so the target is $h_0d_0$.

\textit{The extension representatives.} For the normalized
map, the first finite relation comes from multiplication by $h_0$
on the first lambda quotient. Apply the specified Bockstein
formulas to the independent differentials
$d_2(h_4)=h_0h_3^2$ and $d_3(h_0h_4)=h_0d_0$.
Their naturality constructs full classes $a$ and $\beta$ with
the prescribed labels and $[h_0]a=\lambda\beta$.
The class $d_0$ has a full lift $D$: all its outgoing differentials
vanish, by the bounded Ext calculation and the sphere vanishing line.
The actual first-quotient product shows that $[h_0]D$ has label
$h_0d_0$. Hence $a-\lambda D$ keeps the original source label,
while its image is
\[
\lambda\bigl(\beta-[h_0]D\bigr)\in F^6\pi_{14,18}S^{0,0}.
\]
It therefore has zero target label in filtration five. This constructs the zero-target relation by placing the image in filtration six. The length-one extension from $d_0$ has
nonzero target and is essential in its minimal-length range.

\textit{Crossing and inessentiality.} Express the shorter relation
at weight $18$ using $\lambda D$ and $\lambda[h_0]D$.
Its target grade is nonzero because $h_0d_0\notin B_2$, as follows
from the nonzero classical $d_3$. In the minimal-length range,
the actual no-crossing criterion proves this relation essential.
This proves essentiality in the target coset at weight $18$.
It is the required crossing. The two long relations have the
same source infinite-page class, since their actual sources
differ in higher filtration. Uniqueness of the extension
differential, compared with the zero-target witness, shows that
every long witness with target $h_0d_0$ is inessential.
Its target infinite-page class is nevertheless nonzero.

\textit{The actual zero-image lift.} The boundary construction
of $a$ gives $\lambda a=0$, so its realization is zero.
The class $d_0$ survives nontrivially. Product detection and the
fact that $h_0d_0$ is a classical boundary show that the
realization of $2D$ is infinitely Adams-filtered. Ordinary
sphere separation makes it odd torsion; choose an odd integer
$N$ annihilating it. Replacing $D$ by $ND$ preserves its
first-quotient label. Put
\[
z=[h_0](a-\lambda ND)\in F^6\pi_{14,18}S^{0,0}.
\]
This class has zero realization. In this stem the filtration-six
Ext group is generated by $h_0^2d_0$, a $d_3$ target, and all
higher filtration groups vanish. The corresponding synthetic
window therefore puts $z$ in every filtration. Zero realization
also makes $z$ lambda-torsion.

For weights at most $18$, the actual first-quotient exact sequence
shows that the lambda kernel in this stem has no infinitely
filtered nonzero class. At weights $17$ and $18$, its possible
generators are the boundaries of $h_4$ and $h_0h_4$;
their specified nonzero associated grades exclude infinite
filtration. At lower weights the relevant first-quotient group
vanishes. Induction on a lambda-annihilating power, which only
lowers weight, now gives $z=0$. Thus $a-\lambda ND$ has the
original source label and actual zero $[h_0]$-image. The fixed
sphere and suspension comparisons transport this equation to
the displayed $S^{0,0}$ coordinates.
That same zero image is divisible by every power of lambda,
so it is a maximal-image choice. The nonzero classical $d_3$
from above still excludes the zero conclusion.

This proves the specified zero-layer example.
\end{proof}

\section{The generalized Mahowald trick modulo a power of lambda}

The Mahowald rule uses a cofiber triangle to solve two equations at
once. A boundary equation gives a compatible pair in two finite
quotients; the triangle lifts that pair to a single source with a
prescribed restriction and image. No-crossing is used before this step
to choose the compatible representatives.
Throughout this section, $Q_DU=U/\lambda^D$ for a synthetic object
$U$. We keep quotient length, lambda exponent, Adams filtration and
page number as separate indices.

\begin{lemma}[A common finite-quotient lift with two prescribed images]
\label{sem:ProofTools:finite-common-lift}
\level{coarse}
The finite version of the signed triangle construction in
Lemma~\ref{paper:6.11} takes a compatible pair of classes and lifts
both at once.
Fix a classical triangle of $2$-completed, bounded-below finite-type spectra,
with $e_f+e_g+e_h=1$, and fix its distinguished normalized triangle
\[
 U\xrightarrow F V\xrightarrow G W\xrightarrow H\Sigma U.
\]
Let $A,B>0$, $N=A+B$, and choose an isomorphism
\[
 \varepsilon_V:Q_N(\Sigma V)\cong\Sigma Q_NV
\]
whose square with the quotient inclusions commutes. Use the maps
\[
\begin{gathered}
 \delta_{A,B}:Q_AW\longrightarrow\Sigma^{1,-A}Q_BW,
 \qquad G_B=\Sigma^{1,-A}Q_BG,\\
 \rho:Q_N(\Sigma U)\longrightarrow Q_A(\Sigma U),
 \qquad H_A=Q_AH,\\
 F_N=\varepsilon_V\circ Q_N(\Sigma F),\qquad
 \iota_{A,B}:\Sigma^{1,-A}Q_BV\longrightarrow\Sigma Q_NV,
\end{gathered}
\]
where $\delta_{A,B}$ is the quotient boundary, $\rho$ is restriction,
and $\iota_{A,B}$ is the suspended lambda inclusion.
For every $m,w\in\mathbb Z$ and homotopy classes
\[
 a\in\pi_{m,w}Q_AW,\qquad
 b\in\pi_{m,w}\Sigma^{1,-A}Q_BV,
 \qquad \delta_{A,B}(a)=G_B(b),
\]
there is one class $c\in\pi_{m,w}Q_N(\Sigma U)$ such that
\[
 \boxed{\rho(c)=H_A(a),\qquad F_N(c)=-\iota_{A,B}(b).}
\]
The same homotopy class $c$ solves both the restriction equation and
the signed image equation for the given $a,b$ and suspension comparison.
\uses{paper:3.20,sem:AdamsSynthetic:lambda-quotients,sem:Literature:normalized-finite-quotient-row}
\lean{KIP126.ProofTools.Internal.Mahowald.finite_may_at_normalized_triangle_explicit_lengths}
\end{lemma}
\begin{proof}
\uses{sem:Literature:normalized-finite-quotient-row,sem:Literature:infinite-finite-window,sem:ProofTools:finite-strip,paper:6.11}
Combine the distinguished normalized quotient row at length $N$ with
the coefficient triangle of lengths $A,B$. Naturality of the lambda
inclusions and restrictions gives the two commuting squares of the
finite triangle diagram. The equality $\delta_{A,B}(a)=G_B(b)$ supplies
the compatible pair for its common-lift construction. The finite-window
bounds and separation show that every composite from a sphere through
the shifted length-$B$ quotient of $V$ back to the length-$A$ quotient of
$\Sigma U$ vanishes. This removes the restriction defect.
The resulting single lift therefore has restriction $H_A(a)$ and
signed image $-\iota_{A,B}(b)$.

The quotient row initially supplies its own inclusion-compatible
suspension comparison. The same finite-window comparison identifies its
action with that of the chosen $\varepsilon_V$ on these sphere
representatives, so the second equality holds in the specified
coordinates. The comparison identifies the actions on the specified homotopy classes.
\end{proof}

\subsection{The general Mahowald rule and its core specialization}

\begin{theorem}[The finite-quotient Mahowald trick and its core form]
\label{paper:6.12}
\label{sem:ProofTools:ss-mahowald-specialization}
\source{A finite-quotient revision of AimPaper, Theorem 6.12}
\level{coarse}
\uses{sem:ProofTools:mh-actual-no-crossing,sem:ProofTools:sphere-label-detection,sem:ProofTools:mh-common-source,sem:Literature:normalized-finite-quotient-row,paper:3.20,paper:6.11,sem:ProofTools:finite-common-lift,sem:Literature:finite-quotient-labels,sem:Literature:infinite-finite-window,sem:ProofTools:level-finite-extension}
\lean{KIP126.ProofTools.Standard.generalizedMahowaldAtLevel}
\lean{KIP126.ProofTools.Standard.ssMahowaldZeroInputLevels}
\sketch

The finite-quotient form of \cite[Theorem 6.12]{LWX126} below retains
the source lambda power and the synthetic no-crossing conditions.

Fix a distinguished classical triangle $X\xrightarrow fY\xrightarrow gZ\xrightarrow h\Sigma X$,
with $e_f+e_g+e_h=1$.
Put
\[
 n_1=n-e_f\ge1,\qquad m_1=m-e_g\ge0,\qquad
 k=l-e_h\ge0,\qquad r=n+m+l.
\]
For input exponents $i\ge0$ and $0\le j<n_1$, define
\[
\begin{array}{lll}
 A=n_1+k+i-j,& B=m_1+j+1,&N=A+B=r+i,\\
 K=i+k,&J=A+j,&Q=N-K=r-k.
\end{array}
\]
Here $i,K<A$, $j<B$, and $J<N$. The parameters have the following roles:
\[
\begin{array}{c|l}
 n,m,l&\text{extension lengths for }f,g,h\text{, respectively}\\
 i,j&\text{input powers of }\lambda\\
 A,B&\text{the two input quotient lengths}\\
 N=A+B&\text{the output quotient length}\\
 K=i+k,\ J=A+j&\text{the output source and target powers of }\lambda\\
 Q=N-K&\text{the output classical cycle index}
\end{array}
\]
Use the normalized vertices
$U=\Sigma^{0,e_f}\nu X$, $V=\nu Y$, $W=\Sigma^{0,-e_g}\nu Z$.
Use the normalized triangle
$U\xrightarrow{F}V\xrightarrow{G}W\xrightarrow{H}\Sigma U$,
which is required to be distinguished.
All maps use fixed suspension--quotient isomorphisms
\[
\varepsilon_X^D:Q_D(\Sigma U)\cong\Sigma Q_DU\ (D=A,N),\qquad
\varepsilon_Y^N:Q_N(\Sigma V)\cong\Sigma Q_NV
\]
commuting with quotient inclusions. The assertion holds for every such
choice. Source labels are obtained by shifting the unshifted finite
labels through these quotient isomorphisms, as in
Definition~\ref{sem:ProofTools:level-finite-extension}.

\textit{Inputs.} Take
\[
\begin{array}{ll}
\bar x\in Z_{r-1}^{s,t}(Z),&\bar y\in Z_\infty^{s+r,t+r-1}(Z),\\
x\in Z_{n_1-j}^{s+l,t+l-1}(X),&
y\in Z_{m_1+1}^{s+l+n,t+l+n-1}(Y),
\end{array}
\]
with $d_r\bar x=\bar y$. Require the three finite extension relations
shown here:
\[
\begin{array}{ccc}
\lambda^i\bar x&\xrightarrow{\ d_l^{h_A}\ }&\lambda^Kx\\
{\scriptstyle d_r^{\delta_{A,B}}}\downarrow&&\\
\lambda^{B-1}\bar y&\xleftarrow{\ d_m^{g_B}\ }&\lambda^jy.
\end{array}
\]
\textit{Extension-relation diagram for the general Mahowald trick.}

The two relations from $\lambda^i\bar x$ use the same infinite-page
class, and the other two meet in the same target class.
The maps are the normalized $h_A$, the quotient boundary
$\delta_{A,B}:Q_AW\to\Sigma^{1,-A}Q_BW$, and
$g_B:\Sigma^{1,-A}Q_BV\to\Sigma^{1,-A}Q_BW$;
the $h_A$ target is $\Sigma Q_AU$ via the chosen comparison.
Assume actual no-crossing on either $h_A$ or $\delta_{A,B}$, in the
full range $s<u\le v\le s+l$ or $s<u\le v\le s+r$, respectively,
at stem and weight $(t-s,t-i-e_g)$ in $W$ coordinates
(equivalently $(t-s,t-i)$ before normalizing $Z$).

\textit{Output.} In the quotient of length $N$, one obtains
\[
\boxed{x\in Z_Q,\qquad
 d_n^{\hat f_N}(\lambda^Kx)=[\lambda^Jy].}
\]
The source is in $Z_Q/B_{K+1}$, and the target is in
$Z_{m_1+1}/B_{J+1}$ modulo shorter actual $\hat f_N$ extension images.
More precisely, there are representatives $a,b,c$ for which
\[
\begin{gathered}
a\text{ detects }\lambda^i\bar x,\qquad b\text{ detects }\lambda^jy,\\
\delta_{A,B}(a)=g_B(b)\text{ detects }\lambda^{B-1}\bar y,\\
\rho(c)=h_A(a),\qquad f_N(c)=-\lambda^Ab,
\end{gathered}
\]
where both $c$ and $h_A(a)$ detect the specified $\lambda^Kx$,
and $f_N(c)$ detects the specified $\lambda^Jy$.
Here $\rho$ reduces from $N$ to $A$; the equations use the same
suspension comparisons as the inputs. In these representative equations,
$c$ is transported to $\Sigma Q_NU$ and $f_N$ is the transported map
specified in the proof. A zero target label records an image in higher filtration;
the representative equations retain every power of $\lambda$.

\textit{Core specialization: $i=j=0$.}
Put $r'=n_1+k+1$. The input quotients now have lengths
$A=r'-1=n_1+k$ and $B=m_1+1$, and the output quotient has length $N=r$.
Thus the three relations have source labels $\bar x$ and $y$, common
boundary label $\lambda^{m_1}\bar y$, and $h$-image $\lambda^kx$.
The same actual no-crossing alternative on the $h$ or boundary edge is
required at $(t-s,t-e_g)$. The conclusion is
\[
 x\in Z_{r-k},\qquad d_n^{\hat f_r}(\lambda^kx)
      =[\lambda^{r'-1}y],
\]
with source denominator $B_{k+1}$, target denominator $B_{r'}$ and all
shorter extension images retained. Both input exponents are zero, and the output retains exponent $k$
as the displayed $\lambda^k$ source factor.
\end{theorem}
\begin{proof}
\textit{Compatibility of the representatives.}
For clarity, the maps in the representative equations are
\[
\begin{gathered}
h_A=\varepsilon_X^A\circ Q_AH,\qquad
\rho=\Sigma(Q_NU\longrightarrow Q_AU),\\
f_N=\varepsilon_Y^N\circ Q_N(\Sigma F)\circ(\varepsilon_X^N)^{-1}.
\end{gathered}
\]
Thus $c$ is the image under $\varepsilon_X^N$ of the lift in
$Q_N(\Sigma U)$; $\lambda^A b$ uses the suspended quotient inclusion
$\Sigma^{1,-A}Q_BV\to\Sigma Q_NV$. All detections use the fixed
identifications of ordinary suspension with bigraded shift.

Actual no-crossing gives a common representative $a$ for the two
relations leaving $\lambda^i\bar x$, by
\ref{sem:ProofTools:mh-common-source}. The target
$\lambda^{B-1}\bar y$ is in the highest retained lambda power of $Q_B$;
uniqueness there identifies $\delta_{A,B}(a)$ with $g_B(b)$.
Lemma~\ref{sem:ProofTools:finite-common-lift}, applied to these actual
representatives, gives a single $c$ with
$\rho(c)=h_A(a)$ and $f_N(c)=-\lambda^Ab$.

The chosen suspension comparisons preserve the quotient
inclusions. Their compatibility with reduction $Q_N\to Q_A$ on sphere
representatives follows from those inclusion squares and the vanishing
argument in Lemma~\ref{sem:ProofTools:finite-common-lift}: the relevant
composite through the shifted length-$B$ quotient has zero image in the
length-$A$ quotient of $\Sigma U$. This vanishing also removes the
restriction defect of $c$.
Quotient reflection and detection identify the source as $\lambda^Kx$
and its image as $\lambda^Jy$, retaining $x\in Z_Q$ and the displayed
indeterminacies, retaining every power of $\lambda$.

The relation lies in the quotient of length $N$ and retains the
source factor $\lambda^K$. The comparison below gives the source and target coordinates of the
\href{paper/main.tex.html#L1860}{article's zero-level conclusion}.

The core form takes $i=j=0$ in the displayed lengths and representative
equations. Then $K=k$, $J=r'-1$ and $Q=r-k=n+m+e_h$, giving exactly
its cycle and extension conclusions.
\end{proof}

\paragraph{What the core rule constructs.}
The three input arrows specify a common-source problem and a common-target
problem. Actual no-crossing lets one choose a representative $a$ of
$\bar x$ with both prescribed images. The top lambda power
$\lambda^{m_1}\bar y$ identifies its boundary image with $g(b)$ for a
representative $b$ of $y$. The triangle construction then produces a
single lift $c$ in the quotient of length $r$, whose restriction detects
$\lambda^kx$ and whose $f$ image detects $\lambda^{r'-1}y$.
The signed equations and suspension transports are those in the theorem.

\paragraph{Comparison with the zero-level conclusion.}
The conclusion in \cite[Theorem 6.12]{LWX126} is
\[
d_n^{f,E_{r-k+1}}(x)\equiv y\pmod{B_{r'}}.
\]
For $k>0$, the core rule retains $\lambda^kx$ in the quotient of length $r$.
For $k=0$, its output has the original zero-level shape. In both cases,
the hypotheses specify the synthetic no-crossing conditions at the chosen level.


% AimPaper: 6.12 theorem 1780 1799


\subsection{Full representatives for the nu-extension}

We first construct full synthetic representatives directly from the
independent sphere differentials. Their restrictions give the finite
nu-extension in the cofiber examples below.

\begin{proposition}[The infinite-page nu-extension used in the Leibniz example]
\label{paper:6.24}
\label{exam:stretchext}
\level{coarse}
From \cite{LWX126}, Example 6.24.\par
For $f=\nu:S^3\to S^0$, there is an extension
\[
d_2^{f,E_\infty}(h_0h_4^2)=h_0p.
\]
This means that representatives in the full synthetic spheres have the
prescribed source and image detections. Restricting these same
representatives to the quotient by $\lambda^2$ also gives the finite
relation
\[
d_2^{f,E_3}(h_0h_4^2)=h_0p.
\]
The following cofiber calculation describes this finite relation in
Atiyah--Hirzebruch coordinates.
\uses{paper:5.4}
\lean{KIP126.ProofTools.Standard.paper6_24_infinite_nu_extension}
\end{proposition}
\begin{proof}
\uses{sem:Literature:one-line-differentials,sem:Literature:two-line-classical-inputs,sem:Computation:bounded-comparison,sem:Literature:finite-bockstein,sem:ProofTools:sphere-label-detection}
\textit{Direct construction of the full representatives.}
Write $A=\Sigma^{0,1}\nu S^3$, $B=\nu S^0$ and let
$F:A\to B$ be the normalized $\nu$ map. For a quotient
boundary write
\[
\partial_q:\pi_{a,w}(Q_qV)\longrightarrow\pi_{a-1,w+q}(V).
\]
The independent inputs are $d_2(h_5)=h_0h_4^2$ and
$d_3(h_2h_5)=h_0p$. For the latter, the classical input in
\ref{sem:Literature:two-line-classical-inputs} gives a nonzero target
in degree $(5,38)$. The bounded comparison of
\ref{sem:Computation:bounded-comparison} identifies this one-dimensional
$\mathbb F_2$ space with the span of $h_0p$, so the target is $h_0p$.
The adjacent Bockstein comparison gives classes
\[
\xi\in\pi_{34,36}(Q_1A),\quad
z\in\pi_{34,36}(Q_2B),\quad
a=\partial_1\xi\in\pi_{33,37}(A),\quad
\beta=\partial_2z\in\pi_{33,38}(B).
\]
Here $a$ has the prescribed $h_0h_4^2$ source detector, transported by
the actual threefold suspension and normalization comparisons, and
$\beta$ has detector $h_0p$.

The actual first-quotient map sends $\xi$ to the class $h_2h_5$,
so that $F_1\xi=\rho_{1,2}z$. This is a comparison of operations:
quotient multiplication and its unit identify the action with
multiplication by $h_2$; the one-dimensional first-quotient source
fixes the sphere automorphisms in the suspension comparisons.
Boundary naturality and compatibility with restriction now give
\[
F(a)=\partial_1(F_1\xi)
     =\partial_1(\rho_{1,2}z)
     =\lambda\,\partial_2z=\lambda\beta.
\]
Thus a single full source and its actual image have the required
detectors. The source has filtration $3$, and the image has the
specified filtration-$5$ detector $\lambda h_0p$ at weight $37$.
This is precisely the displayed length-two infinite-page extension.
The classical page classes may vanish; the specified $E_2$ labels
and their synthetic detectors are retained throughout.

\textit{Full representatives.}
In \cite[Example 6.24]{LWX126}, the equation is obtained by applying
Corollary 6.23 at later page $\infty$, using Example 5.13. Here the two
full representatives constructed above give this zero-level extension
directly, for the displayed map and labels. Their finite restrictions
also supply the input for the cofiber calculation below;
Corollary~\ref{paper:6.23} treats the general finite-page transfer.
\end{proof}

\subsection{Cofiber examples of the Mahowald rule}

\begin{proposition}[The nu-cofiber calculation of a finite extension]
\label{paper:6.19}
\label{exam:Mahowald}
\level{coarse}
From \cite{LWX126}, Example 6.19.\par
Consider
\[
S^3\xrightarrow{\nu}S^0\xrightarrow{i}C\nu\xrightarrow{q}S^4,
\qquad (e(\nu),e(i),e(q))=(1,0,0).
\]
In Atiyah--Hirzebruch notation, $a[4]$ has top-cell image $a$ and $b[0]$ denotes the bottom-cell image of $b$. The input data are
\[
\begin{aligned}
d_0^{q,E_2}(h_0h_4^2[4])&=h_0h_4^2,\\
d_2(h_0h_4^2[4])&=h_0p[0],\\
d_0^{i,E_2}(h_0p)&=h_0p[0],
\end{aligned}
\]
where $h_0p[0]$ is a permanent cycle. The classical $d_2$ and the length-zero $q$-extension have no crossing in the sense used in the article. With
\[
n=2,\quad m=l=0,\quad n_1=1,\quad m_1=l_1=0,
\quad r=r'=2,
\]
the output is that $h_0h_4^2\in E_2^{3,36}(S^3)$ reaches $E_3$ and
\[
d_2^{\nu,E_3}(h_0h_4^2)=h_0p,
\qquad h_0p\in E_2^{5,38}(S^0).
\]
The qualification on $E_3$ allows the zero page class: the specified
$E_2$ label $h_0h_4^2$ is itself a $d_2$ boundary. The finite synthetic
extension retains this $E_2$ label and allows a zero $E_3$ representative.
\uses{paper:5.4,sem:Literature:nu-cofiber,paper:4.14}
\lean{KIP126.ProofTools.Standard.paper6_19_finite_nu_extension,KIP126.ProofTools.Standard.paper6_19_actual_cofiber_d2}
\end{proposition}
\begin{proof}
\uses{paper:6.24,paper:6.11,paper:6.12,sem:Literature:finite-bockstein,sem:ProofTools:extension-representatives}
\textit{An independent proof of the finite output.} The direct Bockstein
construction in Proposition~\ref{paper:6.24} gives one full source $a$ and
one full target representative $\beta$, with the prescribed labels and
$\widehat\nu(a)=\lambda\beta$. Project these same representatives to
$Q_2$. Naturality preserves their image equation and their finite
detectors. Since $e(\nu)=1$, the length bound on effective page $3$ is
$2\le3-2+1$, with equality. This is exactly the finite extension
displayed above, in the strong quotient of length $2$. The source
label is a permanent cycle and hence reaches $E_3$, with zero page
class allowed. This proof uses the independent sphere $d_3$ input.

\textit{The cofiber data.} They can also be recovered from the same
full representatives, without assuming the cofiber differential.
Rotate the fixed normalized $\nu$ triangle and pair it with the
actual $\lambda$-cofiber triangle. The common-lift construction
gives one class in $Q_1C\nu$ with the prescribed suspended top-cell
image and boundary equal to the bottom-cell image of the chosen
full target. Its restriction defect is the sphere composite through the
shifted finite quotient shown to vanish in
Lemma~\ref{sem:ProofTools:finite-common-lift}. The specified-lift Bockstein formula then
gives the displayed $d_2$ and proves that its bottom-cell target is
permanent. This constructs an Atiyah--Hirzebruch lift with the required top-cell label.

In \cite[Example 6.19]{LWX126}, the cofiber data come from a
secondary-Steenrod-algebra calculation. The construction above uses
the independent sphere $d_3$ input and the same full representatives.
The page-two crossing ranges are empty. Both input levels and the
output level are zero, so the finite-quotient Mahowald rule uses
quotient length $2$ throughout. Both constructions use the displayed page-two crossing checks.

For positive input level $a$, the corresponding finite target relation
is $\mathcal E^\nu_{2+a,a}(2;h_0h_4^2,h_0p)$.
At that level, the Mahowald theorem uses the input representatives
and comparison equations in the corresponding quotients and weights.
The following theorem treats a length-three extension in the displayed bidegrees.
\end{proof}

\begin{theorem}[The zero-layer nu-cofiber Mahowald instance]
\label{sem:ProofTools:nu-zero-layer}
\level{coarse}
Consider the nu-cofiber triangle
\[
S^3\xrightarrow{\nu}S^0\xrightarrow{i}C\nu\xrightarrow{q}S^4.
\]
Take labels
\[
\begin{gathered}
x\in Z_2^{8,133}(S^3),\qquad y\in Z_1^{11,136}(S^0),\\
\bar x\in Z_2^{8,134}(C\nu),\qquad
\bar y\in Z_\infty^{11,136}(C\nu).
\end{gathered}
\]
Let $\Sigma x\in E_2^{8,134}(S^4)$ be the suspended label of $x$.
Assume
\[
d_0^{q,E_3}(\bar x)=\Sigma x,\qquad
 d_3(\bar x)=\bar y,\qquad d_0^{i,E_2}(y)=\bar y.
\]
Require either that the connecting extension have no crossing in
Definition~\ref{paper:5.9}, or that $d_3(\bar x)=\bar y$ have no
crossing on $E_3$ in Definition~\ref{paper:4.14}.
Then $x\in Z_3$ and there is a nu-extension
\[
d_3^{\nu,E_4}(x)=y',\qquad y'-y\in B_3^{11,136}(S^0).
\]
\uses{sem:Literature:nu-cofiber,sem:Literature:connected-homology,sem:AdamsSynthetic:shift-coherence,sem:AdamsSynthetic:ess-shifts,sem:ProofTools:common-labels,sem:ProofTools:filtration-representatives,sem:ProofTools:extension-representatives,sem:ProofTools:sphere-label-detection,sem:ProofTools:finite-strip,paper:5.4}
\lean{KIP126.ProofTools.Internal.Mahowald.main_nu_mahowald_zero_layer}
\end{theorem}
\begin{proof}
\textit{Proof sketch.}
Use the quotient-boundary diagram of the fixed $\nu$ cofiber, with its specified suspension comparison. The classical $d_3$ supplies the boundary equation, and the two length-zero extensions supply the connecting and bottom-cell equations. The no-crossing alternative synchronizes the required representatives. In this particular finite diagram, the quotient-window bounds identify the source as a $Z_3$ label and identify the resulting length-three extension target modulo $B_3$. Thus the extension is on page $4$ with the stated congruence. This numerical construction retains the $\lambda$ powers and
comparisons of the displayed quotient-boundary diagram.
\end{proof}

\section{Numerical extensions and their crossings}

The full nu representatives and their finite restrictions have now been
constructed. The multiplication-by-two calculation has also supplied
the shorter essential extension from $d_0$. We can therefore collect
these page extensions and analyze their crossings using the definitions
already in place.

\subsection{Page extensions from the constructed representatives}

The following relations collect the preceding constructions in
page-extension notation. For each map and page, we give its extension relation and then
calculate its target coset and crossing conditions.

\begin{proposition}[Page extensions for nu and multiplication by two]
\label{paper:5.8}
\label{exam:extonEn}
\level{coarse}
From \cite{LWX126}, Example 5.8.\par
For $f=\nu:S^3\to S^0$, with $e(f)=1$, the article gives the following extensions:
\[
\begin{aligned}
d_1^{f,E_2}(h_5)&=h_2h_5,\\
d_2^{f,E_3}(h_0h_4^2)&=h_0p,\\
d_2^{f,E_\infty}(h_0h_4^2)&=h_0p.
\end{aligned}
\]
The $E_3$ relation means that a representative of $h_0h_4^2$ in $S^{3,4}/\lambda^2$ has $[h_2]$-image detecting $\lambda h_0p$ in $S^{0,0}/\lambda^2$. The $E_\infty$ relation requires representatives in the full synthetic spheres with the same specified source and image detections.

For $f=2:S^0\to S^0$, normalized by $\lambda[h_0]=2$, the article gives
\[
d_2^{f,E_\infty}(h_0h_3^2)=0,
\qquad d_1^{f,E_\infty}(d_0)=h_0d_0.
\]
The latter extension is essential and makes
\[
d_2^{f,E_\infty}(h_0h_3^2)=h_0d_0
\]
inessential. Thus two target labels can describe the same full extension coset.
\uses{paper:5.4,paper:5.7,sem:Literature:hopf-classes,sem:Literature:synthetic-h-zero}
\lean{KIP126.ProofTools.Standard.paper5_8_nu_e2_extension,KIP126.ProofTools.Standard.paper6_19_finite_nu_extension,KIP126.ProofTools.Standard.paper6_24_infinite_nu_extension,KIP126.ProofTools.Standard.paper6_8_complete_crossing_warning}
\end{proposition}
\begin{proof}
\uses{paper:6.19,paper:6.24,paper:6.8}
The first relation is the $E_2$ product with the detector $h_2$ of $\nu$, using the actual threefold suspension of $h_5$ and the normalized map on the first quotient. The second relation is the finite nu extension proved in Example~\ref{paper:6.19}. For the third, use the actual full representatives constructed in Example~\ref{paper:6.24}, with their specified source and target detections and strict normalized-map equation. These full representatives establish the infinite relation.

For multiplication by two, the full representative calculation in Example~\ref{paper:6.8} gives a representative of $h_0h_3^2$ annihilated by $[h_0]$. The shorter essential extension from $d_0$ to $h_0d_0$ puts the latter label in the length-two target indeterminacy. Adding its representative constructs the displayed nonzero-target length-two relation and proves its inessentiality.

\textit{Weight shifts of the examples.}
For $a\geq0$, the level-$a$ analogues of the $E_2$ and $E_3$
nu-extensions would be
\[
\mathcal E^\nu_{a+1,a}(1;h_5,h_2h_5),\qquad
\mathcal E^\nu_{a+2,a}(2;h_0h_4^2,h_0p).
\]
The suspension of the source labels is the same as in the original
$\nu:S^3\to S^0$ examples. The full nu-extension and the
multiplication-by-two extensions require representatives in the full
normalized map, with
source label multiplied by $\lambda^a$ and target label multiplied
by $\lambda^{a+n-e(f)}$. These equations use representatives for the full normalized map. For multiplication by two, essentiality of the transported shorter
extension is tested in its resulting target coset.
\end{proof}

\subsection{Crossings in the numerical examples}

\begin{proposition}[Crossings in the nu and multiplication-by-two examples]
\label{paper:5.13}
\label{exam:Ercross}
\level{coarse}
From \cite{LWX126}, Example 5.13.\par
The following are the crossing assertions for the maps and labels of Example~\ref{paper:5.8}, using the classical definition in \ref{paper:5.9}. Assume the displayed extension relations exist; in the fourth case assume the displayed length-one extension is essential. Their crossing conditions are as follows:
\begin{enumerate}
\item For $f=\nu:S^3\to S^0$, $d_1^{f,E_2}(h_5)=h_2h_5$ has no crossing.
\item For the same $f$, $d_2^{f,E_3}(h_0h_4^2)=h_0p$ has no crossing.
\item The corresponding $d_2^{f,E_\infty}(h_0h_4^2)=h_0p$ has no crossing.
\item For $f=2:S^0\to S^0$, the essential extension $d_1^{f,E_\infty}(d_0)=h_0d_0$ crosses both the length-two relation with target zero and the inessential length-two relation with target $h_0d_0$.
\end{enumerate}

\uses{paper:5.9}
\lean{KIP126.ProofTools.Standard.paper5_13_nu_e2_no_crossing,KIP126.ProofTools.Standard.paper5_13_nu_e3_no_crossing,KIP126.ProofTools.Standard.paper5_13_nu_infinite_no_crossing,KIP126.ProofTools.Standard.paper6_8_complete_crossing_warning}
\end{proposition}
\begin{proof}
\uses{paper:6.8}
The first relation is in the minimal-length range: its crossing definition would require $0<a\leq0$. For the second and third relations, the length bounds leave only $a=1,b=0$. The corresponding source group is generated by $h_0^2h_4^2$. The actual first-quotient normalized nu operation is multiplication by $h_2$, and $h_2h_0^2h_4^2=0$; hence that operation is zero on the entire candidate source group. In the finite case every possible image is zero, so an essential candidate is impossible. In the full case a zero first-quotient image lies one step higher in the target filtration, again contradicting essentiality.

For the fourth assertion, use the shorter essential extension constructed
in Example~\ref{paper:6.8}. Its target also satisfies
$h_0d_0\notin B_2$, by the independent nonzero classical differential
$d_3(h_0h_4)=h_0d_0$. Together with the displacement bounds in
Definition~\ref{paper:5.9}, this boundary exclusion makes the shorter
extension a crossing.
The same construction puts its target in the length-two indeterminacy
and proves the long nonzero-target relation inessential.

The synthetic argument in \cite[Example 5.13]{LWX126} uses
Proposition 5.10 to compare crossings. Division by lambda can expose
a lower target grade in the kernel of the coefficient inclusion;
Theorem~\ref{sem:ProofTools:finite-lambda-division} excludes these
grades by its target-kernel crossing condition. Applying the
level-indexed Leibniz theorem therefore requires its synthetic
no-crossing hypotheses on the stated quotient and weight, in addition
to the classical crossing assertions above.

At level $a$, the first two finite examples have quotient exponents
$a+1$ and $a+2$. Their actual crossing questions concern the maps
and weights of those quotients. The two infinite examples concern
the full map at the shifted weight. The four crossing assertions apply to the stated zero-level maps and labels.
At a specified further level, the crossing calculation uses the full
shorter-image indeterminacy in that quotient and weight.
\end{proof}

\section{Restriction, indeterminacy and finite stretching}

Quotient reduction carries an extension equation to a shorter quotient,
but it may enlarge the shorter-image indeterminacy and destroy
essentiality. The converse problem is to lift an extension while
preserving its target. We treat these two questions separately:
restriction is functorial, whereas stretching requires ruling out
nonliftable higher-filtration corrections.
The results in this section concern zero-level page extensions;
the corresponding weight and label comparisons specify positive-level coordinates.

\begin{theorem}[Descent of finite-page extensions in the valid length range]
\label{sem:ProofTools:finite-extension-descent}
\level{coarse}
The restriction statement of \cite[Proposition 6.20, first assertion]{LWX126},
within the page and length bounds of Definition~\ref{paper:5.4}.\par

Let $X,Y$ be classical spectra, $f:X\to Y$ a map, and $e=e(f)$. Take finite pages $r_{\rm early},r_{\rm late}\in\mathbb N$, integers $n,s,t$, and common labels
\[
x\in E_2^{s,t}(X),\qquad y\in E_2^{s+n,t+n}(Y).
\]
Suppose that
\[
2\leq r_{\rm early}\leq r_{\rm late},\qquad
n\leq r_{\rm early}-2+e,
\]
and that there is an extension, in the sense of Definition~\ref{paper:5.4},
\[
d_n^{f,E_{r_{\rm late}}}(x)=y.
\]
Then the same map and the same specified labels admit the earlier-page extension
\[
d_n^{f,E_{r_{\rm early}}}(x)=y.
\]
The late extension includes the lower length bound $e\leq n$ and the cycle conditions
\[
x\in Z_{r_{\rm late}-1}^{s,t}(X),\qquad
y\in Z_{r_{\rm late}-1-n+e}^{s+n,t+n}(Y).
\]
The conclusion includes the corresponding earlier-page cycle conditions. The extension need not be essential; no crossing hypothesis is imposed.

\textit{Level comparison.}
This assertion uses zero-level page extensions. At fixed positive
level $k$, the corresponding quotients would have lengths
$r_{\rm late}+k-1$ and $r_{\rm early}+k-1$, at the same weight
$t+e-k$. Positive-level restriction uses a comparison of the prescribed source
and target labels along this quotient reduction.
\uses{paper:5.4}
\lean{KIP126.ProofTools.Paper.finiteExtension_descends}
\end{theorem}
\begin{proof}
Apply quotient reduction to the source and target of the late normalized extension equation. Naturality gives the equation in the earlier quotient. The explicit upper bound $n\le r_{\rm early}-2+e$ places it in the valid length window of Definition~\ref{paper:5.4}; the late cycle conditions supply the corresponding early cycles. Reading the reduced equation gives the stated early extension, with the same map and labels.
\end{proof}


% AimPaper: 6.20 proposition 1981 1993
\begin{proposition}[Fixed-target indeterminacy after restriction]
\label{paper:6.20}
\source{A fixed-target revision of AimPaper Proposition 6.20, second assertion}
\level{coarse}
The fixed-target criterion below is related to
\cite[Proposition 6.20, second assertion]{LWX126}. It uses $b=0$
and imposes no exclusion from the later cycle set.\par
\label{prop:dec738d3}
Let $f:X\to Y$ be a map of $2$-completed, bounded-below finite-type spectra, put $e=e(f)$, and take finite pages $r',r\in\mathbb N$ and integers $n,s,t$ with
\[
2\le r'\le r,\qquad e\le n\le r'-2+e.
\]
Fix common labels
\[
x\in Z_{r-1}^{s,t}(X),\qquad y\in Z_{r-1-n+e}^{s+n,t+n}(Y),
\]
and witnesses $W_r,W_{r'}$ for $d_n^{f,E_r}(x)=y$ and
\begin{equation}\label{eq:dr'}
d_n^{f,E_{r'}}(x)=y.
\end{equation}
Both witnesses include the full cycle and length qualifications of Definition~\ref{paper:5.4}; the earlier relation is supplied by \ref{sem:ProofTools:finite-extension-descent}. Suppose that $W_r$ is essential.

Put $q=r'-1$, $A=\Sigma^{0,e}\nu X$, $B=\nu Y$, and let $F_q:A/\lambda^q\to B/\lambda^q$ be the normalized map. Retain the target class
\[
\bar y_{r'}\in E_\infty^{s+n,t+n,t+e}(B/\lambda^q)
\]
of the chosen witness $W_{r'}$, represented by the specified label $\lambda^{n-e}y$. Then the extension in (\ref{eq:dr'}) is inessential if and only if there exist
\[
a\in\mathbb N,\qquad 0<a\le n-e,\qquad
x'\in Z_{r'-1-a}^{s+a,t+a}(X),
\]
and a class $z\in E_\infty^{s+a,t+a,t+e}(A/\lambda^q)$ represented by the actual correction label $\lambda^a x'$, transported by the shift comparison and quotient isomorphism into the normalized source quotient at weight $t+e$, such that in this same extension spectral sequence
\[
d_{n-a}^{F_q}(z)=\bar y_{r'}.
\]
This is the full extension relation, including its source qualification. The shorter extension may be inessential and its source class may be zero. No exclusion from the later cycle set is imposed; the target is the original specified target, with $b=0$.

\textit{Level comparison.}
At level $k$, the same question belongs in the early quotient
$q=r'+k-1$ and weight $t+e-k$, retaining its actual target class.
The general representative criterion for inessentiality still concerns
a higher-filtration source mapping to that target. Expressing it by
classical correction labels requires the level-$k$ windows, their
boundary denominators and the comparison with the later quotient.
The displayed criterion concerns the zero-level extension with its
fixed target, so $b=0$. Requiring the shorter source to lie outside
the later cycle set leads to the two alternatives in
Theorem~\ref{sem:ProofTools:nonliftable-or-crossing}.
This is the fixed-target, zero-level criterion.
\uses{paper:5.4,sem:ProofTools:finite-extension-descent,sem:ProofTools:extension-representatives}
\lean{KIP126.ProofTools.Paper.finiteExtension_inessential_iff_shorter}
\end{proposition}
\begin{proof}
Compare \cite[Proposition 6.20]{LWX126} and its
\href{paper/main.tex.html#L1995}{proof}.

The page--length domain condition $n\le r'-2+e$ allows restriction to the earlier page. Essentiality of the late witness implies that the specified early target is nonzero. The representative criterion for inessentiality in the actual early ESS identifies an inessential extension with an extension to that same target from a source in higher filtration. The minimal-length bound gives $0<a\le n-e$. The finite-window comparison expresses this higher source by a classical label $x'\in Z_{r'-1-a}$ and its actual correction label, retaining the same quotient, stem and weight. This proves the forward implication, including the source qualification.

Conversely, the shorter extension puts the same target in the earlier
ESS's shorter-extension indeterminacy, so $W_{r'}$ is inessential.
This implication allows $x'$ to lift to the later cycle set. If it
does lift, comparison with the fixed late target can produce a late
crossing; Theorem~\ref{sem:ProofTools:nonliftable-or-crossing} retains
that alternative. The fixed-target criterion uses $b=0$ and allows sources in the later cycle set.
\end{proof}

\begin{theorem}[A nonliftable short extension or a late crossing]
\label{sem:ProofTools:nonliftable-or-crossing}
\level{coarse}
Consider finite witnesses for the same $f,x,y,n$ on pages $r_{\rm early}\leq r_{\rm late}$. If the late witness is essential and the early one is inessential, then either a shorter early extension retains the original target and has a source in the earlier cycle set but outside the later cycle set, or the same late normalized extension spectral sequence has a nonzero higher-source crossing in the specified range. The first branch uses the actual correction label with $b=0$ and all length and cycle bounds.
\uses{sem:ProofTools:extension-representatives,sem:ProofTools:finite-strip,paper:6.20}
\lean{KIP126.ProofTools.Paper.finiteExtension_nonliftableShort_or_lateCrossing}
\end{theorem}
\begin{proof}
\textit{Proof sketch.}
Use the fixed-target shorter-extension criterion in the early quotient. If its higher source has no corresponding later cycle label, it is the first alternative. Otherwise lift the source label and compare the actual image with the fixed late target. Their nonzero discrepancy has a first grade in the finite higher-source range and gives the second alternative, an actual late crossing. Thus imposing the cycle-set difference requires retaining the crossing branch.
\end{proof}

% AimPaper: 6.23 corollary 2074 2081
\begin{corollary}[Finite-page stretching with a same-quotient obstruction]
\label{paper:6.23}
\source{A finite-page sufficient-condition revision of AimPaper Corollary 6.23}
\level{coarse}
The finite-page stretching criterion below uses the same-quotient
obstruction described in \cite[Corollary 6.23]{LWX126}, with all the
shorter relations specified explicitly.\par
\label{cor:dfc6043e}
Let $f:X\to Y$ be a map of $2$-completed, bounded-below finite-type spectra, put $e=e(f)$, and let $r',r\in\mathbb N$ and $n,s,t\in\mathbb Z$ satisfy
\[
2\le r'\le r,\qquad e\le n\le r'-2+e.
\]
Suppose that the specified labels satisfy the later-page cycle conditions
\[
x\in Z_{r-1}^{s,t}(X),\qquad
y\in Z_{r-1-n+e}^{s+n,t+n}(Y),
\]
and that there is an earlier-page extension $d_n^{f,E_{r'}}(x)=y$.

Let $A=\Sigma^{0,e}\nu X$, $B=\nu Y$, and $q=r'-1$. Assume that there is no nonliftable obstruction in the same early quotient. Explicitly, exclude every pair $a,b\in\mathbb N$ with
\[
0<a\le r'-2,\qquad a+b\le n,
\]
and classes
\[
\xi\in E_\infty^{s+a,t+a,t+e}(A/\lambda^q),\qquad
0\ne\eta\in E_\infty^{s+n-b,t+n-b,t+e}(B/\lambda^q),
\]
such that $\xi$ is outside the image of the quotient reduction
\[
\rho_\infty:E_\infty^{s+a,t+a,t+e}(A/\lambda^{r-1})
\longrightarrow E_\infty^{s+a,t+a,t+e}(A/\lambda^q)
\]
and, in the same earlier-quotient ESS at stem $t-s$ and weight $t+e$,
\[
d_{n-a-b}^{\hat f_q}(\xi)=\eta.
\]
These obstructions include inessential extensions, $b=0$ and length zero; source nonliftability is defined by the actual $\rho_\infty$ image. Then the same labels admit the later-page extension
\[
d_n^{f,E_r}(x)=y.
\]
The construction supplies a later finite-page extension with the same source and target labels.

\lean{KIP126.ProofTools.Standard.finitePageExtensionStretching}
\end{corollary}
\begin{proof}
\uses{sem:ProofTools:finite-correction,sem:ProofTools:finite-strip,sem:ProofTools:extension-representatives}
Compare \cite[Corollary 6.23]{LWX126} and its
\href{paper/main.tex.html#L2074}{statement and context}.

Use actual representatives of the earlier extension. The prescribed late source and target cycle qualifications supply lifts of their infinite-page labels along quotient reduction. Correct the difference between the chosen early source and a late lift over the finite filtration interval. At each step, a residual source with nonzero target that cannot be lifted would give exactly one of the excluded same-quotient obstructions. Every other residual can be corrected to the next filtration or already has image above the prescribed target filtration. The finite correction theorem and the finite-quotient cutoff therefore supply a late representative with the specified source and target labels. Reading its actual map equation gives the required extension.

\end{proof}

\subsection{The numerical Leibniz application}

The following calculation in \cite{LWX126} uses the infinite extension
of Example~\ref{paper:6.24}. The comparison conditions for applying
\ref{sem:ProofTools:generalized-leibniz-at-level} are discussed below.

\begin{proposition}[A classical d-three obtained by the Leibniz argument]
\label{paper:6.7}
\label{exam:Leibnizyes}
\level{coarse}
From \cite{LWX126}, Example 6.7.\par
For $f=\nu:S^3\to S^0$, consider
\[
\begin{gathered}
x=h_5\in E_2^{1,35}(S^3),\qquad
y=h_2h_5\in E_2^{2,36}(S^0),\\
x_\infty=h_0h_4^2\in Z_\infty^{3,36}(S^3),\qquad
y_\infty=h_0p\in Z_\infty^{5,38}(S^0).
\end{gathered}
\]
The article applies the classical input and the two extensions
\[
d_2(h_5)=h_0h_4^2,\qquad
d_1^{f,E_2}(h_5)=h_2h_5,\qquad
d_2^{f,E_\infty}(h_0h_4^2)=h_0p
\]
with their specified no-crossing conditions to obtain
\[
d_3(h_2h_5)=h_0p.
\]
Its parameters are $n=r=2$, $m=1$, $l=2$, $e(f)=1$, $s=1$ and $t=35$.
\uses{paper:5.4,paper:4.14}
\lean{KIP126.ProofTools.Standard.paper6_7_classical_d3}
\end{proposition}
\begin{proof}
\uses{sem:Literature:two-line-classical-inputs,sem:Computation:bounded-comparison}
\textit{An independent proof of the numerical conclusion.} The classical
input in \ref{sem:Literature:two-line-classical-inputs} gives a nonzero
$d_3$ from the specified $h_2h_5$ to degree $(5,38)$. The bounded Ext
calculation identifies the unique basis vector in that degree as $h_0p$,
where $p$ is the generator in degree $(4,37)$ fixed by the sphere comparison.
Consequently every nonzero target in this one-dimensional
$\mathbb F_2$ space is $h_0p$, proving the displayed differential.
The classical input and the bounded sphere comparison prove this
numerical differential directly.

\textit{Comparison conditions for the generalized rule.} The classical differential is the Hopf-invariant-one differential. The finite extension is the product with $h_2$; the infinite extension is the assertion of Example~\ref{paper:6.24}. The article's length formula gives $r+l-m=3$.

To apply Theorem~\ref{sem:ProofTools:generalized-leibniz-at-level}
at level zero, require synthetic top-or-left no-crossing in the
quotient of length $1$ and synthetic no-crossing on the propagated
infinite bottom edge. The application uses the synthetic conditions at the displayed quotient
and weight; Example~\ref{paper:5.13} gives the classical crossing assertions.
The direct $d_3$ proof above uses the classical input and bounded comparison.

For an application of the finite-layer surjectivity composition in
Theorem~\ref{sem:ProofTools:cprime}, identify its first-quotient
connecting class with the specified $h_2h_5$, supply its two classical
differentials, and verify every layer of (Surj) in this same
$\nu$ triangle. These are the full hypotheses of the layer-surjectivity composition
criterion. Example~\ref{paper:5.13} states the classical crossing checks.
\end{proof}

\section{The finite-quotient core and classical composition}

\begin{theorem}[Finite-quotient Mahowald comparison at a specified level]
\label{sem:ProofTools:ml-finite-core}
\level{coarse}
The finite May construction of Lemma~\ref{paper:6.11}, together with
the normalized quotient rows and the Bockstein comparison, gives a
representative with an exact image in the top lambda layer.
The underlying May lemma is \ref{paper:3.20}.\par
Work with the mod-$2$ Adams spectral sequence of $2$-complete, connective,
finite-type spectra, fixed coherent synthetic and finite Bockstein comparisons,
and separated finite-quotient filtrations.
Fix one actual cofiber triangle
\[
 X\xrightarrow fY\xrightarrow gZ\xrightarrow h\Sigma X
\]
with its normalized triangle from Proposition~\ref{paper:3.20}, and write
$e_f+e_g+e_h=1$, with each exponent in $\{0,1\}$.
Absorb the classical suspensions into the three objects. The first-quotient
operations of the normalized arrows then have degrees
\[
 f_1:(e_f,e_f),\qquad g_1:(e_g,e_g),\qquad h_1:(e_h,e_h-1).
\]
When an exponent is zero this is the ordinary page map (with the displayed
desuspension for $h$). When it is one this is the normalized
first-quotient operation. They are induced by the three arrows of the cofiber triangle.
Their finite-quotient maps commute with restriction, lambda inclusion
and the quotient boundaries, with the displayed shifts.

Put $Q_nW=\nu W/\lambda^n$. Fix integers $S,T$, a level $c\geq0$ and
an integer $K\geq2$, and set $p=c+K$. Suppose
\[
 z\in Z_{K-1}^{S-e_h,T-e_h+1}(Z),\qquad
 u=h_1z\in E_2^{S,T}(X),\qquad
 v\in E_2^{S+e_f+K-1,T+e_f+K-1}(Y),
\]
and the classical differential is
\[
 d_Kz=g_1v\pmod{B_{K-1}(Z)}.
\]
The cycle qualification of its target is part of this page equation.
Define the actual homotopy groups
\[
 A=\pi_{T-S,T-c}Q_pX,\qquad
 H=\pi_{T-S+1,T-c+1-e_h}Q_pZ,\qquad
 B=\pi_{T-S,T-c+e_f}Q_pY.
\]
Write $\widehat f_p:A\to B$ and $h_p:H\to A$ for the induced normalized
maps, retaining their shifts and the fixed suspension comparison.
Let $\iota_pv\in B$ be the image of the specified first-quotient label
under the actual top-layer inclusion
$\Sigma^{0,-(p-1)}Q_1Y\to Q_pY$; its label is $\lambda^{p-1}v$.
Then there exists an actual representative $a_M\in F^SA$ such that
\[
 \operatorname{gr}^{S}(a_M)=[\lambda^cu],\qquad
 \widehat f_p(a_M)=\iota_pv.
\tag{M}
\]
In particular $u\in Z_K(X)$; no nonzero-label hypothesis is imposed.
If, in addition, the actual higher-filtration ambiguity satisfies
\[
 \widehat f_p(F^{S+1}A)=0,
\tag{A}
\]
then every specified $a\in F^SA$ with
$\operatorname{gr}^{S}(a)=[\lambda^cu]$ satisfies
\[
 \widehat f_p(a)=\iota_pv.
\tag{Cores}
\]
This is an equality in the length-$p$ quotient, with its original target.
The zero-level version is the specialization $c=0$, $p=K$.
The general level is obtained by the actual coefficient inclusion
$\Sigma^{0,-c}Q_K\to Q_{c+K}$, retaining both labels.
\uses{paper:3.20,paper:6.11,sem:AdamsSynthetic:normalized-map,sem:AdamsSynthetic:lambda-quotients,sem:Literature:normalized-finite-quotient-row,sem:Literature:infinite-finite-window}
\lean{KIP126.ProofTools.Standard.finiteQuotientMahowaldLeibnizCore}
\end{theorem}
\begin{proof}
\uses{sem:Literature:finite-bockstein,paper:4.6,sem:ProofTools:finite-common-lift,sem:Literature:normalized-finite-quotient-row,sem:Literature:infinite-finite-window,sem:ProofTools:finite-strip,sem:ProofTools:sphere-label-detection}
First construct a lift of the \emph{specified} $z$ to $Q_{K-1}Z$ whose
adjacent boundary is exactly $g_1v$ in $Q_1Z$. Finite Bockstein comparison
first supplies a lift with boundary label $w$. The input page equation
says $w-g_1v\in B_{K-1}$. Each earlier $d_r$ boundary, $2\leq r<K$,
is the adjacent boundary of a $Q_{r-1}$ representative. Its image under
$\lambda^{K-r}:Q_{r-1}\to Q_{K-1}$ has zero first-quotient restriction
and the same adjacent boundary label. Adding these correction classes
therefore changes $w$ to $g_1v$ while retaining $z$ exactly.
Equivalently, the adjacent boundary of the kernel of first-quotient
restriction has label space $B_{K-1}$. For $K=2$ this correction space
is $B_1=0$.

Apply Lemma~\ref{sem:ProofTools:finite-common-lift} to this actual boundary
square with quotient lengths $K-1$ and $1$, in the same normalized triangle.
It constructs a single $a_0$ in $Q_KX$ whose restriction is the normalized
$h$ image of the chosen lift and whose normalized $f$ image is
$\iota_Kv$. The triangle convention may put a minus sign on the latter;
it disappears on the image of the mod-$2$ first-quotient label.

The finite-window comparison and naturality for the normalized $h$
map identify the first-quotient label of $a_0$ with $h_1z=u$.
Restriction is injective on the retained bottom grade, so finite
filtration reflection puts $a_0$ in $F^S$, identifies its specified
grade, and qualifies $u$ in $Z_K$. This includes the zero-label case.
Now apply the actual coefficient inclusion
$\Sigma^{0,-c}Q_KX\to Q_{c+K}X$ to $a_0$, and call its image $a_M$.
Naturality with the same normalized arrows and the finite label
comparison carry the two labels to $\lambda^cu$ and
$\lambda^{c+K-1}v$. The latter is the highest target grade.
Its actual representative equals the single top-layer inclusion
$\iota_pv$: their difference has filtration at least
$S+e_f+K$, where the target finite quotient is zero.
This proves equality of the homotopy classes.
Hence $a_M$ has precisely the degree, leading term
and image asserted in (M). For $c=0$ this operation is the identity.
For positive $e_h$, this step uses normalized graded naturality.
A positive $e_g$ acts here on $Q_1Y$; it introduces no additional survival
requirement on $v$.

Finally, two representatives of the same associated grade differ by an
element of $F^{S+1}A$. Apply (A) to $a-a_M$ and use additivity:
$\widehat f_p(a)-\widehat f_p(a_M)=0$. Equation (M) gives the specified
image in (Cores), without changing either representative or cancelling
any power of $\lambda$.
\end{proof}

\begin{lemma}[Two finite criteria for annihilating the ambiguity]
\label{sem:ProofTools:ml-ambiguity-criteria}
\level{coarse}
A finite-filtration consequence of \ref{sem:ProofTools:finite-correction}.\par
In the notation of \ref{sem:ProofTools:ml-finite-core}, the finite-window
comparison identifies, for $0\leq j<K$,
\[
 \operatorname{gr}^{S+j}A
 =Z_{K-j}^{S+j,T+j}(X)/B_{c+j+1},\qquad F^{S+K}A=0.
\]
The highest target grade is $S+e_f+K-1$, and
$F^{S+e_f+K}B=0$.
Either of the following is sufficient for (A).

\emph{Surjective graded lifts.} For each $1\leq j<K$ and
$a\in F^{S+j}A$, there exists
$b\in F^{S+j-e_h}H$ such that
\[
 a-h_p(b)\in F^{S+j+1}A.
\]

\emph{Collapsed intermediate grades.} For $1\leq j<K-1$ the grade
$\operatorname{gr}^{S+j}A$ is zero, and the map induced by $\widehat f_p$
from the entire last grade $\operatorname{gr}^{S+K-1}A$ to
$\operatorname{gr}^{S+e_f+K-1}B$ is zero.
The latter criterion controls the image of the entire last source grade.
\uses{sem:ProofTools:ml-finite-core,sem:ProofTools:finite-strip,sem:Literature:infinite-finite-window}
\lean{KIP126.ProofTools.Standard.finiteSurjectiveAmbiguityZero,KIP126.ProofTools.Standard.finiteCollapsedAmbiguityZero}
\end{lemma}
\begin{proof}
\uses{sem:ProofTools:finite-correction}
For the first criterion, start with $a_1\in F^{S+1}A$. Choose successive
corrections $b_j$ so that $a_{j+1}=a_j-h_p(b_j)\in F^{S+j+1}A$.
After $K-1$ steps $a_K=0$, hence
$a_1=h_p(\sum_{j=1}^{K-1}b_j)$. The actual triangle gives
$\widehat f_ph_p=0$, so $\widehat f_p(a_1)=0$.
For the second criterion, the intermediate zero grades give
$F^{S+1}A=F^{S+K-1}A$. The zero map on the last grade sends this subgroup
into $F^{S+e_f+K}B$, which is zero by the finite strip and separation.
For $K=2$, the last-grade condition is the sole grade condition.
\end{proof}

\begin{theorem}[Exactness at the source of the specified $E_2$ connecting operation]
\label{sem:ProofTools:e2-boundary-exact}
\level{coarse}
The source position of the exact sequence in
\ref{sem:Literature:normalized-e2-long-exact}.\par
Fix an actual distinguished triangle in the bounded-below finite-type spectrum range
\[
X\xrightarrow fY\xrightarrow gZ\xrightarrow h\Sigma X,
\qquad e_f=1,\quad e_g=e_h=0,
\]
Its normalized synthetic triangle is distinguished by
Proposition~\ref{paper:3.20}.
Let $Q_1V=V/\lambda$. The normalized map
$\widehat f:\Sigma^{0,1}\nu X\to\nu Y$ induces the specified operation
\[
\partial_f^Q:E_2^{s,t}(X)\longrightarrow E_2^{s+1,t+1}(Y)
\]
by applying $Q_1(\widehat f)$ and the fixed first-quotient comparisons.
Let
\[
h_2:E_2^{s,t+1}(Z)\longrightarrow E_2^{s,t}(X)
\]
be the $E_2$ map of $h$ followed by the Adams desuspension comparison
fixed in Chapter~2. For every $s,t\in\mathbb Z$ and $x\in E_2^{s,t}(X)$,
\[
\partial_f^Q(x)=0
\quad\Longleftrightarrow\quad
\exists z\in E_2^{s,t+1}(Z),\quad h_2z=x.
\]
Thus $\ker\partial_f^Q=\operatorname{im}h_2$ in these precise degrees.
The desuspension on the right is the tower comparison applied to
$E_2(h)z$.
\uses{sem:AdamsSynthetic:normalized-map,sem:AdamsSynthetic:first-quotient-shifts,sem:Literature:normalized-e2-long-exact}
\lean{KIP126.ProofTools.Standard.cofiberE2Boundary_exact}
\end{theorem}
\begin{proof}
\uses{sem:Literature:normalized-e2-long-exact}
Apply the source-position clause of the normalized $E_2$ exact
sequence to the same triangle and normalization witnesses. Its connecting
operation is defined by the displayed quotient map; its image clause
uses the same $h$ map and Adams desuspension comparison.
\end{proof}

\begin{theorem}[Mahowald--Leibniz composition under finite-layer surjectivity]
\label{sem:ProofTools:cprime}
\label{sem:ProofTools:cprime-model}
\level{coarse}
\uses{sem:AdamsSynthetic:lambda-quotients,sem:AdamsSynthetic:normalized-map,sem:AdamsSynthetic:homology-exact,sem:ProofTools:e2-boundary-exact,paper:3.20,sem:Literature:normalized-finite-quotient-row,sem:AdamsSynthetic:filtration-comparison,sem:ProofTools:ml-finite-core,sem:ProofTools:ml-ambiguity-criteria,paper:4.6,sem:Literature:infinite-finite-window,sem:ProofTools:finite-strip,sem:AdamsSynthetic:first-quotient-shifts}
\lean{KIP126.ProofTools.Standard.finiteLayerSurjectivityMahowaldLeibnizE2}
\lean{KIP126.ProofTools.Standard.cofiberFirstQuotientBoundary,KIP126.ProofTools.Standard.MahowaldLeibnizFiniteLayerSurjectivity,KIP126.ProofTools.Standard.FiniteLayerSurjectivityMahowaldLeibnizStatement,KIP126.ProofTools.Standard.finiteLayerSurjectivityMahowaldLeibniz}

This theorem combines the finite Mahowald construction with the
Bockstein comparison of \cite[Theorem A.1 and Corollary A.11]{BHS}.
Its surjectivity hypothesis makes every higher-filtration correction
come from the preceding arrow of the triangle, and hence makes its
$f$ image zero.\par
Work with the mod-$2$ Adams spectral sequence of $2$-complete, connective,
finite-type spectra, with the usual coherent synthetic comparisons and
separated finite-quotient filtrations. Fix one actual cofiber triangle
\[
X\xrightarrow fY\xrightarrow gZ\xrightarrow h\Sigma X,
\qquad e_f=1,\quad e_g=e_h=0.
\]
Here $e_f=1$ means positive Adams filtration, not necessarily filtration
exactly one. Require the normalized synthetic triangle formed from
these same arrows to be distinguished. The same classical triangle then has the
homology short exact sequence
\[
0\longrightarrow H_*Y\xrightarrow{g_*}H_*Z
 \xrightarrow{h_*}\Sigma H_*X\longrightarrow0.
\]
Use the specified $E_2$ connecting operation
\[
\partial_f^Q:E_2^{s,t}(X)\longrightarrow E_2^{s+1,t+1}(Y)
\]
induced by $Q_1(\widehat f)$ and the fixed first-quotient comparisons,
as in \ref{sem:ProofTools:e2-boundary-exact}.
Write $g_2$ for the ordinary $E_2$ map of $g$, and $h_2$ for the ordinary
$E_2$ map of $h$ followed by the fixed Adams desuspension, so that
$h_2$ lowers internal degree by one. Thus $\partial_f^Q$ has degree $(1,1)$ and is induced by the
normalized arrow on the first quotient. The ordinary map $E_2(f)$
has degree $(0,0)$ and is zero here. The kernel--image relation from
homology exactness must therefore be applied to this connecting
operation, with this weight shift.

Use $Z_r,B_r\subseteq E_2$ for cycles and boundaries through $d_r$,
with $Z_1=E_2$ and $B_1=0$; all boundary labels are allowed.
Let $R,K\geq2$ be integers, $s,t\in\mathbb Z$, and put
\[
S=s+R,\qquad T=t+R-1,\qquad D=R+K-1.
\]
Suppose
\[
x\in Z_{R-1}^{s,t}(X),\qquad
z\in Z_{K-1}^{S,T+1}(Z),\qquad
v\in E_2^{S+K,T+K}(Y),
\]
with the following page equations, including qualification of
their target labels on the indicated pages:
\[
d_Rx=h_2z\pmod{B_{R-1}^{S,T}(X)},\qquad
d_Kz=g_2v\pmod{B_{K-1}^{S+K,T+K}(Z)}.
\]
For every $1\leq j\leq K-1$, assume the finite-layer surjectivity condition
\[
Z_{K-j}^{S+j,T+j}(X)
=h_2\bigl(Z_{K-j}^{S+j,T+j+1}(Z)\bigr)
 +B_{R+j-1}^{S+j,T+j}(X).
\tag{Surj}
\]
Then the specified class $y=\partial_f^Q(x)$ satisfies
\[
y\in Z_{D-1}^{s+1,t+1}(Y),\qquad
d_Dy=v\pmod{B_{D-1}^{S+K,T+K}(Y)}.
\]
In particular, if $y=0$, then $v\in B_{D-1}$.
If $v\notin B_{D-1}$, the resulting differential and the source class
on $E_D$ are nonzero.

\textit{Relation with the level-indexed rules.}
The output is a classical Adams differential; the intermediate lambda
level in the proof is fixed by $c=R-2$. The index $j$ in (Surj) runs through the specified finite set
of higher filtration grades.
Its finite Mahowald construction uses an $h$ edge of length zero,
whose crossing range $s<u\leq v\leq s$ is empty. The higher-source
ambiguity is eliminated by (Surj). The argument combines this finite
Mahowald construction with the finite-quotient May lemma.

\textit{Notation for applications.}
For a triangle $T$ in the bounded-below finite-type range with
$e_f=1$ and $e_g=e_h=0$, always take $\partial_f^Q$ to be induced by
$Q_1(\widehat f):Q_1(\Sigma^{0,1}\nu X)\to Q_1(\nu Y)$.
Write $\mathrm{Surj}(T,R,K,s,t)$ for the complete
layer condition above, and $\mathrm{ML}_{\mathrm{surj}}(T)$ for the
displayed cycle and differential conclusion for every qualified
$R,K,s,t,x,z,v$.
Then $\mathrm{ML}_{\mathrm{surj}}(T)$ holds. The output source may be zero. Applications use both input differentials,
every layer of (Surj), and the specified first-quotient boundary.
\end{theorem}
\begin{proof}
First obtain the displayed underlying graded homology short exact
sequence from the same triangle. Positive Adams filtration of $f$,
by the factorization comparison, implies $H_*f=0$.
Tensor the distinguished triangle with $H\mathbb F_2$ and rotate its
homotopy long exact sequence. In every integer degree, exactness
and the vanishing of $H_*f$ and $H_*(\Sigma f)$ show that $g_*$ is
injective, its image is the kernel of $h_*$, and $h_*$ is surjective.
The final target is $H_*(\Sigma X)$, identified with $\Sigma H_*X$
by the canonical tensor--suspension comparison; the arrows are the
actual $g$ and $h$ of the original triangle.

Set $c=R-2$ and $p=c+K=R+K-2$. The quotient boundary of the
first-quotient label of $x$, after undoing its displayed suspension, gives
\[
 a_\delta\in A=\pi_{t-s-1,t+1}Q_pX.
\]
The finite Bockstein comparison for the given $d_R$ puts it in $F^SA$
and gives the specified leading grade $[\lambda^ch_2z]$.
Theorem~\ref{sem:ProofTools:ml-finite-core}, applied to the second
classical differential, constructs $a_M$ with that same grade and
$\widehat f_pa_M=\iota_pv$.

For $1\leq j<K$, finite-window comparison identifies the actual graded
$h_p$ map with
\[
 h_2:Z_{K-j}^{S+j,T+j+1}(Z)/B_{c+j+1}
       \longrightarrow Z_{K-j}^{S+j,T+j}(X)/B_{c+j+1}.
\]
Since $c+j+1=R+j-1$, (Surj) is exactly its surjectivity.
Lift each graded preimage to an actual representative in the same
filtration of $H$. The first criterion of
\ref{sem:ProofTools:ml-ambiguity-criteria} gives (A), including the
$j=K-1$ layer with denominator $B_p$. Thus
$\widehat f_pa_\delta=\iota_pv$.

The coherent quotient-boundary square identifies the left side with
$\delta_Y(\partial^Q_fx)$. Its equality with $\iota_pv$ decodes as
follows. Reduction to $Q_{p-1}Y$ kills this top-layer image; exactness
therefore lifts $y=\partial^Q_fx$ to $Q_pY$, giving $y\in Z_p$.
For a lift with adjacent boundary label $v_0$, naturality gives
$\iota_pv_0=\iota_pv$. The kernel of the actual top-layer inclusion
on these labels is $B_p$, by the finite-window comparison and vanishing
above the top filtration. Hence $v_0-v\in B_p$, and the adjacent
Bockstein criterion gives $d_{p+1}y=v$ modulo $B_p$.
The class decoded in this argument is exactly $\partial^Q_f(x)$:
its representative was obtained by applying the same normalized
map on $Q_1$, with the specified weight shift and first-quotient
comparisons. Thus the coherent quotient-boundary square already has
the required output; no independently chosen connecting arrow is inserted.
The output bidegree is $(S+K,T+K)$, as required.
The construction supplies the representative of $y$ on the output page.
\end{proof}

\begin{theorem}[Classical composition with collapsed intermediate layers]
\label{sem:ProofTools:ml-classical-nc}
\level{coarse}
A second way to eliminate the ambiguity is to make all intermediate
source grades zero and the map on the last source grade zero.
The conditions (NC1) and (NC2) express these two requirements in
classical cycle and boundary groups.\par
Use the triangle, normalized exponents and first-quotient operations of
\ref{sem:ProofTools:ml-finite-core}, allowing all three positions of the
positive exponent. Let $R,K\geq2$, $x\in Z_{R-1}^{s,t}(X)$,
$S=s+R$, $T=t+R-1$ and $p=R+K-2$. Suppose
\[
 z\in Z_{K-1}^{S-e_h,T-e_h+1}(Z),\quad
 u=h_1z\in E_2^{S,T}(X),\quad
 v\in E_2^{S+e_f+K-1,T+e_f+K-1}(Y),
\]
where $u=h_1z$ is an equality of the specified $E_2$ labels, and
\[
 d_Rx=u\pmod{B_{R-1}(X)},\qquad
 d_Kz=g_1v\pmod{B_{K-1}(Z)}.
\]
Assume the whole-space conditions
\[
 Z_{K-j}^{S+j,T+j}(X)\subseteq B_{R+j-1}^{S+j,T+j}(X)
 \quad(1\leq j<K-1),
\tag{NC1}
\]
and
\[
 f_1\bigl(E_2^{S+K-1,T+K-1}(X)\bigr)
 \subseteq B_p^{S+e_f+K-1,T+e_f+K-1}(Y).
\tag{NC2}
\]
Then the specified $y=f_1x\in E_2^{s+e_f,t+e_f}(Y)$ satisfies
\[
 y\in Z_p(Y),\qquad d_{p+1}y=v\pmod{B_p(Y)}.
\]
For $K=2$, (NC1) is empty and (NC2) still concerns the entire final
$E_2$ source. For $e_f=1,e_g=e_h=0$, this is the branch directly
comparable with \ref{sem:ProofTools:cprime}. The other two branches use
a normalized first-quotient operation on $g$ or $h$, respectively;
each branch uses that normalized first-quotient operation.
\uses{sem:ProofTools:ml-finite-core}
\lean{KIP126.ProofTools.Standard.classicalNCComposition}
\end{theorem}
\begin{proof}
\uses{sem:ProofTools:ml-finite-core,sem:ProofTools:ml-ambiguity-criteria,paper:4.6,sem:Literature:infinite-finite-window,sem:ProofTools:finite-strip}
Take $c=R-2$ and the same specified $a_\delta$ and constructed $a_M$
as in the proof of \ref{sem:ProofTools:cprime}, now retaining the general
normalized shifts. Their difference lies in $F^{S+1}A$.
The quotient formula has denominator $B_{c+j+1}=B_{R+j-1}$, so (NC1)
annihilates the intermediate associated grades and gives
\[
 F^{S+1}A=F^{S+K-1}A.
\]
At the last layer, the source label space is $E_2/B_p$.
Normalized graded naturality identifies the map on that layer with
$f_1$ modulo the same $B_p$ in the highest target grade
$S+e_f+K-1$. By (NC2) this map is zero. Its actual image is therefore
in $F^{S+e_f+K}B=0$. This is precisely the second criterion of
\ref{sem:ProofTools:ml-ambiguity-criteria}, giving (A).
The core yields $\widehat f_pa_\delta=\iota_pv$.
The coherent boundary square and the same finite Bockstein decoding
then give $y\in Z_p$ and $d_{p+1}y=v$ modulo $B_p$.
The argument retains all three arrows, including the normalized
$g$ or $h$ operation when its exponent is positive.
\end{proof}

\begin{proposition}[The last-layer distinction between the two criteria]
\label{sem:ProofTools:ml-last-layer}
\level{coarse}
The two preceding composition theorems kill the same ambiguity in
different ways. Their last-layer conditions can be compared by the
following kernel--image calculation.\par
In their common $e_f=1,e_g=e_h=0$ branch, put
\[
 U=E_2^{S+K-1,T+K-1}(X),\quad B_X=B_p(X),\quad B_Y=B_p(Y)
\]
in the corresponding source and target bidegrees, and let $\partial=f_1$.
Exactness gives $\ker\partial=\operatorname{im}h_2$.
The last Surj condition is equivalent to
\[
 \partial U=\partial B_X,
\]
whereas (NC2) only says $\partial U\subseteq B_Y$.
Naturality gives $\partial B_X\subseteq B_Y$, so last-layer Surj implies
(NC2). Under (NC2), last-layer Surj is equivalent to lifting all the
relevant target boundaries through the source boundary group:
\[
 \operatorname{im}\partial\cap B_Y=\partial B_X.
\]
In general the difference is measured by
\[
 \frac{\partial^{-1}(B_Y)}{\operatorname{im}h_2+B_X}
 \cong\frac{\operatorname{im}\partial\cap B_Y}{\partial B_X}.
\]
For an intermediate layer, (NC1) implies the corresponding Surj
condition by taking the $h_2$ preimage to be zero. At the level of these
additive-group conditions, the reverse implication need not hold: a
nonboundary cycle can have a cycle-qualified $h_2$ preimage. These are comparisons of additive-group conditions at the specified layers.
The reverse last-layer implication uses the displayed boundary-lifting condition.
\uses{sem:ProofTools:ml-finite-core,sem:ProofTools:ml-ambiguity-criteria,sem:ProofTools:e2-boundary-exact,sem:AdamsSynthetic:first-quotient-shifts,sem:Literature:normalized-finite-quotient-row}
\lean{KIP126.ProofTools.Standard.mahowaldLeibnizLastLayerComparison,KIP126.ProofTools.Standard.mahowaldLeibnizLastLayerSurjectivity_iff,KIP126.ProofTools.intermediateLayerSurjectivity_not_converse}
\end{proposition}
\begin{proof}
\uses{sem:ProofTools:ml-finite-core,sem:ProofTools:e2-boundary-exact,sem:AdamsSynthetic:first-quotient-shifts,sem:Literature:normalized-finite-quotient-row}
Exactness of the same normalized first-quotient triangle gives
$\ker\partial=\operatorname{im}h_2$ after the fixed suspension comparisons.
The actual top-layer inclusion vanishes on a label exactly when that
label lies in $B_p$. Naturality of that inclusion under the normalized
map, together with the positive weight-shift comparison, therefore gives
$\partial B_X\subseteq B_Y$. At the last layer $Z_1=E_2$, and the
actual desuspension relation identifies its Surj condition with the
whole-group decomposition used below.

If $a=h_2z+b$ with $b\in B_X$, then $\partial a=\partial b$.
Conversely, if $\partial a=\partial b$ for such a $b$, exactness gives
$a-b=h_2z$. This proves the first equivalence. The restriction of
$\partial$ to $\partial^{-1}(B_Y)$ is onto
$\operatorname{im}\partial\cap B_Y$; taking its quotient by the
preimage of $\partial B_X$ gives the displayed isomorphism.
The intermediate-layer statement follows directly from the specified
cycle-qualified formulation of Surj. For its additive-group nonconverse,
take a nonzero group with identity $h_2$, all elements as cycles and zero
boundary subgroup. Every element has a cycle-qualified preimage, while
the cycle subgroup is not contained in the boundary subgroup. This gives an example of the distinction between the additive-group conditions.
\end{proof}


\section{Adams differential transport and classical realization of finite-level extensions}

The differential-transport rules have the shape
\[
d_bx=u,\qquad x\longmapsto y,\qquad u\longmapsto v
\quad\Longrightarrow\quad d_{b+L-M}y=v,
\]
where $M$ and $L$ are image filtration shifts. The quotient-boundary
square of generalized Leibniz, including the cutoff in
Proposition~\ref{paper:6.5}, supplies sufficient hypotheses for this
conclusion. The complementary Mahowald construction takes an $h$
preimage, a differential, and a $g$ preimage, with $r=n+l+m$.
Classical extensions additionally require full representatives and
the realization of Proposition~\ref{paper:5.14}.

A finite extension equation is understood modulo the shorter-extension
images of its map, as in
Definition~\ref{sem:ProofTools:level-finite-extension}.
To obtain a classical extension, lift its source to full synthetic
homotopy and show that the image of this lift has the required target
grade. Source permanence supplies the lift; no-crossing controls its
image; lambda inversion then gives the classical equation.
The next proposition carries out these steps. It applies to the finite-level output of the Mahowald rule
\ref{sem:ProofTools:ss-mahowald-specialization}. The composition theorem
\ref{sem:ProofTools:cprime} proves its classical differential directly
in finite quotients.

\subsection{Finite-level extensions and classical realization}


\begin{proposition}[Full target lifts, permanence and classical realization]
\label{sem:ProofTools:ss-full-lift-permanence}
\label{sem:ProofTools:ss-classical-soundness}
\level{coarse}
\uses{sem:ProofTools:finite-strip,sem:ProofTools:filtration-representatives,sem:Literature:infinite-finite-window,sem:ProofTools:level-finite-extension,sem:ProofTools:extension-level-zero,paper:5.14,sem:ProofTools:mh-actual-no-crossing,sem:Literature:infinite-nu-window}
\lean{KIP126.ProofTools.finite_target_from_full_lift_is_permanent}
\lean{KIP126.ProofTools.Standard.finiteLevelExtensionToClassical}

\textit{A full lift of the finite target.}
A finite target label is permanent whenever its detecting class lifts
to $\nu Y$ in the visible strip. The strip allows filtration to be
reflected from the quotient to the full object.

Let $N>0$, $y\in E_2^{s,t}(Y)$ and $0\le t-w<N$. Suppose an actual
full class
\[
\beta\in\pi_{t-s,w}(\nu Y)
\]
has quotient projection $q_N\beta$ detected by the specified finite
label $\lambda^{t-w}y$ in $Q_NY$. Then the original label satisfies
\[
y\in Z_\infty^{s,t}(Y).
\]
Neither nonvanishing nor prior permanence of $y$ is assumed. The second part constructs such a full target as the image of a full
source and then applies Proposition~\ref{paper:5.14}.

\textit{Realization of a finite-level extension.}
Apply this observation to the image of a full source lift. A permanent
source label supplies the lift, while no-crossing makes its finite
image have the target grade prescribed by the extension.
Use the realization isomorphisms and Adams comparison of
Proposition~\ref{paper:5.14}.

Let $f:X\to Y$, $e=e(f)$, $A=\Sigma^{0,e}\nu X$, $B=\nu Y$,
and let $F:A\to B$ be the normalized map. Take $R\ge2$,
$k\in\mathbb N$, and put
\[
q=R+k-1,\qquad p=(t-s,t+e-k).
\]
Suppose
\[
e\le m\le R-2+e,\qquad
x\in Z_\infty^{s,t}(X),\qquad
y\in E_2^{s+m,t+m}(Y),
\]
and assume the specified actual finite-level extension
\[
\mathcal E^f_{q,k}(m;x,y),\qquad
d_m^{F_q}(\lambda^kx)=\lambda^{k+m-e}y.
\]
Here the source is the intrinsic shifted finite label of
\ref{sem:ProofTools:level-finite-extension}; the equation includes
the indicated infinite-page classes and the complete shorter-image
indeterminacy of this same map.

Require $NC(F_q;p,s,m)$: there is no actual essential extension
$d_{v-u}^{F_q}(\xi)=\eta$ at this fixed stem and weight with
\[
s<u\le v\le s+m.
\]
This condition includes length zero and the terminal target
filtration $s+m$. Then
\[
y\in Z_\infty^{s+m,t+m}(Y),
\]
and there is one $a\in\pi_{t-s}(X)$ such that $x$ detects $a$ and
$y$ detects its actual image $f_*a$. Equivalently, the specified
permanent labels define the classical extension $d_m^f(x)=y$.
Neither essentiality of the given extension nor nonzero source
or target labels is required.

At $k=0$ this is a statement about a finite extension in
$Q_{R-1}$; an original $(f,E_R)$-extension supplies its finite
input by \ref{sem:ProofTools:extension-level-zero}. Positive
levels retain quotient length $R+k-1$ and weight $t+e-k$.
The proof retains the same full source representative through
the finite extension and classical realization at every level.
\end{proposition}
\begin{proof}
\textit{Full-target criterion.}
The prescribed finite detection puts $q_N\beta$ in filtration at
least $s$. Reflection along the full-to-finite projection in this
visible window puts $\beta$ in the same filtration. Its associated
grade is a full infinite-page class. Naturality and injectivity of
the finite-window comparison identify its classical label with $y$
up to the corresponding boundary group. The full label and these
boundary labels lie in $Z_\infty$, so $y$ lies there as well.
This uses the specified full class $\beta$. For the second part,
construct it as follows.

\textit{Finite extension and realization.}
\textit{The full source at the prescribed level.}
Permanence of $x$ and the full $\nu$-window comparison give
$\alpha_0\in\pi_{t-s,t}(\nu X)$ detected by $x$, whose first-quotient
label is $x$. Shift it by $(0,e)$ to obtain
$\widetilde\alpha_0\in\pi_{t-s,t+e}A$.
The zero-level source comparison identifies its detector with the
shifted label of $x$. Lambda inversion therefore gives a classical
class $a$ detected by $x$.

Set $\alpha=\lambda^k\widetilde\alpha_0$.
Lambda multiplication commutes with the weight shift. The lambda-label
comparison and the quotient-shift inclusion square identify the detector
of $q_q^A\alpha$ with the finite source label $\lambda^kx$.
Thus the level-$k$ source is obtained from the level-zero lift by the
lambda and shift maps. Since lambda multiplication preserves classical
realization, $\alpha$ still realizes the same $a$.

\textit{The actual image and its target detection.}
Choose the source filtration representative supplied by the
detection of $q_q^A\alpha$. Uniqueness of the infinite-page
representative identifies its grade with the source of the given
finite extension. Actual no crossing, including length zero and
the terminal target filtration, permits this particular source
representative while preserving the specified target grade.
Naturality of quotient projection then gives
\[
q_q^B(F_*\alpha)=F_{q*}(q_q^A\alpha)
\quad\text{detected by}\quad \lambda^{k+m-e}y.
\]
Set $\beta=F_*\alpha$. This constructs the required full target
lift from the same source. The bounds
\[
0\le k+m-e\le R+k-2=q-1
\]
place the target label in the visible finite window. By
the full-target criterion above, $y$ is permanent;
filtration and detection reflection show that $\beta$ has the
specified full target detection at weight $t+e-k$.

\textit{Classical realization.}
The arbitrary-weight $\nu Y$ comparison identifies the realization
of $\beta$ as a class detected by $y$. The realization square for
the normalized map identifies this class with $f_*a$,
where $a$ is the realization of the same $\alpha$ above.
This proves both conclusions for every $k\ge0$.
\end{proof}


The finite Mahowald theorem \ref{paper:6.12} constructs the finite
extension used in this proposition. In its notation the quotient
length and source level are $N$ and $K$, so the effective page
parameter here is $R=N-K+1$; its target level is
$J=K+n-e_f$. For realization, take a permanent source and impose the stated actual
no-crossing condition on the output map. In particular, setting
$i=j=0$ in the Mahowald theorem can still leave the output at
level $k=l-e_h>0$. The general-level realization comparison handles this weight. The finite-layer surjectivity composition
\ref{sem:ProofTools:cprime} instead proves its classical Adams
differential entirely in finite quotients.

\section{Finite comparison criteria}
\label{sec:ProofTools:computational-corollaries}

The finite-quotient constructions produce the desired labels together
with errors in higher lambda layers. The following criteria remove
these errors by controlling the subspaces in which their labels lie.
For differential transport, compare compatible finite representatives
and their coefficient boundaries. For permanence, compare a finite
representative with a full lift, and then identify the realized image.

The common-lift and Bockstein comparisons give these criteria under
the subspace hypotheses stated below. For each conclusion, the
hypotheses control the corresponding higher-layer error.

The following reading guide separates the differential and permanence
applications.
\begin{enumerate}
\item To propagate a classical differential along a triangle, start with
Corollary~\ref{sem:ProofTools:coverage-finite-zigzag-transport}. Its
page-class and selective variants allow boundary-congruent labels and
smaller spaces of possible intermediate errors.
\item When the source image is itself obtained from a finite extension,
use Corollary~\ref{sem:ProofTools:coverage-finite-source-image-transport}.
Its independently surviving-source variant permits a sharper intermediate
bound when the resulting source is known to reach the output page.
\item If a connecting image is an early boundary, the results beginning
with Corollary~\ref{sem:ProofTools:coverage-second-page-boundary-corrections}
construct compatible representatives before propagating the main
differential.
\item For permanence and the image of a prescribed classical homotopy
class, use the results beginning with
Corollary~\ref{sem:ProofTools:coverage-two-layer-permanent}. Their
middle and terminal conditions control the difference between a full
lift and a finite representative.
\item If the classical source may change in higher filtration, use
Proposition~\ref{sem:ProofTools:coverage-ess-source-adjustment}. The final
residual criterion instead retains the prescribed source by controlling
the torsion introduced during insertion into the larger quotient.
\end{enumerate}
The subspace hypotheses in these criteria apply to complete subspaces,
so an application must verify their images and kernels on the whole
indicated domain.

\subsection{Differential transport and finite source images}

\paragraph{Notation and the correction spaces.}
Work with $2$-completed, bounded-below finite-type spectra, as in
Definition~\ref{sem:AdamsSynthetic:paper-spectrum-scope}.
Fix a cofiber triangle
\[
X\xrightarrow fY\xrightarrow gZ\xrightarrow h\Sigma X,
\qquad \alpha=e_f,\quad\beta=e_g,\quad\gamma=e_h,
\qquad \alpha+\beta+\gamma=1.
\]
Write its normalized synthetic triangle as
\[
 A=\nu X,\qquad B=\Sigma^{0,-\alpha}\nu Y,\qquad
 C=\Sigma^{0,-\alpha-\beta}\nu Z,\qquad
 A\xrightarrow{\varphi}B\xrightarrow{\psi}C
    \xrightarrow{\chi}\Sigma A.
\]
For a synthetic object $D$, write $Q_qD=D/\lambda^q$ and
$\epsilon_q:D\to Q_qD$ for quotient projection. A subscript on
$\varphi,\psi,\chi$ denotes the induced quotient map, with the fixed
suspension identifications. Write $\epsilon a$ for the first-quotient
label of a full or finite homotopy class $a$: first project or restrict
to $Q_1D$, then use the appropriate classical $E_2$ comparison.
The absolute coefficient boundary and coefficient inclusion are
\[
\begin{gathered}
 \partial_q:\pi_{m,w}Q_qD\longrightarrow\pi_{m-1,w+q}D,\\
 \iota_{j,q}:\Sigma^{0,-j}Q_{q-j}D\longrightarrow Q_qD
 \qquad(0\le j<q).
\end{gathered}
\]
These conventions apply throughout the comparison criteria below,
including when $D$ is already a shifted synthetic object.

Under these degree identifications, the first-quotient operations
$F_2,G_2,H_2$ have degrees
$(\alpha,\alpha)$, $(\beta,\beta)$ and $(\gamma,\gamma-1)$,
respectively. These operations and the finite-quotient maps use the
same suspension identifications and classical Adams comparisons.
The three exponent arrangements are $(1,0,0),(0,1,0),(0,0,1)$.
Write
\[
Z_{<q}=Z_{q-1},\qquad B_{<q}=B_{q-1}\qquad(q\ge2)
\]
in the common-label notation of Definition~\ref{sem:ProofTools:common-labels}.
Thus $Z_{<q}$ is the space of labels reaching page $q$, while
$B_{<q}$ consists of actual boundaries from pages strictly below $q$.
All differential equations below are equations on the actual indicated
page; their source and target labels are allowed to represent zero.
All kernels, inverse images and inclusions concern whole subspaces.

For integers $a,b$, $q\ge2$ and $1\le j\le q-2$, set
\[
\begin{aligned}
\mathcal Z_j(a,b;q)&=Z_{<q-j+1}^{a+j,b+j}(X),\\
\mathcal L_j(a,b;q)&=B_{<j+2}^{a+j,b+j}(X)
 +H_2\!\left(Z_{<q-j+1}^{a+j-\gamma,b+j-\gamma+1}(Z)\right).
\end{aligned}
\]
Let $\mathsf M(a,b;q)$ mean that
$\mathcal Z_j(a,b;q)\subseteq\mathcal L_j(a,b;q)$ for every such $j$.
For subspaces $P_j\subseteq E_2^{a+j+\alpha,b+j+\alpha}(Y)$,
let $\mathsf M_P(a,b;q)$ mean
\[
\mathcal Z_j(a,b;q)\cap F_2^{-1}(P_j)
 \subseteq\mathcal L_j(a,b;q)
\qquad(1\le j\le q-2).
\]
These conditions control the intermediate lambda layers. At layer $j$,
the residual label lies in $\mathcal Z_j(a,b;q)$.
The condition $\mathsf M$ writes it as an earlier boundary plus an
$H_2$ image lifting through the remaining quotient. A Bockstein
correction removes the boundary, while a connecting-image correction
removes the second term without changing the image under the first
arrow. The condition $\mathsf M_P$ requires this decomposition only
when the $F_2$ image lies in $P_j$.
For $q=2$ there are no intermediate layers, so both conditions are empty.

\begin{corollary}[Finite zigzag transport: labels, page classes and selective elimination]
\label{sem:ProofTools:coverage-finite-zigzag-transport}
\level{coarse}
\uses{sem:ProofTools:common-labels,sem:ProofTools:ml-finite-core,sem:ProofTools:finite-common-lift,sem:Literature:finite-bockstein}
\lean{KIP126.ProofTools.Internal.SyntheticCoverage.differential_of_middle_images_and_candidate_kernel,KIP126.ProofTools.Internal.SyntheticCoverage.differential_of_middle_images_and_candidate_kernel_second,KIP126.ProofTools.Internal.SyntheticCoverage.differential_of_middle_images_and_candidate_kernel_third}
\lean{KIP126.ProofTools.Internal.SyntheticCoverage.differential_of_page_classes_and_candidate_kernel,KIP126.ProofTools.Internal.SyntheticCoverage.differential_of_page_classes_and_candidate_kernel_second,KIP126.ProofTools.Internal.SyntheticCoverage.differential_of_page_classes_and_candidate_kernel_third}
\lean{KIP126.ProofTools.Internal.SyntheticCoverage.differential_of_selective_middle_and_candidate_kernel_first,KIP126.ProofTools.Internal.SyntheticCoverage.differential_of_selective_middle_and_candidate_kernel_second,KIP126.ProofTools.Internal.SyntheticCoverage.differential_of_selective_middle_and_candidate_kernel_third}

\textit{Label input.}
With the preceding triangle and notation, let $r,K\ge2$, $D=r+K-1$,
$S=s+r$, $T=t+r-1$, $n=t-s$, and $b=s+\alpha+D$.
Choose
\[
\begin{gathered}
x\in Z_{<r}^{s,t}(X),\quad u\in Z_{<r}^{S,T}(X),
\quad y\in E_2^{s+\alpha,t+\alpha}(Y),\\
z\in Z_{<K}^{S-\gamma,T-\gamma+1}(Z),\quad
v\in P,\quad G_2v\in Z_{<K}(Z),\\
Z_\infty^{b,b+n-1}(Y)\subseteq P\subseteq E_2^{b,b+n-1}(Y).
\end{gathered}
\]
Suppose
\[
d_r[x]=[u],\qquad F_2x=y,\qquad F_2u=0,\qquad
H_2z=u,\qquad d_K[z]=[G_2v],
\]
and assume the complete intermediate and terminal conditions
\[
\mathsf M(S,T;K),\qquad
\ker(G_2|_P)\subseteq B_{<D}^{b,b+n-1}(Y).
\]
Then the original labels satisfy
\[
y,v\in Z_{<D}(Y),\qquad d_D[y]=[v].
\]

\textit{Page-class input.}
Retain the objects, degrees, $r,K,D$, candidate space $P$, condition
$\mathsf M(S,T;K)$ and entire terminal kernel condition of
the label-input form above.
Retain $x,u\in Z_{<r}(X)$ and $z,G_2v\in Z_{<K}(Z)$, and require
$y\in Z_{<r}^{s+\alpha,t+\alpha}(Y)$ and $v\in P$.
Replace its five zigzag relations by
\[
d_r[x]=[u],\qquad F_r[x]=[y],\qquad F_r[u]=0,
\qquad H_2z-u\in B_{<r}^{S,T}(X),\qquad d_K[z]=[G_2v].
\]
Here $F_r$ is the map of page classes induced by the same filtered map;
in common labels its source equality means $y-F_2x\in B_{<r}(Y)$.
Then
\[
y,v\in Z_{<D}(Y),\qquad d_D[y]=[v].
\]
The page-class formulation allows $z$ outside the cycles reaching
page $r$, while retaining the stated condition on $H_2z-u$.

\textit{Restricted intermediate directions.}
Retain all the objects, labels, degrees, zigzag relations and terminal
conditions of the label-input form above.
For each $1\le j\le K-2$, choose
\[
Z_\infty^{S+j+\alpha,T+j+\alpha}(Y)
 \subseteq P_j\subseteq E_2^{S+j+\alpha,T+j+\alpha}(Y).
\]
Replace the full intermediate condition $\mathsf M(S,T;K)$ by
$\mathsf M_P(S,T;K)$. Then still
\[
y,v\in Z_{<D}(Y),\qquad d_D[y]=[v],\qquad D=r+K-1.
\]
\end{corollary}
\begin{proof}
\textit{Label input.}
Write the normalized triangle as $A\xrightarrow\varphi B\to C\to\Sigma A$,
with its weight shifts understood. The specified main differential gives
one $\zeta\in Q_{r-1}A$ with first label $x$ and absolute boundary $U$
with first label exactly $u$. Earlier-boundary corrections retain $x$
while making this second label exact. The bottom differential similarly
gives a representative with adjacent boundary exactly $G_2v$.
The finite common-lift lemma then produces one $c\in Q_KA$ with first
label $u$ and $\varphi_Kc=\iota_{K-1,K}v$. The sign disappears on this
image of a mod-two first-quotient label.

The difference of $U_K$ and $c$ begins in the first lambda layer.
At layer $j$, its residual label lies in $\mathcal Z_j(S,T;K)$.
The condition $\mathsf M$ expresses it as a true boundary from pages
at most $j+1$ plus an $H_2$ image with a lift through the remaining
$K-j$ layers. The boundary has an absolute lift killed by $\lambda^j$;
the lifted $H_2$ image has zero $\varphi$ image. Subtracting them raises
the residual by one layer without changing the image of the same $U$.
At the last layer this yields
\[
\varphi_KU_K=\iota_{K-1,K}(v+F_2e).
\]
Exactness gives $\varphi U=\lambda^{K-1}V$. The coefficient pullback
applied to the same $\varphi\zeta$ and $V$ produces a length-$D-1$
representative of $y$ with boundary label $w$, hence $d_D[y]=[w]$.
The full boundary label $w$ lies in $Z_\infty(Y)$, and comparison in
$Q_KB$ gives
\[
w-v-F_2e\in B_{<K+1}(Y)\subseteq B_{<D}(Y).
\]
Since $w,v$ and true boundaries lie in $P$, so does $F_2e$.
The triangle identity $G_2F_2=0$ and the entire terminal kernel
condition imply $F_2e\in B_{<D}$, giving the asserted original target.
The same normalized construction applies to all three exponent positions.

\textit{Page-class input.}
Put $u'=H_2z$ and $y'=F_2x$. The earlier-boundary congruences preserve
the input equation $d_r[x]=[u']$, and $F_2u'=0$ follows from
$F_2H_2=0$. Apply the label-input criterion to the literal labels
$x,u',y',z,v$. It gives $d_D[y']=[v]$ and the required cycle
qualifications. Since $D>r$ and $y-y'\in B_{<r}\subseteq B_{<D}$,
the same conclusions hold for the original $y$.

\textit{Restricted intermediate directions.}
Use the same absolute boundary $U$, finite representative $c$ and
image-preserving residual induction as in the label-input proof.
At stage $j$, projection to $Q_jB$ gives
$\varphi U=\lambda^jV_j$ in the full object. Thus the first label
$\epsilon V_j$ is in $Z_\infty(Y)\subseteq P_j$.
Projection to $Q_{j+1}B$ kills the specified last-layer target and gives
\[
\epsilon V_j-F_2\epsilon e_j\in B_{<j+2}(Y).
\]
True boundaries also lie in $P_j$, so the actual residual label belongs
to $\mathcal Z_j\cap F_2^{-1}(P_j)$. The restricted inclusion therefore
performs every needed elimination. The last-layer comparison and
terminal kernel argument retain the original $y$ and $v$.
\end{proof}


\begin{corollary}[Transport through a finite source image]
\label{sem:ProofTools:coverage-finite-source-image-transport}
\level{coarse}
\uses{sem:ProofTools:coverage-finite-zigzag-transport,sem:ProofTools:common-labels,sem:ProofTools:finite-common-lift,sem:Literature:finite-bockstein}
\lean{KIP126.ProofTools.Internal.SyntheticCoverage.source_finite_differential_positive_first_selective_bottom,KIP126.ProofTools.Internal.SyntheticCoverage.source_finite_differential_ordinary_second_selective_bottom,KIP126.ProofTools.Internal.SyntheticCoverage.source_finite_differential_ordinary_third_selective_bottom,KIP126.ProofTools.Internal.SyntheticCoverage.source_finite_differential_positive_first_bottom_one,KIP126.ProofTools.Internal.SyntheticCoverage.source_finite_differential_ordinary_second_bottom_one,KIP126.ProofTools.Internal.SyntheticCoverage.source_finite_differential_ordinary_third_bottom_one}
\lean{KIP126.ProofTools.Internal.SyntheticCoverage.known_source_differential_positive_first,KIP126.ProofTools.Internal.SyntheticCoverage.known_source_differential_ordinary_second,KIP126.ProofTools.Internal.SyntheticCoverage.known_source_differential_ordinary_third,KIP126.ProofTools.Internal.SyntheticCoverage.known_source_differential_ordinary_second_top_one,KIP126.ProofTools.Internal.SyntheticCoverage.known_source_differential_ordinary_third_top_one,KIP126.ProofTools.Internal.SyntheticCoverage.known_source_differential_positive_first_bottom_one,KIP126.ProofTools.Internal.SyntheticCoverage.known_source_differential_ordinary_second_bottom_one,KIP126.ProofTools.Internal.SyntheticCoverage.known_source_differential_ordinary_third_bottom_one,KIP126.ProofTools.Internal.SyntheticCoverage.known_source_differential_ordinary_top_one_bottom_one}

\textit{General finite-source form.}
With the preceding triangle and notation, let
\[
r>K_1\ge2,\qquad K_2\ge1,\qquad D=r+K_2-K_1,
\qquad S=s+r,\quad T=t+r-1.
\]
Take $x\in Z_{<r}^{s,t}(X)$ and $u\in Z_{<r}^{S,T}(X)$,
with $d_r[x]=[u]$. Choose $z_1$ of degree
$(s-\gamma,t-\gamma+1)$ and labels $y,v$ of degrees
\[
\deg y=(s+\alpha+K_1-1,t+\alpha+K_1-1),\qquad
\deg v=(S+\alpha+K_2-1,T+\alpha+K_2-1).
\]
Assume
\[
F_2x=0,\quad H_2z_1=x,\quad d_{K_1}[z_1]=[G_2y],
\quad z_1,G_2y\in Z_{<K_1}(Z),
\]
and the complete upper conditions
\[
\mathsf M(s,t;K_1),\qquad
F_2\!\left(E_2^{s+K_1-1,t+K_1-1}(X)\right)
 \subseteq Z_{<D+1}^{\deg y}(Y).
\]
If $K_2=1$, require $v=F_2u$ and no further lower condition.
If $K_2\ge2$, choose $z_2\in Z_{<K_2}^{S-\gamma,T-\gamma+1}(Z)$
with
\[
H_2z_2=u,\qquad d_{K_2}[z_2]=[G_2v],\qquad G_2v\in Z_{<K_2}(Z).
\]
For $1\le j\le K_2-2$, choose
$Z_\infty^{S+j+\alpha,T+j+\alpha}(Y)\subseteq P_j$
inside the same $E_2$ group and require $\mathsf M_P(S,T;K_2)$.
At degree $\deg v$, choose $Z_\infty(Y)\subseteq P\subseteq E_2(Y)$
with $v\in P$ and $\ker(G_2|_P)\subseteq B_{<D}(Y)$.
Then in both cases
\[
y,v\in Z_{<D}(Y),\qquad d_D[y]=[v].
\]

\textit{Refinement with independently known source reach.}
Use the source and target degrees, main equation $d_r[x]=[u]$ and
the lower cases of
the general finite-source form above, with
$r>K_1$, $K_2\ge1$ and $D=r+K_2-K_1$.
For $K_1\ge2$, retain its full upper zigzag, $F_2x=0$,
$\mathsf M(s,t;K_1)$, and the whole correction-image inclusion into
$Z_{<D+1}(Y)$. Also allow $K_1=1$ when $\alpha=0$: in that case
replace all upper-zigzag conditions by the ordinary equality $y=F_2x$;
there is no upper correction term.

Assume independently that the original specified $y$ lies in
$Z_{<D}(Y)$. For $1\le j\le K_2-2$, put
\[
p_j=r-K_1+1+j.
\]
Replace the lower candidate bounds by arbitrary subspaces
\[
B_{<p_j}^{S+j+\alpha,T+j+\alpha}(Y)
 \subseteq P_j\subseteq E_2^{S+j+\alpha,T+j+\alpha}(Y),
\]
and require the complete restricted inclusions
$\mathsf M_P(S,T;K_2)$. Retain the lower zigzag and the entire
terminal candidate-kernel conditions when $K_2\ge2$; when $K_2=1$
retain only $v=F_2u$ on the lower side. Then
\[
y,v\in Z_{<D}(Y),\qquad d_D[y]=[v].
\]
The representative of $y$ on $E_D$ is an independent hypothesis.
\end{corollary}
\begin{proof}
\textit{General finite-source form.}
Choose one main representative $\zeta\in Q_{r-1}A$ with first label
$x$ and absolute boundary $U$ with first label $u$.
For $K_2\ge2$, the lower finite construction and selective elimination
give $\varphi U=\lambda^{K_2-1}V$ and a representative
$\eta\in Q_{r+K_2-2}B$ restricting to $\varphi\zeta$ with absolute
boundary $V$. Its boundary label $w_v$ differs from $v$ by a true
$B_{<D}$ boundary, by the entire terminal kernel condition.
For $K_2=1$, take $V=\varphi U$, $\eta=\varphi\zeta$ and $w_v=v$.

The upper zigzag supplies a common representative in $Q_{K_1}A$.
Since $r-1\ge K_1$, the full upper condition compares it to the
restriction of the same $\zeta$. The result is
\[
\eta_{K_1}=\iota_{K_1-1,K_1}(y+F_2e)
\quad\text{for }e\in E_2^{s+K_1-1,t+K_1-1}(X).
\]
Consequently $\eta$ vanishes in $Q_{K_1-1}B$. Divide by the
corresponding coefficient inclusion to obtain a representative in
$Q_{D-1}B$ with boundary $V$ and first label $w_y$ satisfying
\[
d_D[w_y]=[w_v],\qquad w_y-y-F_2e\in B_{<K_1+1}(Y).
\]
The whole image condition puts $F_2e$ in $Z_{<D+1}$, so its
$d_D$ vanishes. True boundary labels are also permanent zero labels.
The source correction therefore retains $d_D[y]=[v]$, including the cases where $K_1+1>D$.

\textit{Refinement with independently known source reach.}
Perform the upper comparison first. It produces the same divided
source $\eta_0\in Q_{r-K_1}B$ with boundary $\varphi U$ and first
label $w_y$, where $w_y-y$ is a sum of the complete upper correction
image and a true boundary. Thus $w_y\in Z_{<D}$ by the independent
assumption that $y$ reaches $E_D$ and the upper image bound. For $K_1=1$, simply
take $w_y=y$ and the ordinary image of the main source representative.

At lower stage $j$, the residual induction gives
$\varphi U=\lambda^jV_j$. The coefficient pullback with the same
$\eta_0$ produces a length-$p_j-1$ representative of $w_y$ whose
absolute boundary is $V_j$. Hence
\[
d_{p_j}[w_y]=[\epsilon V_j].
\]
Since $w_y$ reaches $E_D$ and $p_j<D$, this differential is zero, so
$\epsilon V_j\in B_{<p_j}(Y)$. The residual-image comparison
also gives $\epsilon V_j-F_2\epsilon e_j\in B_{<j+2}(Y)$.
Because $r>K_1$, one has $j+2\le p_j$, proving that the actual
residual image lies in $P_j$. The stated restricted inclusion therefore
permits every lower elimination, and the unchanged terminal argument
gives the original target. If $K_2=1$ the lower induction is empty;
when also $K_1=1$, ordinary naturality proves the conclusion directly.
\end{proof}


\begin{corollary}[Bounding new boundaries over a finite set of differentials]
\label{sem:ProofTools:coverage-finite-boundary-branches}
\level{coarse}
Fix a spectrum $X$ and a bidegree $\mathbf a$ in its actual mod-two
Adams spectral sequence. Let $\mathcal T$ be a finite set of linear maps
\[
T:Z_{<3}^{\mathbf a}(X)\longrightarrow
E_2^{\mathbf a+(3,2)}(X)/B_{<3}^{\mathbf a+(3,2)}(X)
\]
containing the quotient-valued $d_3$. Write $K_T=\ker T$.
For each $T$, let $\mathcal C_T$ be a collection of linear maps
\[
L:K_T\longrightarrow E_2^{\mathbf a+(4,3)}(X).
\]
Assume that if $T$ is the quotient-valued $d_3$ and $L$ represents
$d_4$ on every vector of $K_T$, then $L\in\mathcal C_T$.
Fix a subspace $B$ and suppose
\[
B=B_{<4}^{\mathbf a+(4,3)}(X),\qquad
\forall T\in\mathcal T\;\forall L,
\quad L\in\mathcal C_T\Longrightarrow\operatorname{im}L\subseteq B.
\]
Then
\[
B_{<5}^{\mathbf a+(4,3)}(X)
 =B_{<4}^{\mathbf a+(4,3)}(X)=B.
\]
In applications, the complete $d_2$ determines $Z_{<3}$, and ordinary
naturality supplies such collections of admissible maps. The constraints
use each branch's cycle domain $K_T$ and valid upper bounds for earlier boundaries.
\uses{sem:ProofTools:common-labels}
\lean{KIP126.ProofTools.Internal.SyntheticCoverage.standard_finite_boundary_branches}
\end{corollary}
\begin{proof}
\uses{sem:ProofTools:common-labels,sem:ProofTools:finite-strip,sem:AdamsSynthetic:e2-comparison,sem:Computation:bounded-comparison}
Let $T_0\in\mathcal T$ be the quotient-valued $d_3$, so
$K_{T_0}=Z_{<4}^{\mathbf a}(X)$.
The finite-quotient comparison and first-quotient sphere invariance
identify the scalars in degree $(0,0)$ with $\mathbb F_2$.
We may therefore lift the quotient-valued $d_4$ linearly to
\[
L_0:K_{T_0}\longrightarrow E_2^{\mathbf a+(4,3)}(X).
\]
The assumption gives $L_0\in\mathcal C_{T_0}$, and the universal image
condition gives $\operatorname{im}L_0\subseteq B$.
Since $B$ is already $B_{<4}$ in the target degree, $d_4$ contributes
no new boundary there, proving $B_{<5}=B_{<4}=B$.
The universal hypothesis concerns every admissible linear map on
each candidate's whole kernel.
\end{proof}


\subsubsection{Early boundaries and compatible representatives}

If the first image is a boundary, a preliminary
rotation of the triangle supplies a torsion correction. Before using
the main differential, we must make its representative compatible
with this early-boundary construction. The next results give conditions
for doing so without changing the prescribed final target.
Use the notation and normalized triangle of
Corollary~\ref{sem:ProofTools:coverage-finite-zigzag-transport}, including
the subspace conditions $\mathsf M$ and $\mathsf M_P$.
Continue with the normalized vertices and coefficient maps fixed at
the start of this section. Every equality involving them uses the same
degree and suspension identifications.

\begin{corollary}[Second-page boundary corrections at the target and source]
\label{sem:ProofTools:coverage-second-page-boundary-corrections}
\level{coarse}
\uses{sem:ProofTools:coverage-finite-zigzag-transport,sem:ProofTools:finite-common-lift,sem:Literature:finite-bockstein,sem:Literature:normalized-finite-quotient-row,sem:ProofTools:coverage-finite-source-image-transport}
\lean{KIP126.ProofTools.Internal.SyntheticCoverage.differential_of_early_boundary_positive_page_source,KIP126.ProofTools.Internal.SyntheticCoverage.ordinary_d2_early_target_second,KIP126.ProofTools.Internal.SyntheticCoverage.ordinary_d2_early_target_third}
\lean{KIP126.ProofTools.Internal.SyntheticCoverage.early_source_differential_positive_first,KIP126.ProofTools.Internal.SyntheticCoverage.ordinary_d2_early_source_second,KIP126.ProofTools.Internal.SyntheticCoverage.ordinary_d2_early_source_third}

\textit{Target correction.}
Let $r,L\ge3$, $D=r+L-2$, $S=s+r$, $T=t+r-1$ and $n=t-s$.
Take
\[
x\in Z_{<r}^{s,t}(X),\qquad u\in Z_{<r}^{S,T}(X),\qquad
 y\in Z_{<r}^{s+\alpha,t+\alpha}(Y),
\]
with $d_r[x]=[u]$ and $y-F_2x\in B_{<r}(Y)$.
Take
\[
c\in E_2^{S+\alpha-2,T+\alpha-1}(Y),\qquad
z=G_2c\in Z_{<L}^{S-\gamma-1,T-\gamma}(Z).
\]
At $\mathbf v=(s+\alpha+D,t+\alpha+D-1)$ choose a subspace
$Z_\infty^{\mathbf v}(Y)\subseteq P\subseteq E_2^{\mathbf v}(Y)$
and $v\in P$ with $G_2v\in Z_{<L}(Z)$. For each $1\le j\le L-3$
choose a middle candidate subspace
\[
Z_\infty^{S+j+\alpha,T+j+\alpha}(Y)\subseteq P_j
 \subseteq E_2^{S+j+\alpha,T+j+\alpha}(Y).
\]
Assume
\[
d_2c=F_2u,\qquad d_L[z]=[G_2v],
\]
\[
H_2\bigl(E_2^{S-\gamma,T-\gamma+1}(Z)\bigr)
 \subseteq B_{<3}^{S,T}(X),\qquad
\mathsf M_P(S,T;L-1),\qquad
\ker(G_2|_P)\subseteq B_{<D}^{\mathbf v}(Y).
\]
Then the originally specified labels satisfy
\[
y,v\in Z_{<D}(Y),\qquad d_D[y]=[v].
\]
The early $d_2$ equation specifies $F_2u$ as an $E_2$ boundary. When $L=3$ the middle range is empty.

\textit{Source correction.}
Let $r\ge L\ge3$ and $D=r-L+2$. Take
\[
x\in Z_{<r}^{s,t}(X),\qquad
u\in Z_{<r}^{s+r,t+r-1}(X),\qquad d_r[x]=[u],\qquad v=F_2u.
\]
Take
\[
c\in E_2^{s+\alpha-2,t+\alpha-1}(Y),\qquad
z=G_2c\in Z_{<L}^{s-\gamma-1,t-\gamma}(Z),
\]
\[
y\in E_2^{s+\alpha+L-2,t+\alpha+L-2}(Y),\qquad
G_2y\in Z_{<L}(Z).
\]
Assume
\[
d_2c=F_2x,\qquad d_L[z]=[G_2y],\qquad \mathsf M(s,t;L-1),
\]
\[
H_2\bigl(E_2^{s-\gamma,t-\gamma+1}(Z)\bigr)
 \subseteq B_{<3}^{s,t}(X),\qquad
F_2\bigl(E_2^{s+L-2,t+L-2}(X)\bigr)
 \subseteq Z_{<D+1}^{\deg y}(Y).
\]
Then $y,v\in Z_{<D}(Y)$ and $d_D[y]=[v]$.
There is no terminal candidate-space or kernel hypothesis.
\end{corollary}
\begin{proof}
\textit{Target correction.}
Apply the common lift to the rotated early equation. It gives a
representative $R\in Q_2C$ with first label $z$ and
$\chi_2R=\iota_{1,2}u$. A representative of the specified later
differential in $Q_{L-1}C$ differs from $R$ on $Q_2$ by a top-layer
class. The entire $H_2$-image condition kills its image in $Q_2A$.
Using this representative in the common-lift construction gives
$c_A\in Q_{L-1}A$ and a full element $W$ such that
\[
u'=\epsilon c_A,\qquad u'-u\in B_{<3}(X),\qquad
\lambda W=0,\qquad
\varphi_{L-1}c_A=\epsilon_{L-1}W+\iota_{L-2,L-1}v.
\]
Here $W=\partial_1\widehat c$ for a suitable $Q_1B$ class;
$\widehat c$ need not be the originally given $c$.
Since $r\ge3$, choose the main representative $\zeta\in Q_{r-1}A$
with first label $x$ and absolute boundary $U'$ with first label $u'$.
Apply restricted middle elimination to the same full error
$\Delta=\varphi U'-W$. It gives
\[
\Delta=\lambda^{L-2}V,\qquad
\epsilon V-v-F_2e\in B_{<L}(Y).
\]
Subtract $\iota_{r-2,r-1}\widehat c$ from
$\varphi_{r-1}\zeta$. Its first label stays $F_2x$ and its absolute
boundary becomes $\Delta$. Coefficient pullback gives the differential
$d_D[F_2x]=[\epsilon V]$. Both terminal labels and their boundary
error lie in $P$; $G_2F_2=0$ and the whole restricted-kernel condition
replace $\epsilon V$ by $v$ modulo $B_{<D}$. Finally
$y-F_2x\in B_{<r}\subseteq B_{<D}$ preserves the specified source.

\textit{Source correction.}
Use the two compatible lift constructions in the target-correction proof with
$x$ in place of $u$. They produce $c_A\in Q_{L-1}A$, a full element
$W$, and $x'=\epsilon c_A$ with
\[
x'-x\in B_{<3}(X),\qquad
\varphi_{L-1}c_A=\epsilon_{L-1}W+\iota_{L-2,L-1}y.
\]
The main differential still has the specified target $u$. Choose its
representative $\zeta\in Q_{r-1}A$ with first label $x'$ and absolute
boundary $U$. Since $r\ge L$, the full middle condition compares
$\zeta_{L-1}$ with $c_A$ and gives
\[
\varphi_{L-1}\zeta_{L-1}-\epsilon_{L-1}W
 =\iota_{L-2,L-1}(y+F_2e).
\]
The class $\Gamma=\varphi_{r-1}\zeta-\epsilon_{r-1}W$ has absolute
boundary $\varphi U$, since a full class has zero boundary after
projection. Coefficient exactness writes
$\Gamma=\iota_{L-2,r-1}\eta$ with $\eta\in Q_{D-1}B$.
Its first label $w$ satisfies $d_D[w]=[F_2u]$ and
$w-y-F_2e\in B_{<L}(Y)$. True boundaries have zero outgoing
differential, and the entire $F_2$ correction image reaches $E_{D+1}$
by hypothesis. Thus the conclusion holds for the original $y$.
\end{proof}


\begin{corollary}[Higher early-boundary corrections and their torsion refinement]
\label{sem:ProofTools:coverage-higher-early-boundary-corrections}
\level{coarse}
\uses{sem:ProofTools:coverage-second-page-boundary-corrections,sem:ProofTools:coverage-finite-zigzag-transport,sem:ProofTools:finite-common-lift,sem:Literature:finite-bockstein,sem:Literature:normalized-finite-quotient-row,sem:ProofTools:coverage-finite-source-image-transport}
\lean{KIP126.ProofTools.Internal.SyntheticCoverage.differential_of_higher_early_boundary_positive_page_source,KIP126.ProofTools.Internal.SyntheticCoverage.ordinary_early_target_second,KIP126.ProofTools.Internal.SyntheticCoverage.ordinary_early_target_third}
\lean{KIP126.ProofTools.Internal.SyntheticCoverage.higher_early_source_differential_positive_first,KIP126.ProofTools.Internal.SyntheticCoverage.ordinary_early_source_second,KIP126.ProofTools.Internal.SyntheticCoverage.ordinary_early_source_third}
\lean{KIP126.ProofTools.Internal.SyntheticCoverage.differential_of_torsion_higher_early_boundary_positive_page_source,KIP126.ProofTools.Internal.SyntheticCoverage.ordinary_early_torsion_target_second,KIP126.ProofTools.Internal.SyntheticCoverage.ordinary_early_torsion_target_third}

\textit{Target correction.}
Let $r>a\ge2$, $L>a$, $K=L-a+1$, $D=r+K-1$,
$S=s+r$ and $T=t+r-1$. Take
\[
x\in Z_{<r}^{s,t}(X),\quad u\in Z_{<r}^{S,T}(X),\quad
 y\in Z_{<r}^{s+\alpha,t+\alpha}(Y),\quad
 d_r[x]=[u],\quad y-F_2x\in B_{<r}(Y).
\]
The following rotation conditions concern entire subspaces. Take
\[
c\in Z_{<a}^{S+\alpha-a,T+\alpha-a+1}(Y),\qquad
 d_a[c]=[F_2u],\qquad z=G_2c\in Z_{<L}(Z),\qquad
 d_L[z]=[G_2v].
\]
The target qualifications $F_2u\in Z_{<a}$ and
$G_2v\in Z_{<L}$ are part of these page equations.
Put $K=L-a+1$ and $\mathbf z=\deg z=(S+1-\gamma-a,T+2-\gamma-a)$.
At degree $\mathbf z+(j,j)$ require, for every $1\le j\le a-2$,
\[
 Z_{<a-j+1}(Z)\subseteq B_{<j+2}(Z)
   +G_2\bigl(Z_{<a-j+1}(Y)\bigr)+Z_{<L-j+1}(Z).
\tag{R}
\]
At degree $\mathbf z+(a-1,a-1)$ require
\[
 H_2\bigl(E_2(Z)\bigr)\subseteq B_{<a+1}^{S,T}(X).
\tag{H}
\]
The source degree of each $G_2$ is determined by its fixed degree
$(\beta,\beta)$. The last summand in (R) permits adjustment of a
representative of the later differential while preserving its first label and
its prescribed differential target. At $a=2$, (R) is empty and (H)
is precisely the earlier whole $H_2$-image condition.

Impose the early and later page equations and all conditions (R) and (H). At
$\mathbf v=(s+\alpha+D,t+\alpha+D-1)$ take
\[
 Z_\infty^{\mathbf v}(Y)\subseteq P\subseteq E_2^{\mathbf v}(Y),
 \qquad v\in P,\qquad \ker(G_2|_P)\subseteq B_{<D}^{\mathbf v}(Y).
\]
For each $1\le j\le K-2$ choose
\[
Z_\infty^{S+j+\alpha,T+j+\alpha}(Y)\subseteq P_j
 \subseteq E_2^{S+j+\alpha,T+j+\alpha}(Y),
\]
and assume $\mathsf M_P(S,T;K)$ for these entire subspaces. Then
\[
y,v\in Z_{<D}(Y),\qquad d_D[y]=[v].
\]

\textit{Source correction.}
Let $L>a\ge2$, $K=L-a+1$, $r-1\ge K$ and $D=r-K+1\ge2$.
Take the original main differential
\[
x\in Z_{<r}^{s,t}(X),\qquad u\in Z_{<r}^{s+r,t+r-1}(X),
\qquad d_r[x]=[u],\qquad v=F_2u.
\]
Take
\[
c\in Z_{<a}^{s+\alpha-a,t+\alpha-a+1}(Y),\qquad
z=G_2c\in Z_{<L}(Z),\qquad
 y\in E_2^{s+\alpha+K-1,t+\alpha+K-1}(Y)
\]
with $d_a[c]=[F_2x]$ and $d_L[z]=[G_2y]$, including the target
cycle qualifications. Impose (R) and (H) with $(S,T)=(s,t)$ and
with $x,y$ in place of $u,v$. Assume
\[
\mathsf M(s,t;K),\qquad
F_2\bigl(E_2^{s+K-1,t+K-1}(X)\bigr)
 \subseteq Z_{<D+1}^{\deg y}(Y).
\]
Then $y,v\in Z_{<D}(Y)$ and $d_D[y]=[v]$.

\textit{Main-boundary torsion refinement.}
In the target-correction form above,
retain the main differential, the early $d_a$ and later $d_L$
equations, all rotation conditions, the original page source $y$,
and the entire terminal candidate and kernel conditions.
For $1\le j\le K-2$, retain the original spaces
\[
Z_\infty^{S+j+\alpha,T+j+\alpha}(Y)\subseteq P_j
\]
and choose any further subspaces $Q_j$ such that
\[
B_{<r+j+1}^{S+j+\alpha,T+j+\alpha}(Y)\subseteq Q_j.
\]
Replace the middle hypothesis by the entire inclusion
\[
Z_{<K-j+1}^{S+j,T+j}(X)\cap F_2^{-1}(P_j\cap Q_j)
\subseteq B_{<j+2}^{S+j,T+j}(X)
 +H_2\bigl(Z_{<K-j+1}^{S+j-\gamma,T+j-\gamma+1}(Z)\bigr).
\]
Then the same conclusion holds:
\[
y,v\in Z_{<D}(Y),\qquad d_D[y]=[v],\qquad D=r+K-1.
\]
In particular $a=2$ gives the second-page early-boundary version.
\end{corollary}
\begin{proof}
\textit{Target correction.}
Rotate the early $d_a$ equation to obtain $R\in Q_aC$ with
$\chi_aR=\iota_{a-1,a}u$. Compare it with the $Q_a$ restriction of
a representative $\widehat z\in Q_{L-1}C$ of the specified later
differential. At layer $j$, condition (R) decomposes the error label
into a true early boundary, a $G_2$ image, and a class in
$Z_{<L-j+1}$. The boundary has a full lift annihilated by
$\lambda^j$, and the $G_2$ image is killed by $\chi$.
For the last summand choose a lift in $Q_{L-j}C$ and first restrict
it to $Q_{L-j-1}C$ before inserting it into $\widehat z$.
This correction has zero adjacent boundary, so it changes neither
$z$ nor the prescribed later target $G_2v$. At the last layer (H)
kills the remaining connecting image. The adjusted representative
therefore satisfies $\chi_a\widehat z_a=\iota_{a-1,a}u$.

The common lift of that representative, followed by coefficient
exactness, produces
\[
c_A\in Q_KA,\quad u'=\epsilon c_A,\quad
u'-u\in B_{<a+1}(X),\quad \lambda^{a-1}W=0,\quad
\varphi_Kc_A=\epsilon_KW+\iota_{K-1,K}v,
\]
where $W=\partial_{a-1}\widehat c$ is a full absolute boundary.
Since $r>a$, the main differential still has target $u'$.
For its same representative $\zeta\in Q_{r-1}A$ and absolute
boundary $U'$, restricted middle elimination gives
\[
\varphi U'-W=\lambda^{K-1}V,\qquad
\epsilon V-v-F_2e\in B_{<K+1}(Y).
\]
Subtract $\iota_{r-a,r-1}\widehat c$ from
$\varphi_{r-1}\zeta$. The strict inequality $r>a$ preserves the
first label. Coefficient pullback then produces
$d_D[F_2x]=[\epsilon V]$. The terminal kernel condition replaces
$\epsilon V$ by $v$, since $B_{<K+1}\subseteq B_{<D}$, and the
original congruence replaces $F_2x$ by $y$.

\textit{Source correction.}
The compatible-representative construction in the target-correction proof
produces $c_A\in Q_KA$, a full $W$, and $x'=\epsilon c_A$ with
\[
x'-x\in B_{<a+1}(X),\qquad
\varphi_Kc_A=\epsilon_KW+\iota_{K-1,K}y.
\]
A true boundary supports zero outgoing differential on every page:
before its killing page it is a cycle, on that page use $d^2=0$,
and after that page its class is zero. Thus $d_r[x']=[u]$ still
holds for $x'$, including when the boundary is not earlier than $r$.
Choose the corresponding main representative $\zeta\in Q_{r-1}A$
and its specified absolute boundary $U$. Since $r-1\ge K$, the
full middle condition gives
\[
\varphi_K\zeta_K-\epsilon_KW=\iota_{K-1,K}(y+F_2e).
\]
Subtracting the projection of the full class $W$ preserves
the absolute boundary. Exactness therefore yields
\[
\varphi_{r-1}\zeta-\epsilon_{r-1}W
 =\iota_{K-1,r-1}\eta,\qquad
\eta\in Q_{D-1}B,\qquad \partial_{D-1}\eta=\varphi U.
\]
Its first label $w$ satisfies $d_D[w]=[F_2u]$ and
$w-y-F_2e\in B_{<K+1}(Y)$. The true boundary term and the entire
correction image in $Z_{<D+1}$ have zero $d_D$, proving the
assertion for the originally specified $y$ and $v$.

\textit{Main-boundary torsion refinement.}
The same specified main absolute boundary obeys
$\lambda^{r-1}U'=0$. The correction $W$ obeys
$\lambda^{a-1}W=0$, and $a<r$. Hence the actual full error
$\Delta=\varphi U'-W$ is annihilated by $\lambda^{r-1}$.
At the $j$th elimination step write $\Delta=\lambda^jV_j$.
Then $\lambda^{r+j-1}V_j=0$, so coefficient Bockstein comparison
gives $\epsilon V_j\in B_{<r+j+1}(Y)$.
The difference between this first label and $F_2\epsilon e_j$
lies in $B_{<j+2}(Y)$, which is contained in the same boundary
space. Thus the actual middle error lies in $P_j\cap Q_j$.
The original induction applies to this smaller inverse image;
its synchronized source, prescribed target and final kernel
comparison remain unchanged.
\end{proof}


\subsubsection{A complete relation at the bottom three layers}

An intermediate residual need not be a boundary or a connecting image.
The next criterion allows a remaining subspace $\mathcal K_0$, provided
each of its elements occurs in a compatible two-layer relation.
The projected target of that relation contributes a final error, which
is controlled together with the ordinary last-layer image. Surjectivity
of the full relation supplies one compatible triple for each residual.

\begin{corollary}[A bottom three-layer relation with independently known page reach]
\label{sem:ProofTools:coverage-bottom-recursive-differential}
\level{coarse}

Fix $d_r[x]=[u]$ with $x\in Z_{<r}^{s,t}(X)$ and
$u\in Z_{<r}^{S,T}(X)$, where $S=s+r$ and $T=t+r-1$.
Let $K_1\ge1$, $r>K_1$, $D=r+3-K_1\ge4$, and $p=D-1$.
For $K_1=1$ put $y=F_2x$, using the normalized first-quotient
operation also when $\alpha=1$. For $K_1\ge2$ impose the full upper
conditions of Corollary~\ref{sem:ProofTools:coverage-finite-source-image-transport}:
\[
F_2x=0,\quad H_2z_1=x,\quad d_{K_1}[z_1]=[G_2y],\quad
z_1,G_2y\in Z_{<K_1}(Z),\quad \mathsf M(s,t;K_1),
\]
\[
F_2\bigl(E_2^{s+K_1-1,t+K_1-1}(X)\bigr)
 \subseteq Z_{<D+1}^{\deg y}(Y).
\]
In both cases $\deg y=(s+\alpha+K_1-1,t+\alpha+K_1-1)$.
Require the lower relations
\[
H_2z_2=u,\qquad d_3[z_2]=[G_2v],\qquad
z_2,G_2v\in Z_{<3}(Z),
\]
Here $\deg z_1=(s-\gamma,t-\gamma+1)$ when $K_1\ge2$,
and $\deg z_2=(S-\gamma,T-\gamma+1)$. Independently assume
that the original output source reaches $E_D$, namely
$y\in Z_{<D}(Y)$.
Put
\[
\mathbf x=(S+1,T+1),\quad
\mathbf z=(S+1-\gamma,T+2-\gamma),\quad
\mathbf y=(S+1+\alpha,T+1+\alpha),\quad
\mathbf v=(S+2+\alpha,T+2+\alpha).
\]
Choose subspaces
\[
Z_\infty^{\mathbf y}(Y)\subseteq P_{\rm mid},\qquad
B_{<p}^{\mathbf y}(Y)\subseteq\overline B_{<p}^{\mathbf y}(Y)
 \subseteq E_2^{\mathbf y}(Y),
\]
and define
\[
\mathcal K=\ker(d_2:E_2^{\mathbf x}(X)\to E_2(X))
 \cap F_2^{-1}(P_{\rm mid}\cap\overline B_{<p}(Y)).
\]
Assume there is a subspace
$\mathcal K_0\subseteq\ker d_2^{\mathbf x}(X)\cap\ker F_2$
with
\[
\mathcal K\subseteq B_{<p}^{\mathbf x}(X)
 +H_2(\ker d_2^{\mathbf z}(Z))+\mathcal K_0.
\]
Use the complete linear relation
\[
\mathcal R=\{(\omega,k,z):\ k\in\mathcal K_0,\quad
 H_2\omega=k,\quad d_2\omega=G_2z\},
\]
where $\omega\in E_2^{\mathbf z}(Z)$ and
$z\in E_2^{\mathbf v}(Y)$, and require
$\mathcal R\to\mathcal K_0$ to be surjective. Finally choose
\[
Z_\infty^{\mathbf v}(Y)\subseteq P_{\rm fin}\subseteq E_2^{\mathbf v}(Y),
\qquad v\in P_{\rm fin},
\]
with the whole-space condition
\[
\left(\operatorname{pr}_z\mathcal R+
 F_2(E_2^{S+2,T+2}(X))\right)\cap P_{\rm fin}
 \subseteq B_{<D}^{\mathbf v}(Y).
\]
Then the original labels satisfy $y,v\in Z_{<D}(Y)$ and
$d_D[y]=[v]$. The hypotheses specify a representative on $E_D$ and a surjective
complete linear relation on the entire residual space.
\uses{sem:ProofTools:coverage-finite-source-image-transport,sem:ProofTools:coverage-finite-zigzag-transport}
\lean{KIP126.ProofTools.Internal.SyntheticCoverage.recursive_differential_ordinary_third_top_one,KIP126.ProofTools.Internal.SyntheticCoverage.recursive_differential_ordinary_second_top_one,KIP126.ProofTools.Internal.SyntheticCoverage.recursive_differential_positive_first_top_one,KIP126.ProofTools.Internal.SyntheticCoverage.recursive_differential_ordinary_third_finite_top,KIP126.ProofTools.Internal.SyntheticCoverage.recursive_differential_ordinary_second_finite_top,KIP126.ProofTools.Internal.SyntheticCoverage.recursive_differential_positive_first_finite_top}
\end{corollary}
\begin{proof}
\uses{sem:ProofTools:finite-common-lift,sem:ProofTools:coverage-finite-source-image-transport,sem:Literature:finite-bockstein}
First keep the main representative $\zeta\in Q_{r-1}A$ and its
absolute boundary $U$ synchronized. The upper comparison gives
\[
\varphi_{r-1}\zeta=\iota_{K_1-1,r-1}\eta_0,\qquad
\eta_0\in Q_{r-K_1}B,\qquad \partial\eta_0=\varphi U.
\]
For $K_1=1$ its first label $w_y$ is $y$; otherwise
$w_y-y-F_2e_{\rm top}\in B_{<K_1+1}(Y)$.
The hypothesis that $y$ reaches $E_D$, the complete correction-image
condition, and zero outgoing differentials on true boundaries imply
$w_y\in Z_{<D}$.

The common lift for the lower relations gives $c_3\in Q_3A$ with first label $u$
and $\varphi_3c_3=\iota_{2,3}v$. Write
$U_3-c_3=\iota_{1,3}e_2$ and $\varphi U=\lambda V_1$.
The synchronized boundary pullback gives
$d_p[w_y]=[\epsilon V_1]$. As $w_y$ reaches $E_D$, where $D=p+1$,
$\epsilon V_1\in B_{<p}(Y)$. The coefficient-kernel comparison
places $F_2\epsilon e_2-\epsilon V_1$ in $B_{<3}(Y)$, so the
actual error label belongs to $\mathcal K$.

Absorb each true boundary component $b=d_q[a]$, $q<p$, by
simultaneously replacing $\zeta$ and $U$ by
$\zeta-\zeta_b$ and $U-\lambda T_b$, where
$\partial\zeta_b=\lambda T_b$ and $\epsilon T_b=b$.
The correcting representative lies in layer at least $r+1-q$.
This preserves the original first labels and, when $K_1\ge2$,
the entire protected upper quotient $Q_{K_1}$.
Absorb the $H_2(\ker d_2)$ component through its two-layer
connecting lift. The remaining error label is one actual
$k\in\mathcal K_0$.

Surjectivity supplies a single $(\omega,k,z)\in\mathcal R$.
Its common two-layer lift compares with $e_2$ to give
\[
\varphi_3U_3=\iota_{2,3}(v+z+F_2e),\qquad
\varphi U=\lambda^2V.
\]
The same upper division and main boundary pullback produce
$d_D[w_y]=[w]$, where $w=\epsilon V\in Z_\infty(Y)$ and
$w-v-z-F_2e\in B_{<4}(Y)\subseteq B_{<D}(Y)$.
The complete terminal intersection condition puts $z+F_2e$
in $B_{<D}(Y)$. Thus $w$ can be replaced by the originally
specified $v$. The preserved upper correction has zero $d_D$,
so $w_y$ can likewise be replaced by the originally specified $y$.
\end{proof}


\subsection{Permanence and the image of a specified classical source}

To determine the image of a given classical class $\xi$, begin with a
full synthetic lift of $\xi$ and compare it with the finite lift supplied
by the differential data. Remove the difference one lambda layer at a
time. Full lambda-power-torsion corrections preserve $\xi$, while
connecting-image corrections preserve the finite lift's image.
The last-layer comparison then determines the filtration and associated
grade of $f_*\xi$.

Unless stated otherwise, use the compatible triangle and first-quotient
operations $F_2,G_2,H_2$ of this section, with
$\alpha+\beta+\gamma=1$. In shifted coordinates, write
\[
 A=\nu X,\qquad B=\Sigma^{0,-\alpha}\nu Y,\qquad
 C=\Sigma^{0,-\alpha-\beta}\nu Z,\qquad
 A\xrightarrow{\varphi}B\longrightarrow C\longrightarrow\Sigma A.
\]
The first-quotient operations have degrees $(\alpha,\alpha)$,
$(\beta,\beta)$ and $(\gamma,\gamma-1)$. Write
$Q_jD=D/\lambda^j$, $\epsilon_j:D\to Q_jD$,
$\epsilon:Q_jD\to Q_1D$ for the indicated first restriction, and
$\iota_{k,N}:\Sigma^{0,-k}Q_{N-k}D\to Q_ND$ for coefficient inclusion.

Here $B_\infty$ always means all true Adams boundaries and
$Z_\infty$ includes them. A label $x\in E_2^{s,t}(X)$ strictly detects
$\xi\in\pi_{t-s}X$ when $x$ is a nonboundary permanent label and its
infinite-page class is the associated grade of $\xi$. Statements that
a target $v$ is permanent permit $[v]_\infty=0$. The equations
\[
 f_*\xi\in F^p\pi_nY,\qquad
 \operatorname{gr}^{p}(f_*\xi)=[v]_\infty
\]
mean strict detection when $v\notin B_\infty$ and membership in
$F^{p+1}\pi_nY$ when $v\in B_\infty$. Realization and detection
use the Adams comparisons of Proposition~\ref{paper:5.14}.

\begin{corollary}[Permanence from a two-layer fold]
\label{sem:ProofTools:coverage-two-layer-permanent}
\level{coarse}
Put $n=t-s$, assume $s\ge\gamma$, and set $p=s+\alpha+1$.
Let $x\in Z_\infty^{s,t}(X)$ strictly detect a classical class, and
suppose $F_2x=0$. Choose
\[
 U\in E_2^{s-\gamma,t-\gamma+1}(Z),\qquad
 Z_\infty^{p,p+n}(Y)\subseteq P\subseteq E_2^{p,p+n}(Y),
 \qquad v\in P.
\]
If
\[
 H_2U=x,\qquad d_2U=G_2v,\qquad
 \ker(G_2|_P)\subseteq B_\infty^{p,p+n}(Y),
\]
then $v\in Z_\infty^{p,p+n}(Y)$ and, for every classical
$\xi\in\pi_nX$ strictly detected by $x$,
\[
 f_*\xi\in F^p\pi_nY,\qquad
 \operatorname{gr}^{p}(f_*\xi)=[v]_\infty.
\]
The $d_2$ equation is an equality of full $E_2$ vectors and may be zero.
The kernel condition concerns the entire subspace $P$.
\uses{sem:ProofTools:common-labels,sem:Literature:lambda-inversion-adams-comparison}
\lean{KIP126.ProofTools.Internal.SyntheticCoverage.two_layer_permanence_positive_h,KIP126.ProofTools.Internal.SyntheticCoverage.two_layer_permanence_positive_g,KIP126.ProofTools.Internal.SyntheticCoverage.two_layer_permanence_positive_f}
\end{corollary}
\begin{proof}
\uses{sem:ProofTools:finite-common-lift,sem:ProofTools:finite-strip,sem:Literature:finite-bockstein,sem:Literature:complete-permanent-lifts,sem:Literature:lambda-inversion-detection}
The common-lift construction gives $c\in\pi_{n,t}Q_2A$ with
$\epsilon c=x$ and $\varphi_2c=\iota_{1,2}v$.
Fix the prescribed classical class $\xi$ and take its full synthetic
lift $a$ with first label $x$. Exactness gives
$\epsilon_2a-c=\iota_{1,2}e$, with
$e\in E_2^{s+1,t+1}(X)$.
Since $F_2x=0$, write $\varphi a=\lambda b$ and put $w=\epsilon_1b$.
The full-object comparison makes $w$ permanent and identifies
$[w]_\infty$ with the grade of this same $f_*\xi$ at filtration $p$.
Comparison in $Q_2B$ gives
$w-v-F_2e\in B_{<3}(Y)$. As $P$ contains $w,v$ and all true
boundaries, $F_2e\in P$; the triangle identity $G_2F_2=0$ and the
entire-kernel hypothesis imply $w-v\in B_\infty(Y)$.
Thus $v$ is permanent with the required grade. No nonzero target
assumption is used.
\end{proof}

\begin{corollary}[Correcting a boundary first image]
\label{sem:ProofTools:coverage-boundary-lift}
\level{coarse}
For this statement let $\alpha,\beta\ge0$ be arbitrary integers.
It suffices to have actual synthetic maps
\[
 F:\Sigma^{0,\alpha}\nu X\longrightarrow\nu Y,\qquad
 G:\Sigma^{0,\beta}\nu Y\longrightarrow\nu Z,
 \qquad G\circ\Sigma^{0,\beta}F=0,
\]
whose first-quotient operations are $F_2,G_2$.
For the given classical map $f:X\to Y$, require the actual realization
square
\[
 \theta_Y\circ L(F)=f\circ\theta_{X,\alpha},
\]
where $L$ inverts $\lambda$ and $\theta_Y:L(\nu Y)\cong Y$,
$\theta_{X,\alpha}:L(\Sigma^{0,\alpha}\nu X)\cong X$ are the
specified standard Adams comparison isomorphisms. A compatible
normalized triangle supplies these hypotheses, but no condition on
the third exponent or their sum is imposed here.

Let $r\ge2$, $q=r+1$, $n=t-s$, $p=s+\alpha+1$, and let
$x\in Z_\infty^{s,t}(X)$ strictly detect a classical class. Assume
\[
 a\in Z_{<r}^{s+\alpha-r,t+\alpha-r+1}(Y),\qquad
 d_r[a]=[F_2x],\qquad
 U=G_2a\in Z_{<q}^{s+\alpha+\beta-r,t+\alpha+\beta-r+1}(Z).
\]
Take $Z_\infty^{p,p+n}(Y)\subseteq P\subseteq E_2^{p,p+n}(Y)$
and $v\in P$, with $G_2v\in Z_{<q}(Z)$. Suppose
\[
 d_q[U]=[G_2v],\qquad
 \{z\in P:G_2z\in B_{<q}(Z)\}\subseteq B_\infty(Y).
\]
Both differential equations use their actual pages and all true earlier
boundaries. Then $v$ is permanent and, for every $\xi\in\pi_nX$
strictly detected by $x$,
\[
 f_*\xi\in F^p\pi_nY,\qquad
 \operatorname{gr}^{p}(f_*\xi)=[v]_\infty.
\]
\uses{sem:ProofTools:common-labels,sem:Literature:lambda-inversion-adams-comparison}
\lean{KIP126.ProofTools.Internal.SyntheticCoverage.boundary_lift_prescribed_source_general}
\end{corollary}
\begin{proof}
\uses{sem:Literature:finite-bockstein,sem:Literature:complete-boundary-lift,sem:Literature:complete-permanent-lifts,sem:Literature:lambda-inversion-detection}
In coordinates $A=\nu X$, $B=\Sigma^{0,-\alpha}\nu Y$ and
$C=\Sigma^{0,-\alpha-\beta}\nu Z$, let $\varphi,\psi$ be the
corresponding maps. Lift the specified $a$ to
$\zeta\in\pi_{n+1,t-r+1}Q_{r-1}B$ with adjacent boundary exactly
$F_2x$. Earlier-boundary corrections make the target exact while
retaining the first label. Its absolute boundary $T$ satisfies
$\epsilon_1T=F_2x$ and $\lambda^{r-1}T=0$.
For a full lift $\widetilde\xi$ of the prescribed $\xi$, exactness
gives $\varphi\widetilde\xi-T=\lambda b$.
The label $w=\epsilon_1b$ is permanent and realizes the same image
$f_*\xi$, by the stated realization square.

Since $\psi\varphi=0$, the finite image $\psi_{r-1}\zeta$ extends
to $\eta\in Q_rC$ with first label $U$. Naturality of the absolute
boundaries gives $\lambda(\psi b+\partial_r\eta)=0$.
The full coefficient triangle therefore gives a class $e\in Q_1C$
with $\psi b+\partial_r\eta=\partial_1e$; no cancellation of
$\lambda$ is used. On first labels,
$G_2w+\delta_r\eta=d_2e$. The prescribed $d_q$ equation implies
$G_2(w-v)\in B_{<q}(Z)$. As $w-v\in P$, the full quotient-kernel
condition makes it a final boundary. This proves permanence and the
claimed grade of the same classical image.
\end{proof}

\begin{corollary}[Permanence from finite lifts with intermediate-space control]
\label{sem:ProofTools:coverage-finite-lift-permanence}
\level{coarse}
\uses{sem:ProofTools:common-labels,sem:Literature:lambda-inversion-adams-comparison,sem:ProofTools:finite-common-lift,sem:ProofTools:finite-strip,sem:Literature:finite-bockstein,sem:Literature:infinite-finite-window,sem:Literature:lambda-inversion-detection,sem:Literature:complete-boundary-lift,sem:Literature:complete-permanent-lifts}
\lean{KIP126.ProofTools.Internal.SyntheticCoverage.finite_permanence_positive_h,KIP126.ProofTools.Internal.SyntheticCoverage.finite_permanence_positive_g,KIP126.ProofTools.Internal.SyntheticCoverage.finite_permanence_positive_f}
\lean{KIP126.ProofTools.Internal.SyntheticCoverage.selective_permanence_positive_h,KIP126.ProofTools.Internal.SyntheticCoverage.selective_permanence_positive_g,KIP126.ProofTools.Internal.SyntheticCoverage.selective_permanence_positive_f}

\textit{Vanishing intermediate third pages.}
Use the compatible triangle with $\alpha+\beta+\gamma=1$.
Let $q\ge2$, $s\ge\gamma$, $n=t-s$ and $p=s+\alpha+q-1$.
Let $x\in Z_\infty^{s,t}(X)$ strictly detect a classical class,
with $F_2x=0$, and suppose
\[
 U\in Z_{<q}^{s-\gamma,t-\gamma+1}(Z),\qquad H_2U=x.
\]
Take $Z_\infty^{p,p+n}(Y)\subseteq P\subseteq E_2^{p,p+n}(Y)$
and $v\in P$, with $G_2v\in Z_{<q}(Z)$ and
$d_q[U]=[G_2v]$. Assume
\[
 E_3^{s+j,t+j}(X)=0\quad(1\le j\le q-2),\qquad
 F_2(E_2^{s+q-1,t+q-1}(X))\cap P\subseteq B_\infty^{p,p+n}(Y).
\]
Then $v$ is permanent and, for every $\xi$ strictly detected by $x$,
\[
 f_*\xi\in F^p\pi_nY,\qquad
 \operatorname{gr}^{p}(f_*\xi)=[v]_\infty.
\]
The vanishing hypothesis concerns the entire $E_3$ group
in every displayed degree.
For $q=2$ the intermediate conditions are empty.

\textit{Restricted intermediate spaces.}
Keep the triangle, $q,s,t,n,p,x,U,v,P$, first-image equation,
main differential and terminal intersection condition of
the vanishing-intermediate-pages form above.
For each $1\le j\le q-2$, put $L=q-j$ and choose
\[
 Z_\infty^{s+j+\alpha,t+j+\alpha}(Y)
 \subseteq P_j\subseteq E_2^{s+j+\alpha,t+j+\alpha}(Y).
\]
Replace intermediate $E_3$ vanishing by the whole-space condition
\[
 Z_{<L+1}^{s+j,t+j}(X)\cap F_2^{-1}(P_j)
 \subseteq B_\infty^{s+j,t+j}(X)
 +H_2\bigl(Z_{<L+1}^{s+j-\gamma,t+j-\gamma+1}(Z)\bigr).
\]
Then the same permanence and image conclusion holds.
In fact, for permanence itself $x$ need only be permanent: the image
conclusion holds for every strictly detected classical source whenever
one is specified, without assuming that such a source exists.
For $q=3$ the only intermediate cycle spaces are $Z_{<3}=\ker d_2$.
\end{corollary}
\begin{proof}
\textit{Vanishing intermediate third pages.}
Correct the finite lift of $U$ so that its boundary is the specified
$G_2v$. The common-lift theorem then constructs
$c_q\in\pi_{n,t}Q_qA$ with first label $x$ and
$\varphi_qc_q=\iota_{q-1,q}v$.
Fix a full lift $a$ of the prescribed classical source.
The difference $\epsilon_qa-c_q$ has zero first restriction.
The finite-window description and the vanishing of every intermediate
$E_3$ group place this difference in the last coefficient layer:
$\epsilon_qa-c_q=\iota_{q-1,q}e$.
Thus $\varphi a=\lambda^{q-1}b$ for a full class $b$.
Its first label $w$ is permanent, realizes the same $f_*\xi$, and
satisfies $w-v-F_2e\in B_{<q+1}(Y)$.
Since $w,v$ and true boundaries belong to $P$, the terminal
intersection condition makes $F_2e$ a final boundary.
Consequently $[v]_\infty=[w]_\infty$, as required.

\textit{Restricted intermediate spaces.}
Start with the same $c_q$ as in the vanishing-intermediate-pages proof and a full lift $a$
of $x$. At stage $j$, write
$\epsilon_qa-c_q=\iota_{j,q}e_j$, where $e_j\in Q_LA$.
Its first label $z_j$ lies in $Z_{<L+1}(X)$.
The full image is divisible by $\lambda^j$, and comparison in
$Q_{j+1}B$ shows that $F_2z_j$ differs from a permanent label by
a member of $B_{<j+2}(Y)$. Hence $F_2z_j\in P_j$.

Decompose $z_j=b_j+H_2h_j$ by the stated inclusion.
A full lambda-power-torsion lift of the true boundary $b_j$ changes
$a$ by a class divisible by $\lambda^j$; this preserves its first
label and its specified classical realization. A finite lift of $h_j$
changes $c_q$ by a connecting image, preserving both its first label
and its entire $\varphi_q$ image. The new residual has zero first
label and enters stage $j+1$.
After finitely many steps only the last coefficient layer remains.
The same terminal intersection argument as in
the vanishing-intermediate-pages form above proves the result.
For permanence alone begin with any full lift of the permanent $x$;
for the image assertion begin with a lift of the prescribed source.
\end{proof}


\begin{corollary}[Permanence with a full torsion correction]
\label{sem:ProofTools:coverage-torsion-affine-permanent}
\level{coarse}
Use the compatible triangle, with $\alpha+\beta+\gamma=1$.
Let $q\ge2$, $n=t-s$, $p=s+\alpha+q-1$ and choose
$x\in Z_\infty^{s,t}(X)$ and
$v\in P$, where
$Z_\infty^{p,p+n}(Y)\subseteq P\subseteq E_2^{p,p+n}(Y)$.
Suppose there are actual classes
\[
 c\in\pi_{n,t}Q_qA,\qquad T\in\pi_{n,t}B,\qquad
 \lambda^NT=0\quad\text{for some }N\ge1,
\]
such that
\[
 \epsilon c=x',\qquad x'-x\in B_\infty^{s,t}(X),\qquad
 \varphi_qc=\epsilon_qT+\iota_{q-1,q}v.
\]
For every $1\le j\le q-2$, put $L=q-j$ and take
$P_j\subseteq E_2^{s+j+\alpha,t+j+\alpha}(Y)$ containing all
permanent labels in that degree. Require
\[
 Z_{<L+1}^{s+j,t+j}(X)\cap F_2^{-1}(P_j)
 \subseteq B_\infty^{s+j,t+j}(X)
 +H_2\bigl(Z_{<L+1}^{s+j-\gamma,t+j-\gamma+1}(Z)\bigr),
\]
and the terminal condition
\[
 F_2(E_2^{s+q-1,t+q-1}(X))\cap P
 \subseteq B_\infty^{p,p+n}(Y).
\]
Then $v$ is permanent. For every $\xi\in\pi_nX$ strictly detected
by $x$, its same image satisfies
\[
 f_*\xi\in F^p\pi_nY,\qquad
 \operatorname{gr}^{p}(f_*\xi)=[v]_\infty.
\]
Neither the existence of a strictly detected source nor nonvanishing
of the target infinite-page class is required for permanence.
\uses{sem:ProofTools:common-labels,sem:Literature:lambda-inversion-adams-comparison}
\lean{KIP126.ProofTools.Internal.SyntheticCoverage.torsion_affine_permanence_positive_h,KIP126.ProofTools.Internal.SyntheticCoverage.torsion_affine_permanence_positive_g,KIP126.ProofTools.Internal.SyntheticCoverage.torsion_affine_permanence_positive_f}
\end{corollary}
\begin{proof}
\uses{sem:Literature:complete-boundary-lift,sem:Literature:complete-permanent-lifts,sem:ProofTools:finite-strip,sem:Literature:lambda-inversion-detection}
Choose a full lift $a$ of $x$. A full torsion lift of $x'-x$ changes
its first label to $x'$ without changing its classical realization.
Now $\epsilon_qa-c=\iota_{1,q}e_1$.
If the residual is in layer $j$, the equation for $\varphi_qc$
implies that the full class $\varphi a-T$ is divisible by
$\lambda^j$. Its first-label comparison shows
$F_2\epsilon e_j\in P_j$, since $P_j$ contains permanent labels
and all true boundaries. The middle inclusion decomposes
$\epsilon e_j$ into a final boundary and a finite connecting image.
Correct $a$ by the corresponding full torsion lift and $c$ by the
finite connecting image. Both the classical source and the equation
$\varphi_qc=\epsilon_qT+\iota_{q-1,q}v$ are retained.

At the last layer,
$\varphi a-T=\lambda^{q-1}b$ and
$w-v-F_2e\in B_{<q+1}(Y)$, where $w=\epsilon_1b$ is permanent.
The terminal intersection makes $w-v$ a final boundary.
All full torsion corrections, including $T$, vanish after lambda
inversion, so $b$ realizes the same $f_*\xi$ when the initial lift
was chosen for a prescribed $\xi$. An arbitrary full lift gives permanence.
\end{proof}

\begin{corollary}[A full two-layer relation controls a three-layer error]
\label{sem:ProofTools:coverage-recursive-three-layer-permanent}
\level{coarse}
Use the compatible triangle with $\alpha+\beta+\gamma=1$, and put
$n=t-s$, $p=s+\alpha+2$.
Let $x\in Z_\infty^{s,t}(X)$ and
$v\in P\subseteq E_2^{p,p+n}(Y)$, with $P$ containing every
permanent label in that degree. Suppose a single finite class satisfies
\[
 c_3\in\pi_{n,t}Q_3A,\qquad \epsilon c_3=x,\qquad
 \varphi_3c_3=\iota_{2,3}v.
\]
Write
\[
\begin{gathered}
 X_1=E_2^{s+1,t+1}(X),\quad X_2=E_2^{s+2,t+2}(X),\quad
 Z_1=E_2^{s+1-\gamma,t+2-\gamma}(Z),\\
 Y_1=E_2^{s+1+\alpha,t+1+\alpha}(Y),\quad
 Y_2=E_2^{p,p+n}(Y).
\end{gathered}
\]
For these named graded pieces, notation such as $Z_{<3}(X_1)$ or
$B_\infty(Y_2)$ means the cycle or boundary subspace of the original
spectral sequence in the displayed degree.
Choose $P_1\subseteq Y_1$ containing all permanent labels and put
$K=Z_{<3}(X_1)\cap F_2^{-1}(P_1)$.
Let $K_0\subseteq\ker(d_2|_{X_1})\cap\ker(F_2|_{X_1})$ satisfy
\[
 K\subseteq B_\infty(X_1)+H_2(Z_{<3}(Z_1))+K_0.
\]
For the entire linear relation
\[
 \mathcal R=\{(U,k,z)\in Z_1\oplus K_0\oplus Y_2:
                  H_2U=k,\ d_2U=G_2z\},
\]
assume that $\mathcal R\to K_0$ is surjective and
\[
 (\operatorname{pr}_z\mathcal R+F_2(X_2))\cap P
 \subseteq B_\infty(Y_2).
\]
Then $v$ is permanent and, for every $\xi$ strictly detected by $x$,
\[
 f_*\xi\in F^p\pi_nY,\qquad
 \operatorname{gr}^{p}(f_*\xi)=[v]_\infty.
\]
The permanence assertion allows a zero infinite-page target class;
the image formula applies to every strictly detected source.
\uses{sem:ProofTools:common-labels,sem:Literature:lambda-inversion-adams-comparison}
\lean{KIP126.ProofTools.Internal.SyntheticCoverage.recursive_positive_h_permanence,KIP126.ProofTools.Internal.SyntheticCoverage.recursive_positive_g_permanence,KIP126.ProofTools.Internal.SyntheticCoverage.recursive_positive_f_permanence}
\end{corollary}
\begin{proof}
\uses{sem:ProofTools:finite-common-lift,sem:Literature:complete-boundary-lift,sem:Literature:complete-permanent-lifts,sem:Literature:lambda-inversion-detection}
Choose a full lift $a$ of $x$ and write
$\epsilon_3a-c_3=\iota_{1,3}e_2$. The actual first residual belongs
to $K$. Decompose it by the stated inclusion and remove its final
boundary and finite connecting-image parts by compatible corrections
to $a$ and $c_3$. The corrected full lift has the same classical
realization, while $c_3$ retains its first label and entire image.
The residual now has first label $k\in K_0$.

Surjectivity supplies one triple $(U,k,z)\in\mathcal R$ for this
actual $k$. The two-layer common-lift construction gives
$c_k\in Q_2A$ with first label $k$ and
$\varphi_2c_k=\iota_{1,2}z$.
Thus $e_2-c_k=\iota_{1,2}u$ for $u\in X_2$, and comparison
inside the same $Q_3$ gives
$\varphi_3\epsilon_3a=\iota_{2,3}(v+z+F_2u)$.
Write $\varphi a=\lambda^2b$ and let $w=\epsilon_1b$.
Then $w$ is permanent and
$w-v-z-F_2u\in B_{<4}(Y_2)$.
Since $w,v$ and true boundaries lie in $P$, the entire error
$z+F_2u$ belongs to the displayed terminal intersection and is a
final boundary. Hence $[w]_\infty=[v]_\infty$.
The realization comparison identifies this with the grade of the
same prescribed $f_*\xi$; an arbitrary initial full lift suffices
for permanence. Surjectivity chooses this compatible triple for the actual residual.
\end{proof}

\begin{corollary}[Comparison of leading terms in one finite filtration]
\label{sem:ProofTools:coverage-finite-general}
\level{coarse}
Let $A_N,B_N$ be filtered abelian groups with decreasing filtrations,
and let $\phi_N:A_N\to B_N$ be a homomorphism. Fix source and target
filtrations $s,p$ and suppose
\[
 \phi_N(F^{s+1}A_N)\subseteq F^{p+1}B_N.
\]
For any specified $a_N,c_N\in A_N$ such that
$a_N-c_N\in F^{s+1}A_N$ and $\phi_N(c_N)\in F^pB_N$, one has
\[
 \phi_N(a_N)\in F^pB_N,\qquad
 \phi_N(a_N)-\phi_N(c_N)\in F^{p+1}B_N.
\]
Their target grades therefore agree, and one belongs to $F^{p+1}B_N$
if and only if the other does. In particular, this preserves a specified
nonzero leading term. If $a_N$ is the projection of a full source and
the projection square commutes, this identifies the leading term of its finite image with
that of $\phi_N(c_N)$.
The displayed control quantifies over the entire higher-source subgroup,
including all finite directions.
\uses{sem:SpectralSequence:filtration,sem:ProofTools:filtration-representatives}
\lean{KIP126.ProofTools.Internal.SyntheticCoverage.finite_filtered_image_comparison}
\end{corollary}
\begin{proof}
Apply the full subgroup inclusion to $a_N-c_N$ and use additivity:
$\phi_N(a_N)-\phi_N(c_N)\in F^{p+1}B_N$.
As $F^{p+1}B_N\subseteq F^pB_N$, the first membership follows,
and reduction modulo $F^{p+1}B_N$ gives equality of the two grades.
This proves the comparison by filtered-group operations. When applied to a full source, the same projected image is
identified by the stated naturality square. Applications to permanence use the finite-to-full detection comparison
with its stated visible-window hypotheses.
\end{proof}


\subsection{Extension equations with adjustable higher source}

The preceding criteria remove higher-layer errors while preserving the
classical source. We now allow an error component represented by a
higher-filtration permanent class. Removing that component changes the
source from $\xi$ to $\xi-\eta$, with $\eta\in F^{s+1}\pi_*X$,
so its leading grade remains unchanged. The conclusion is a classical
extension equation for this corrected representative.

The argument allows all three placements of the positive normalization
exponent and zero source or target grades. The target grade concerns
$\xi-\eta$. The final proposition gives a stronger residual criterion
that preserves the original $\xi$.

\begin{proposition}[Extension equations after a higher-filtration source correction]
\label{sem:ProofTools:coverage-ess-source-adjustment}
\level{coarse}
\lean{KIP126.ProofTools.Internal.SyntheticCoverage.source_adjusted_affine_permanence_positive_h,KIP126.ProofTools.Internal.SyntheticCoverage.source_adjusted_affine_permanence_positive_g,KIP126.ProofTools.Internal.SyntheticCoverage.source_adjusted_affine_permanence_positive_f}
Use the normalized Adams triangle, strong convergence, and the full
and finite-quotient comparisons of
Corollary~\ref{sem:ProofTools:coverage-torsion-affine-permanent}:
\[
 A=\nu X\xrightarrow{\varphi}
 B=\Sigma^{0,-\alpha}\nu Y\longrightarrow
 C=\Sigma^{0,-\alpha-\beta}\nu Z\longrightarrow\Sigma A,
 \qquad \alpha,\beta,\gamma\geq0,\quad \alpha+\beta+\gamma=1.
\]
Write $F_2,G_2,H_2$ for its normalized first-label maps, and
$Q_mD=D/\lambda^m$. The notation $Z_{<r}$ means the common labels
reaching page $r$, namely those on which all earlier outgoing Adams
differentials vanish; $B_\infty$ consists of true Adams
boundaries. Put
\[
 q\geq2,\qquad n=t-s,\qquad p=s+\alpha+q-1.
\]
Let $x\in Z_\infty^{s,t}(X)$, and choose
$v\in P\subseteq E_2^{p,p+n}(Y)$ with
$Z_\infty^{p,p+n}(Y)\subseteq P$. Suppose a single finite construction gives
\[
 c\in\pi_{n,t}Q_qA,\qquad T\in\pi_{n,t}B,\qquad \lambda^NT=0
 \quad\text{for some }N\geq1,
\]
such that
\[
 \epsilon c=x',\qquad x'-x\in B_\infty^{s,t}(X),\qquad
 \varphi_qc=\epsilon_qT+\iota_{q-1,q}v.
\]
For $1\leq j\leq q-1$, let
$C_j\subseteq Z_\infty^{s+j,t+j}(X)$ be subspaces of true permanent labels.
For $1\leq j\leq q-2$, put $L=q-j$, choose
\[
 Z_\infty^{s+j+\alpha,t+j+\alpha}(Y)
 \subseteq P_j\subseteq E_2^{s+j+\alpha,t+j+\alpha}(Y),
\]
and require the whole-subspace inclusion
\[
 Z_{<L+1}^{s+j,t+j}(X)\cap F_2^{-1}(P_j)
 \subseteq B_\infty^{s+j,t+j}(X)
 +H_2\!\left(Z_{<L+1}^{s+j-\gamma,t+j-\gamma+1}(Z)\right)+C_j.
\]
At the terminal layer require
\[
 F_2\!\left(E_2^{s+q-1,t+q-1}(X)\right)\cap P
 \subseteq B_\infty^{p,p+n}(Y)+F_2(C_{q-1}).
\]
Then $v\in Z_\infty(Y)$. For every
$\xi\in F^s\pi_nX$ with $\operatorname{gr}^s\xi=[x]_\infty$, there exists
$\eta\in F^{s+1}\pi_nX$ such that
\[
 f_*(\xi-\eta)\in F^p\pi_nY,\qquad
 \operatorname{gr}^p f_*(\xi-\eta)=[v]_\infty.
\]
Consequently the classical $f$-extension spectral sequence satisfies
\[
 d^f_{p-s}[x]_\infty=[v]_\infty.
\]
This equation includes zero source or target classes and may be inessential.
If $[x]_\infty\neq0$, the corrected source is still strictly detected by
the same label $x$. The target-grade formula applies to the corrected source $\xi-\eta$. Taking every $C_j=0$ recovers the
source-preserving conclusion of the preceding affine comparison.
\uses{sem:ProofTools:coverage-torsion-affine-permanent,sem:ProofTools:common-labels,sem:ProofTools:extension-representatives}
\end{proposition}
\begin{proof}
\uses{sem:Literature:complete-boundary-lift,sem:Literature:complete-permanent-lifts,sem:Literature:finite-bockstein,sem:Literature:lambda-inversion-detection,sem:ProofTools:extension-representatives}
\uses{sem:AdamsSynthetic:completed-adams-scope}
Fix such a $\xi$ and choose its compatible complete lift $a$ with first
label $x$. This also applies when $[x]_\infty=0$: separatedness gives every
nonzero $\xi$ a finite exact filtration, where the prescribed strict lift
is available. Multiplication by a lambda power lowers its weight to the
required one without changing $\xi$, and a complete torsion correction
matches $x$. For $\xi=0$ use the torsion lift of its boundary label.
Add a complete lambda-power-torsion lift of $x'-x$; this changes
neither its classical realization nor its source grade and makes
$\epsilon a=x'$. Thus $a_q-c=\iota_{1,q}e_1$.

Inductively retain a current complete lift $a$ and a current finite element
$c$ satisfying
\[
 a_q-c=\iota_{j,q}e_j,\qquad e_j\in\pi_{n,t+j}Q_LA,
 \qquad \varphi_qc=\epsilon_qT+\iota_{q-1,q}v.
\]
The complete element $\varphi a-T$ is zero in $Q_jB$, hence equals
$\lambda^jb_j$. Comparing in $Q_{j+1}B$ gives
\[
 \epsilon b_j-F_2(\epsilon e_j)\in B_{<j+2}(Y).
\]
The finite comparison puts $z_j=\epsilon e_j$ in $Z_{<L+1}(X)$.
The first label of the complete element $b_j$ and every true boundary
belong to $P_j$, so $F_2z_j\in P_j$.
Use the middle inclusion to write
\[
 z_j=d_j+H_2h_j+c_j,\qquad d_j\in B_\infty(X),\quad c_j\in C_j.
\]
Choose a complete lambda-power-torsion lift $D_j$ of $d_j$ and a complete
permanent lift $W_j$ of $c_j$. Lift $h_j$ in $Q_LC$ and let $k_j\in Q_LA$
be its connecting image, with the indicated shifts. Then
$\epsilon k_j=H_2h_j$ and $\varphi_Lk_j=0$.
Make the simultaneous replacements
\[
 a\longleftarrow a-\lambda^j(D_j+W_j),\qquad
 c\longleftarrow c+\iota_{j,q}k_j.
\]
The first label of the new difference is zero, so it has the form
$\iota_{j+1,q}e_{j+1}$. The first label and the entire image of $c$ stay
fixed. The torsion correction changes no classical class. The correction
from $W_j$ changes the classical source by an element of $F^{s+j}\pi_nX$,
which lies in $F^{s+1}\pi_nX$ since $j\geq1$. Thus the original source
grade remains $[x]_\infty$ throughout the induction.

After finitely many steps,
\[
 a_q-c=\iota_{q-1,q}e,\qquad
 \varphi a-T=\lambda^{q-1}b,
 \qquad w=\epsilon b\in Z_\infty(Y).
\]
The terminal comparison yields
$w-v-F_2e\in B_{<q+1}(Y)$. Since $w,v$ and the boundary term lie in $P$,
the terminal condition writes $F_2e=d+F_2c_{q-1}$ with
$d\in B_\infty(Y)$ and $c_{q-1}\in C_{q-1}$.
Choose a complete lift $W_{q-1}$ with that first label and set
$a^*=a-\lambda^{q-1}W_{q-1}$. Then
\[
 \varphi a^*-T=\lambda^{q-1}(b-\varphi W_{q-1}),\qquad
 \epsilon(b-\varphi W_{q-1})-v\in B_\infty(Y).
\]
The complete first label on the left is permanent, so $v$ is permanent.
All classical source changes sum to a single
$\eta\in F^{s+1}\pi_nX$. Inverting lambda kills $T$; the same compatible
comparison gives the asserted filtration and grade of $f_*(\xi-\eta)$.
Strong convergence supplies representatives of $[x]_\infty$, including
when this class is zero. The filtered two-term complex
$\pi_nX\xrightarrow{f_*}\pi_nY$ now gives the stated extension equation
from the representative $\xi-\eta$. The same construction applies to
zero source and target classes.
\end{proof}


\paragraph{Obtaining the finite input from a differential.}
For the initial finite construction in
Proposition~\ref{sem:ProofTools:coverage-ess-source-adjustment}, it is
enough to supply a common label
\[
 U\in Z_{<q}^{s-\gamma,t-\gamma+1}(Z),\qquad
 H_2U-x\in B_\infty^{s,t}(X),\qquad
 G_2v\in Z_{<q}(Z),\qquad d_q[U]=[G_2v].
\]
The last equality is in the actual page $q$, with its true earlier-boundary
denominator. Keep all its middle and terminal subspace conditions.
Then the same permanence and classical extension conclusions hold.
The case $C_j=0$ also gives the source-preserving affine conclusion.

To verify this sufficient input,
apply the target-preserving common-lift construction to this one $U$ and
the specified $v$. It gives
$c\in\pi_{n,t}Q_qA$ with first label $x'=H_2U$ and
$\varphi_qc=\iota_{q-1,q}v$. Thus the affine hypothesis holds with $T=0$;
the difference $x'-x$ is the supplied true boundary. The composite
$F_2H_2=0$ holds on this same triangle, so no separate literal equality
$F_2x=0$ is required. Apply the proposition.

\paragraph{Normalization of the target by true boundaries.}
This differential input may instead use a label $v'$ with
$v'-v\in B_\infty(Y)$ and $d_q[U]=[G_2v']$, where
$G_2v'\in Z_{<q}(Z)$. Since $P$ contains all true boundaries, $v'\in P$.
Apply the same comparison with the specified finite target $v'$.
It proves that $v'$ is permanent and yields target grade $[v']_\infty$.
The original label $v$ is then permanent with the same grade.
Thus changing the target by a strictly earlier true boundary
preserves the extension equation. A complete torsion
lift $D$ of $v'-v$ supplies the affine term
$T=\lambda^{q-1}D$, using the same finite quotient and layer inclusion.

\paragraph{Composition with an ordinary filtered map.}
The source-preserving comparison has the following non-strict formulation:
for every $\zeta\in F^u\pi_*Y$ with a prescribed grade $[w]_\infty$,
the image $h_*\zeta$ has the asserted filtration and grade. Its proof uses
only a complete lift with first label $w$. If an initial complete lift has
another first label representing $[w]_\infty$, add a complete torsion lift
of their boundary difference. The same proof therefore works when
$[w]_\infty=0$; strict detection is needed only to assert nonzero grades.

In particular, let $g:X\to Y$ have filtration at least $a$, with the
specified first-label map $g_2$, and suppose $g_2x=w$. If a previously
proved source-preserving comparison for $h:Y\to Z$ applies to every
$\zeta\in F^{s+a}\pi_*Y$ of grade $[w]_\infty$ and gives target grade
$[v]_\infty$ in filtration $p$, then it applies to $\zeta=g_*\xi$ for
every $\xi\in F^s\pi_*X$ of grade $[x]_\infty$. Thus the same target
grade is obtained for $(h\circ g)_*\xi$. This argument uses the actual
filtered composite and its fixed labels. The source-preserving comparison for $h$ applies to this specified
representative $[w]_\infty$ in the image of $g_*$.

\paragraph{A higher early boundary as the initial finite input.}
Let $a\geq2$, put $L=q+a-1$, and suppose that a single label
\[
 c_0\in Z_{<a}^{s+\alpha-a,t+\alpha-a+1}(Y),\qquad
 d_a[c_0]=[F_2x]
\]
gives $z=G_2c_0\in Z_{<L}(Z)$ with
$G_2v\in Z_{<L}(Z)$ and $d_L[z]=[G_2v]$.
All equations use the actual page and its full earlier-boundary denominator.
For $1\leq j\leq a-2$, in degree $\deg z+(j,j)$ require
\[
 Z_{<a-j+1}(Z)\subseteq B_{<j+2}(Z)
 +G_2\bigl(Z_{<a-j+1}(Y)\bigr)+Z_{<L-j+1}(Z),
\]
where the source degree of $G_2$ is obtained using its actual degree shift.
Also require
\[
 H_2\!\left(E_2^{s-\gamma,t-\gamma+1}(Z)\right)
 \subseteq B_{<a+1}^{s,t}(X).
\]
The higher-early comparison then gives one
$c\in\pi_{n,t}Q_qA$ and one complete $T\in\pi_{n,t}B$ with
\[
 \epsilon c-x\in B_{<a+1}(X),\qquad
 \varphi_qc=\epsilon_qT+\iota_{q-1,q}v,\qquad
 \lambda^{a-1}T=0.
\]
These satisfy the initial finite image equation of the proposition. Together with the $q$-layer middle and terminal conditions,
this gives the input for the proposition.

For completeness, the early differential gives a rotated element
$R\in Q_aC$ with first label $z$ and
$\chi_aR=\iota_{a-1,a}x$. Choose a representative of the later differential
$\widehat z\in Q_{L-1}C$ with the same first label and the specified
adjacent boundary $G_2v$. Compare $R$ and $\widehat z_a$ layer by layer.
A $B_{<j+2}$ term lifts to a complete element killed by $\lambda^j$,
so its insertion has no effect. A $G_2$ term lifts to a genuine
$\psi$-image, whose connecting image is zero. For the last summand choose
a lift in $Q_{L-j}C$ and restrict it to $Q_{L-1-j}C$ before inserting it
into the representative of the later differential. The one-layer-longer lift makes its actual
adjacent boundary zero, preserving both $z$ and the prescribed $G_2v$.
After these corrections, the final $H_2$ condition removes the remaining
top discrepancy modulo $\ker\iota_{a-1,a}=B_{<a+1}$.
The common-lift construction applied to this representative supplies $a_L\in Q_LA$ with zero
$Q_{a-1}A$ projection and image $\iota_{L-1,L}v$.
Divide $a_L$ through $\iota_{a-1,L}$ to obtain $c\in Q_qA$.
The kernel of the same coefficient inclusion is the projection of an
absolute boundary $T$ killed by $\lambda^{a-1}$, giving the stated image
formula. Every comparison uses this same corrected representative;
the original target label is not replaced.

\begin{proposition}[A finite residual whose torsion vanishes after insertion]
\label{sem:ProofTools:coverage-torsion-residual}
\level{coarse}
\lean{KIP126.ProofTools.Internal.SyntheticCoverage.torsion_residual_permanence_positive_h,KIP126.ProofTools.Internal.SyntheticCoverage.torsion_residual_permanence_positive_g,KIP126.ProofTools.Internal.SyntheticCoverage.torsion_residual_permanence_positive_f}
Use the same normalized triangle, strong convergence, and full
and finite-quotient comparisons as in
Corollary~\ref{sem:ProofTools:coverage-torsion-affine-permanent}.
Put $n=t-s$ and $p=s+\alpha+3$.
Let $x\in Z_\infty^{s,t}(X)$ and choose subspaces
\[
 Z_\infty^{s+\alpha+1,t+\alpha+1}(Y)\subseteq P_1
 \subseteq E_2^{s+\alpha+1,t+\alpha+1}(Y),\qquad
 Z_\infty^{p,p+n}(Y)\subseteq P\subseteq E_2^{p,p+n}(Y).
\]
Let $v\in P$ and suppose one finite construction provides
\[
 c\in\pi_{n,t}Q_4A,\qquad T_0\in\pi_{n,t}B,\qquad
 \lambda^NT_0=0\quad\text{for some }N\geq1,
\]
with
\[
 \epsilon c=x',\qquad x'-x\in B_\infty^{s,t}(X),\qquad
 \varphi_4c=\epsilon_4T_0+\iota_{3,4}v.
\]
Suppose
\[
 K:=Z_{<4}^{s+1,t+1}(X)\cap F_2^{-1}(P_1)
 \subseteq B_\infty^{s+1,t+1}(X)
 +H_2\!\left(Z_{<4}^{s+1-\gamma,t+2-\gamma}(Z)\right)+K_0,
\]
where $K_0\subseteq Z_{<4}^{s+1,t+1}(X)$ is a subspace, and assume
\[
 E_2^{s+2,t+2}(X)=0.
\]
Choose a subspace $R\subseteq E_2^{p,p+n}(Y)$. For every $k\in K_0$,
require the existence of
\[
 d_k\in\pi_{n,t+1}Q_3A,\qquad
 T_k\in\pi_{n,t+1}B,\qquad z_k\in R
\]
satisfying
\[
 \epsilon d_k-k\in B_\infty^{s+1,t+1}(X),\qquad\lambda T_k=0,\qquad
 \varphi_3d_k=\epsilon_3T_k+\iota_{2,3}z_k.
\]
Choose the representatives separately for each $k$. When $K_0$ is finite
dimensional, basis witnesses supply witnesses for every $k$: choose
coefficients in $\{0,1\}$ and sum the finite representatives, torsion
elements, and terminal labels. Finally suppose
\[
 \left(R+F_2\!\left(E_2^{s+3,t+3}(X)\right)\right)\cap P
 \subseteq B_\infty^{p,p+n}(Y).
\]
Then $v\in Z_\infty(Y)$, and for every
$\xi\in F^s\pi_nX$ with $\operatorname{gr}^s\xi=[x]_\infty$ one has
\[
 f_*\xi\in F^p\pi_nY,\qquad
 \operatorname{gr}^p f_*\xi=[v]_\infty.
\]
In particular the classical extension equation
$d^f_{p-s}[x]_\infty=[v]_\infty$ holds. The same prescribed classical
source $\xi$ is retained. The elements of $K_0$ need only admit the stated
finite representatives; they are not required to be permanent.
\uses{sem:ProofTools:coverage-torsion-affine-permanent,sem:ProofTools:common-labels,sem:ProofTools:extension-representatives}
\end{proposition}
\begin{proof}
\uses{sem:Literature:complete-boundary-lift,sem:Literature:complete-permanent-lifts,sem:Literature:finite-bockstein,sem:Literature:lambda-inversion-detection,sem:ProofTools:extension-representatives}
\uses{sem:AdamsSynthetic:completed-adams-scope}
Fix $\xi$ and choose a compatible complete lift $a$, including a zero
infinity class by the separated-filtration argument above. Correct its first
label by complete lambda-power torsion so that $\epsilon a=x'$ without
changing $\xi$. Then $a_4-c=\iota_{1,4}e$ for some
$e\in\pi_{n,t+1}Q_3A$. The finite comparison places its first label in
$Z_{<4}(X)$.
Since $\varphi a-T_0$ is zero in $Q_1B$, write it as $\lambda b_1$.
Comparison in $Q_2B$ gives
\[
 \epsilon b_1-F_2\epsilon e\in B_{<3}(Y).
\]
Both the complete first label and the true boundary lie in $P_1$;
therefore $\epsilon e\in K$. Write $\epsilon e=b+H_2h+k$ with
$b\in B_\infty(X)$, $h\in Z_{<4}(Z)$ in the indicated degree, and
$k\in K_0$. Choose a complete torsion lift $D$ of $b$, lift $h$ in
$Q_3C$, and denote its connecting image by $\chi_3\widehat h\in Q_3A$.
Make the replacements
\[
 a\longleftarrow a-\lambda D,\qquad
 c\longleftarrow c+\iota_{1,4}\chi_3\widehat h.
\]
The classical source stays $\xi$. The first label and entire image of
$c$ are unchanged because $\varphi_3\chi_3=0$.
The new difference has a representative $e$ with first label exactly $k$.

Choose the supplied $d_k$ and a complete torsion lift $D_k$ of
$\epsilon d_k-k$. Replace $a$ by $a+\lambda D_k$ and $e$ by
$e+\epsilon_3D_k$. The prescribed classical source is unchanged, and
now $e$ has exactly the same first label as $d_k$. Hence $e-d_k$ belongs
to the image of $\iota_{1,3}$ from $Q_2A$.
The first-label group of that remaining $Q_2A$ is
$E_2^{s+2,t+2}(X)=0$. A second application of the same coefficient
exact sequence therefore gives
\[
 e-d_k=\iota_{2,3}u,\qquad u\in E_2^{s+3,t+3}(X).
\]
After this boundary adjustment of the complete lift of the same $\xi$,
the argument applies to every finite representative $e$ of the matched
first label, not only the chosen $d_k$. It controls the entire affine
family of possible finite images.
Applying $\varphi$ and inserting in $Q_4B$ now gives
\[
 \begin{split}
 \varphi_4(a_4-c)
 &=\iota_{1,4}\epsilon_3T_k
   +\iota_{3,4}(z_k+F_2u)\\
 &=\iota_{3,4}(z_k+F_2u),
 \end{split}
\]
because $\iota_{1,4}\epsilon_3T_k=\epsilon_4(\lambda T_k)=0$.
Thus $\varphi a-T_0$ is zero in $Q_3B$ and has the form
$\lambda^3b_3$ for a complete element $b_3$.
Its first label $w=\epsilon b_3$ is permanent, and the top-layer
comparison yields
\[
 w-v-z_k-F_2u\in\ker\iota_{3,4}=B_{<5}(Y).
\]
As $w,v$ and the boundary term belong to $P$, the element
$z_k+F_2u$ lies in $\left(R+\operatorname{im}F_2\right)\cap P$.
The terminal condition makes it a true boundary, hence $w-v\in B_\infty$.
This proves that $v$ is permanent. Inverting lambda kills $T_0$ and the
source torsion correction; the complete comparison gives precisely the
stated filtration and grade of the original $f_*\xi$. The representative
description of the filtered two-term complex gives the extension equation,
including zero source or target grades.
\end{proof}


\paragraph{Producing a three-layer residual from an early boundary.}
For a single residual label $k$, suppose
$d_2c_0=F_2k$, put $z=G_2c_0$, and suppose
$z,G_2r\in Z_{<4}(Z)$ with $d_4[z]=[G_2r]$.
The entire ambiguity in the first source label must satisfy the strong early condition
$H_2(E_2(Z))\subseteq B_{<3}(X)$ in the corresponding degree.
The early-boundary construction produces $d\in Q_3A$ with
first label $k'$ satisfying $k'-k\in B_{<3}(X)$, and
$\varphi_3d=\epsilon_3T+\iota_{2,3}r$ with $\lambda T=0$.
These already satisfy the boundary-congruent finite input in the
proposition, with $\lambda T=0$. The proof adjusts the complete source lift to match $k'$, retaining $T$
and the prescribed terminal label $r$. The finite input allows
nonpermanent $k$.
